\section{Supplement on \'etale morphisms. Henselian local rings and strictly local rings}
\label{section:IV.18}
\egalabel{IV.18}{18}

In the present section we study various properties peculiar to \'etale morphisms.  Moreover, the notion of an \'etale morphism makes it possible to develop, in a very natural way, Nagata's theory of Henselian rings, as well as that of strictly local rings.  These rings play an important role in many recent developments, appearing whenever one needs a process of ``localisation'' finer than that furnished by the Zariski topology (cf.\ for example~[43]), pending the appearance of the chapter of our Treatise devoted to the study of the ``\'etale topology.''

\subsection{A remarkable equivalence of categories.}
\label{subsection:IV.18.1}

\begin{proposition}[18.1.1]
\label{IV.18.1.1}
Let $S$ be a prescheme, $S_0$ a closed subprescheme of $S$, $X_0$ an $S_0$-prescheme smooth \textup{(resp.}\ \'etale\textup{)} over $S_0$, and $x_0$ a point of $X_0$.  Then there exist an open neighbourhood $U_0$ of $x_0$ in $X_0$, an $S$-prescheme $U$ smooth \textup{(resp.}\ \'etale\textup{)} over $S$, and an $S_0$-isomorphism
\[
 U\times_SS_0\xrightarrow{\sim}U_0.
\]
\end{proposition}

\begin{proof}
Observe that if $X_0$ is \'etale over $S_0$ at $x_0$, it is \emph{a fortiori} unramified over $S_0$ at that point.  If an $S$-prescheme $U$ smooth over $S$ has been constructed such that $U\times_SS_0$ is isomorphic to $U_0$, then the fibres of the morphisms $U_0\to S_0$ and $U\to S$ containing $x_0$ are isomorphic.  It follows that $U$ is unramified over $S$ at $x_0$ \sref{IV.17.4.1},~\emph{d)}, hence also in a neighbourhood of $x_0$; replacing $U$ by this neighbourhood, we conclude that $U$ is \'etale over $S$.  It is therefore enough to prove the proposition when $X_0$ is merely supposed smooth over $S_0$.

The question being local on $S$ and on $X_0$, we may suppose that $S=\operatorname{Spec}(A)$ and $X_0=\operatorname{Spec}(C_0)$ are affine, so that $S_0=\operatorname{Spec}(A_0)$, where $A_0$ is a quotient ring of $A$, and
\[
 C_0=B_0/\mathfrak{J}_0,\qquad B_0=A_0[T_1,\ldots,T_n],
\]
where $\mathfrak{J}_0$ is an ideal of finite type in $B_0$; finally, $C_0$ is a formally smooth $A_0$-algebra for the discrete topologies.  Let $\mathfrak p_0$ be the ideal $\mathfrak j_{x_0}$ in $C_0$; then $\mathfrak p_0=\mathfrak q_0/\mathfrak{J}_0$, where $\mathfrak q_0$ is a prime ideal of $B_0$.  The Jacobian criterion ($\mathbf{0}$, 22.6.4), together with ($\mathbf{0}$, 19.1.12), shows that there exist in $\mathfrak{J}_0$ a family of $r$ polynomials $u_i$ $(1\leq i\leq r)$ and indices $j_h$ $(1\leq h\leq r)$ such that the images of the $u_i$ in
\[
 (\mathfrak{J}_0)_{\mathfrak q_0}/(\mathfrak{J}_0^2)_{\mathfrak q_0}
\]
generate this $(B_0)_{\mathfrak q_0}$-module and
\begin{equation}
 \det(\partial u_i/\partial T_{j_h})\notin\mathfrak q_0.
 \tag{18.1.1.1}\label{IV.18.1.1.1}
\end{equation}

Since $(B_0)_{\mathfrak q_0}$ is a local ring, Nakayama's lemma shows that we may suppose that the images of the $u_i$ in $(\mathfrak{J}_0)_{\mathfrak q_0}$ generate this $(B_0)_{\mathfrak q_0}$-module and then, replacing $X_0$ if necessary by an affine open neighbourhood of $x_0$, that the $u_i$ generate $\mathfrak{J}_0$ ($\mathbf{0}_{\mathbf I}$, 5.2.2).  Put $B=A[T_1,\ldots,T_n]$; then $B_0$ is a quotient ring
\oldpage[IV-4]{110}
of $B$, $\mathfrak q_0$ is the image of a prime ideal $\mathfrak q$ of $B$, and $\mathfrak q$ is the inverse image of $\mathfrak q_0$.  For every $i$, let $v_i\in B$ be an element whose image in $B_0$ is $u_i$, and let $\mathfrak{J}$ be the ideal of $B$ generated by the $v_i$, so that $\mathfrak{J}_0$ is the image of $\mathfrak{J}$ in $B_0$.

The proposition will be proved by taking for $U$ an open neighbourhood of the point of $\operatorname{Spec}(B/\mathfrak{J})$ corresponding to the prime ideal $\mathfrak p=\mathfrak q/\mathfrak{J}$, provided that we show that $(B/\mathfrak{J})_{\mathfrak q/\mathfrak{J}}$ is a formally smooth $A$-algebra for the discrete topologies.  This follows from the Jacobian criterion: the images of the $v_i$ in $\mathfrak{J}_{\mathfrak q}/\mathfrak{J}_{\mathfrak q}^2$ generate this $B_{\mathfrak q}$-module, and \eqref{IV.18.1.1.1} gives
\[
 \det\left(\frac{\partial v_i}{\partial T_{j_h}}\right)\notin\mathfrak q.
\]
\end{proof}

\begin{theorem}[18.1.2]
\label{IV.18.1.2}
Let $S$ be a prescheme, $S_0$ a closed sub-prescheme of $S$ whose underlying space is identical to that of $S$.
Then the functor
\[
X\leadsto X\times_S S_0
\]
from the category of $S$-preschemes \'etale over $S$, to the category of $S_0$-preschemes \'etale over $S_0$, is an equivalence of categories.
\end{theorem}

\begin{proof}
Let us first show that this functor is \emph{fully faithful}.
Let $X$, $Y$ be two $S$-preschemes \'etale over $S$, and set $X_0=X\times_S S_0$, $Y_0=Y\times_S S_0$.
If $Z=X\times_S Y$, the set $\operatorname{Hom}_S(X,Y)$ is in canonical bijective correspondence with the set of $X$-sections $\Gamma(Z/X)$, and likewise $\operatorname{Hom}_{S_0}(X_0,Y_0)$ is in canonical bijective correspondence with $\Gamma(Z_0/X_0)$, where $Z_0=Z\times_S S_0=X_0\times_{S_0}Y_0$.
Now $Z$ is \'etale over $X$, $Z_0$ is \'etale over $X_0$, and $X_0$ (resp.\ $Z_0$) is a closed subprescheme of $X$ (resp.\ $Z$) having the same underlying space.  The open subsets of $Z$ on which the restriction of $Z\to X$ is surjective and radicial are therefore the \emph{same} as the open subsets of $Z_0$ having the corresponding properties, and our assertion follows from \sref{IV.17.9.3}.

To complete the proof, it suffices to show that for every $S_0$-prescheme $X_0$ \'etale over $S_0$, there exists an $S$-prescheme $X$ \'etale over $S$ and an $S_0$-isomorphism $X_0\xrightarrow{\sim}X\times_S S_0$.
By virtue of Proposition~\sref{IV.18.1.1}, there exists an open covering $(U_\alpha)$ of $X_0$ and for each $\alpha$, an $S$-prescheme $V_\alpha$ which is \'etale over $S$, and finally an $S_0$-isomorphism $\theta_\alpha:U_\alpha\xrightarrow{\sim}V_\alpha\times_S S_0$.
Furthermore, by virtue of the first part of the proof, the identity of $U_\alpha\cap U_\beta$ lifts uniquely, via $\theta_\alpha$ and $\theta_\beta$, to an $S$-isomorphism $\varphi_{\alpha\beta}$ between the corresponding open subpreschemes of $V_\alpha$ and $V_\beta$, and it is immediate, for the same reason, that these isomorphisms satisfy the glueing condition \sref[0\textsubscript{\mdseries I}]{0.4.1.7}.
Hence there exists an $S$-prescheme $X$ such that the $V_\alpha$ identify canonically with sub-preschemes induced on the open subsets of $X$, the $\theta_\alpha$ identify with $S_0$-isomorphisms which coincide on the intersections $U_\alpha\cap U_\beta$, and define an $S_0$-isomorphism $X_0\xrightarrow{\sim}X\times_S S_0$.
It is clear that $X$ is \'etale over $S$ \sref{IV.17.3.2}, which completes the proof.
\end{proof}

\begin{corollary}[18.1.3]
\label{IV.18.1.3}
Let $S$ be a prescheme, $X$ an $S$-prescheme \'etale over $S$, $S'$ an $S$-prescheme, $S_0'$ a closed sub-prescheme of $S'$ having the same underlying space.
Then the canonical map $X(S')_S\to X(S_0')_S$ \sref[I]{I.3.4.3} is bijective.
\end{corollary}

\begin{proof}
Indeed, set $X'=X\times_S S'$, $X_0'=X\times_S S_0'$, so that $X(S')_S=\Gamma(X'/S')$ and $X(S_0')_S=\Gamma(X_0'/S_0')$;
the corollary follows from the fact that the functor defined in \sref{IV.18.1.2} (with $S$ and $S_0$ replaced by $S'$ and $S_0'$ respectively) is fully faithful (or directly from \sref{IV.17.9.3}).
\end{proof}

\subsection{\'Etale coverings.}
\label{subsection:IV.18.2}

\oldpage[IV-4]{111}

\begin{env}[18.2.1]
\label{IV.18.2.1}
Given a ring $A$ and a commutative $A$-algebra $B$ which is \emph{finite} and is a \emph{free} $A$-module, recall (Bourbaki, \emph{Alg.}, ch.~VIII, \textsection~12, no.~2) that one defines on $B$ an \emph{$A$-linear form} $\operatorname{Tr}_{B/A}$, the ``trace form''; from it one obtains a \emph{symmetric $A$-bilinear form} (also called the ``trace form'')
\[
\label{IV.18.2.1.1}
(x,y)\leadsto\operatorname{Tr}_{B/A}(xy)
\tag{18.2.1.1}
\]
whose datum is equivalent to that of the associated $A$-linear map $\operatorname{astr}_{B/A}:B\to\check{B}$ from the $A$-module $B$ into its \emph{dual} $\check{B}$; this map is equal to its transpose.  When $A$ is a \emph{field}, saying that this bilinear form is \emph{non-degenerate} is equivalent to saying that $B$ is a \emph{separable} $A$-algebra (Bourbaki, \emph{Alg.}, ch.~IX, \textsection~2, prop.~5).

Let $f:A\to A'$ be a ring homomorphism;
set $B'=B\otimes_A A'$, and let $g:B\to B'$ be the canonical homomorphism;
the image under $g$ of a basis of the $A$-module $B$ is then a basis of the $A'$-module $B'$, and it follows from the definitions that, for every $x\in B$,
\[
\label{IV.18.2.1.2}
\operatorname{Tr}_{B'/A'}(g(x))=f(\operatorname{Tr}_{B/A}(x)).
\tag{18.2.1.2}
\]
\end{env}

\begin{env}[18.2.2]
\label{IV.18.2.2}
Consider now a \emph{ringed space} $(X,\sh{O}_X)$ and let $\sh{B}$ be an $\sh{O}_X$-\emph{algebra} which, as an $\sh{O}_X$-\emph{Module}, is \emph{locally free of finite rank}.  For every open subset $U\subset X$ such that $\sh{B}|_U$ is, as an $\sh{O}_U$-Module, isomorphic to $\sh{O}_U^n$ (for some $n$ depending on $U$), $\Gamma(U,\sh{B})$ is a $\Gamma(U,\sh{O}_X)$-algebra which, as a $\Gamma(U,\sh{O}_X)$-module, is free of finite rank.  It therefore defines a $\Gamma(U,\sh{O}_X)$-linear form $\operatorname{Tr}_{\Gamma(U,\sh{B})/\Gamma(U,\sh{O}_X)}$, which we also denote by $\operatorname{Tr}_{\sh{B}/\sh{O}_X,U}$, and hence an associated linear map
\[
\operatorname{astr}_{\sh{B}/\sh{O}_X,U}:\Gamma(U,\sh{B})\to\Gamma(U,\sh{B})^\vee=\Gamma(U,\check{\sh{B}}).
\]

Furthermore, it follows from \sref{IV.18.2.1.2} that these linear maps are compatible with restriction from $U$ to a smaller open subset and therefore define, on the one hand, an \emph{$\sh{O}_X$-Module homomorphism}, again called the \emph{trace homomorphism},
\[
\label{IV.18.2.2.1}
\operatorname{Tr}_{\sh{B}/\sh{O}_X}:\sh{B}\to\sh{O}_X
\tag{18.2.2.1}
\]
and, on the other hand, an \emph{$\sh{O}_X$-Module homomorphism}
\[
\label{IV.18.2.2.2}
\operatorname{astr}_{\sh{B}/\sh{O}_X}:\sh{B}\to\check{\sh{B}}
\tag{18.2.2.2}
\]
called the \emph{homomorphism associated to the trace}, and equal to its transpose.
It also follows from \sref{IV.18.2.1.2} that for every $x\in X$, we have
\[
\label{IV.18.2.2.3}
(\operatorname{Tr}_{\sh{B}/\sh{O}_X})_x=\operatorname{Tr}_{\sh{B}_x/\sh{O}_{X,x}}
\tag{18.2.2.3}
\]
\[
\label{IV.18.2.2.4}
(\operatorname{astr}_{\sh{B}/\sh{O}_X})_x=\operatorname{astr}_{\sh{B}_x/\sh{O}_{X,x}}.
\tag{18.2.2.4}
\]

Finally, under the conditions of \sref{IV.18.2.1}, if we set $X=\operatorname{Spec}(A)$ and $\sh{B}=\widetilde{B}$, the form $\operatorname{Tr}_{\sh{B}/\sh{O}_X}$ (resp.\ the homomorphism $\operatorname{astr}_{\sh{B}/\sh{O}_X}$) corresponds to the form $\operatorname{Tr}_{B/A}$ (resp.\ the $A$-module homomorphism $\operatorname{astr}_{B/A}$), as also follows from \sref{IV.18.2.1.2}.
\end{env}

\begin{proposition}[18.2.3]
\label{IV.18.2.3}
Let $f:X\to Y$ be a \textup{finite} morphism of preschemes and set $\sh{B}=f_*(\sh{O}_X)$.
The following conditions are equivalent:
\begin{enumerate}
  \item[\textup{a)}] $f$ is \'etale.
  \item[\textup{a$'$)}] $f$ is a flat morphism of finite presentation, and for every $x\in X$, if $y=f(x)$, the ring $\sh{O}_{X,x}/\mathfrak{m}_y\sh{O}_{X,x}$ is a field which is a finite separable extension of $\boldsymbol{k}(y)$.
  \item[\textup{b)}] $\sh{B}$ is a locally free $\sh{O}_Y$-Module and, for every $y\in Y$, $\sh{B}_y\otimes_{\sh{O}_{Y,y}}\boldsymbol{k}(y)$ is a finite separable $\boldsymbol{k}(y)$-algebra \textup{(hence a direct product of finitely many fields which are finite separable extensions of $\boldsymbol{k}(y)$)}.
  \item[\textup{c)}] $\sh{B}$ is a locally free $\sh{O}_Y$-Module and the homomorphism $\operatorname{astr}_{\sh{B}/\sh{O}_Y}:\sh{B}\to\check{\sh{B}}$ \sref{IV.18.2.2} is bijective.
\end{enumerate}
\end{proposition}

\begin{proof}
\oldpage[IV-4]{112}
Since $f$ is quasi-compact, the equivalence of \emph{a)} and \emph{a$'$)} has already been shown \sref{IV.17.6.2}.
To prove the rest of the proposition, we may restrict to the case where $Y=\operatorname{Spec}(A)$ and $X=\operatorname{Spec}(B)$ are affine, $B$ being a finite $A$-algebra and $\sh{B}=\widetilde{B}$.
To say that $f$ is a morphism of finite presentation is equivalent to saying that $B$ is an $A$-\emph{module of finite presentation} \sref{IV.1.4.7}.
If, in addition, $f$ is flat, then $B$ is a flat $A$-module; it is known (Bourbaki, \emph{Alg.\ comm.}, ch.~II, \textsection~5, no.~2, cor.~2 to th.~1) that $B$ is then a \emph{projective} $A$-module.  Thus $\sh{B}$ is a locally free $\sh{O}_Y$-Module (\emph{loc.\ cit.}, no.~2, th.~1), and the converse is immediate.
On the other hand, $f^{-1}(y)$ is nothing but the spectrum of the $\boldsymbol{k}(y)$-algebra $\sh{B}(y)=\sh{B}_y\otimes_{\sh{O}_{Y,y}}\boldsymbol{k}(y)$, which completes the proof of the equivalence of \emph{a$'$)} and \emph{b)}.
To see that \emph{b)} is equivalent to \emph{c)}, note that the second assertion of \emph{b)} is equivalent to the bijectivity of
\[
 \operatorname{astr}_{\sh{B}(y)/\kres(y)}:\sh{B}(y)\longrightarrow\sh{B}(y)^\vee.
\]
Since
\[
 \sh{B}(y)=\sh{B}_y/\mathfrak{m}_y\sh{B}_y,
 \qquad
 \sh{B}(y)^\vee=\check{\sh{B}}_y/\mathfrak{m}_y\check{\sh{B}}_y,
\]
and $\sh{B}_y$ and $\check{\sh{B}}_y$ are free $\sh{O}_{Y,y}$-modules, it follows from \sref{IV.18.2.2.4} and Bourbaki, \emph{Alg.\ comm.}, ch.~II, \textsection~3, no.~3, cor.\ to prop.~6, that the homomorphism $\operatorname{astr}_{\sh{B}_y/\sh{O}_{Y,y}}:\sh{B}_y\to\check{\sh{B}}_y$ is itself bijective;
the converse being evident, this completes the proof.

When an $\sh{O}_X$-Algebra $\sh{B}$ satisfies the equivalent conditions \emph{b)} and \emph{c)} of \sref{IV.18.2.3}, we say that $\sh{B}$ is a \emph{finite \'etale $\sh{O}_X$-Algebra}.
When $X=\operatorname{Spec}(A)$ is affine and $\sh{B}=\widetilde{B}$, this amounts, by virtue of \sref{IV.18.2.3}, to saying that $B$ is a finite \'etale $A$-algebra, or that $B$ is a finite $A$-algebra that is \'etale (in the sense of \sref{IV.17.3.2}).
\end{proof}

\begin{corollary}[18.2.4]
\label{IV.18.2.4}
Let $f:X\to Y$ be a finite morphism of finite presentation, and set $\sh{B}=f_*(\sh{O}_X)$.
Let $y$ be a point of $Y$.
The following conditions are equivalent:
\begin{enumerate}[label=\emph{\alph*)}]
  \item There exists an open neighbourhood $U$ of $y$ in $Y$ such that the restriction $f^{-1}(U)\to U$ of $f$ is an \'etale morphism.
  \item $\sh{B}_y$ is a free $\sh{O}_{Y,y}$-module of finite type and $\sh{B}_y\otimes_{\sh{O}_{Y,y}}\boldsymbol{k}(y)$ is a separable $\boldsymbol{k}(y)$-algebra.
\end{enumerate}
\end{corollary}

\begin{proof}
It is clear that \emph{a)} implies \emph{b)} by virtue of \sref{IV.18.2.3}.
On the other hand, $\sh{B}$ is an $\sh{O}_Y$-Module of finite presentation \sref{IV.1.6.3} and \sref{IV.1.4.7}; hence, if $\sh{B}_y$ is a free $\sh{O}_{Y,y}$-module, there exists an open neighbourhood $U$ of $y$ in $Y$ such that $\sh{B}|_U$ is a locally free $\sh{O}_U$-Module ($\mathbf{0}_{\mathbf I}$, 5.2.7).  Moreover, since by hypothesis the homomorphism $\operatorname{astr}_{\sh{B}_y/\sh{O}_{Y,y}}:\sh{B}_y\to\check{\sh{B}}_y$ is bijective, the same result shows that $U$ may be chosen so that
\oldpage[IV-4]{113}
the homomorphism $\operatorname{astr}_{\sh{B}|_U/\sh{O}_U}$ is bijective.
The fact that \emph{b)} implies \emph{a)} then follows from \sref{IV.18.2.3}.
\end{proof}

\begin{corollary}[18.2.5]
\label{IV.18.2.5}
Let $Y$ be a quasi-compact or locally Noetherian prescheme, and let $f:X\to Y$ be a finite morphism of finite presentation;
set $\sh{B}=f_*(\sh{O}_X)$.
Suppose that for every \emph{closed} point $y$ of $Y$, $\sh{B}_y$ is a free $\sh{O}_{Y,y}$-module and $\sh{B}_y\otimes_{\sh{O}_{Y,y}}\boldsymbol{k}(y)$ is a separable $\boldsymbol{k}(y)$-algebra.
Then $f$ is \'etale.
\end{corollary}

\begin{proof}
Indeed, it follows from \sref{IV.18.2.4} that every closed point $y$ of $Y$ has an open neighbourhood $U$ such that the restriction $f^{-1}(U)\to U$ of $f$ is \'etale;
the conclusion follows from the fact that in the two cases considered, every non-empty closed subset of $Y$ contains a closed point (\sref{IV.5.1.11} and \sref[0\textsubscript{\mdseries I}]{0.2.1.3}).
\end{proof}

\begin{corollary}[18.2.6]
\label{IV.18.2.6}
If $f:X\to Y$ is a surjective finite \'etale morphism, and if we set $\sh{B}=f_*(\sh{O}_X)$, then the $\sh{O}_Y$-homomorphism $\operatorname{Tr}_{\sh{B}/\sh{O}_Y}:\sh{B}\to\sh{O}_Y$ \sref{IV.18.2.2} \textup{(also denoted $\operatorname{Tr}_f$)} is surjective.
\end{corollary}

\begin{proof}
The question being local, we may, by virtue of \sref{IV.18.2.3}, suppose that $Y=\operatorname{Spec}(A)$, $X=\operatorname{Spec}(B)$, with $B$ a free $A$-module.
The trace pairing, perfect by \sref{IV.18.2.3}, identifies $B$ with $\check{B}$ and sends $1$ to $\operatorname{Tr}_{B/A}$; since $f$ is surjective, $1$ is unimodular, so the trace is surjective.
\end{proof}

\begin{remarks}[18.2.7]
\label{IV.18.2.7}
\begin{enumerate}
  \item[(i)] When $f:X\to Y$ is a finite morphism such that $f_*(\sh{O}_X)$ is a locally free $\sh{O}_Y$-Module (resp.\ locally free of rank $n$), we say further that $f$ is a \emph{finite locally free} morphism (resp.\ \emph{locally free of rank $n$}).
This condition, by virtue of \sref{IV.18.2.3}, is satisfied if $f$ is a finite \'etale morphism, but does not imply by itself that $f$ is \'etale, as shown by the example where $X=\operatorname{Spec}(K)$ and $Y=\operatorname{Spec}(k)$ are spectra of fields, $K$ being a finite non-separable extension of $k$.
When $f:X\to Y$ is a finite \'etale morphism, we say further that $X$ is an \emph{\'etale covering} of $Y$.
One notes that in this case, $f$ is universally open and universally closed, and in particular $f(X)$ is a subset that is both open and closed of $Y$.

We say that an \'etale covering $X$ of $Y$ is \emph{trivial} if $X$ is a \emph{sum} of a finite number of preschemes \emph{isomorphic to} $Y$.
We say that an \'etale covering $X$ of $Y$ is \emph{locally trivial} if the morphism $f:X\to Y$ is such that for every point $y\in Y$ there is an open neighbourhood $U$ for which the covering $f^{-1}(U)$ of $U$ is trivial.

  \item[(ii)] Let $f:X\to Y$ be a finite morphism, locally free of rank $n$;
set $\sh{B}=f_*(\sh{O}_X)$ and let $u=\operatorname{astr}_{\sh{B}/\sh{O}_Y}:\sh{B}\to\check{\sh{B}}$;
we deduce from this an \emph{$n$-th} exterior power homomorphism
$\bigwedge^n u:\bigwedge^n\sh{B}\to\bigwedge^n\check{\sh{B}}=(\bigwedge^n\sh{B})^\vee$ of \emph{invertible} $\sh{O}_Y$-Modules, and by \sref[0\textsubscript{\mdseries I}]{0.5.4.2} an element
\[
\label{IV.18.2.7.1}
d_{X/Y}\in\Gamma(Y,(\bigwedge^n\check{\sh{B}})\otimes_{\sh{O}_Y}(\bigwedge^n\check{\sh{B}}))
\tag{18.2.7.1}
\]
which we call the \emph{discriminant} of $X$ over $Y$.
Furthermore, since $(\bigwedge^n\check{\sh{B}})\otimes_{\sh{O}_Y}(\bigwedge^n\check{\sh{B}})$ is the dual of $(\bigwedge^n\sh{B})\otimes_{\sh{O}_Y}(\bigwedge^n\sh{B})$, $d_{X/Y}$ can also be identified with a homomorphism
\[
\label{IV.18.2.7.2}
(\bigwedge^n\sh{B})\otimes_{\sh{O}_Y}(\bigwedge^n\sh{B})\to\sh{O}_Y
\tag{18.2.7.2}
\]
and we denote by $\mathscr{D}_{X/Y}$ the quasi-coherent ideal of $\sh{O}_Y$ of finite type, \emph{image} of the homomorphism \sref{IV.18.2.7.2}, also called the \emph{discriminant ideal} of $X$ over $Y$.
\end{enumerate}
\end{remarks}

\oldpage[IV-4]{114}
For the homomorphism $u$ to be bijective, it is necessary and sufficient that $\bigwedge^n u$ be bijective, or again that the section $d_{X/Y}$ have an \emph{invertible} germ at every point $y\in Y$, which is also written $d_{X/Y}(y)\neq 0$ for every $y$ \sref[0\textsubscript{\mdseries I}]{0.5.5.2}.
It amounts to the same to say that the discriminant ideal $\mathscr{D}_{X/Y}$ is \emph{equal} to $\sh{O}_Y$.

The terminology ``covering'' introduced in \sref{IV.18.2.7}, (i), is justified by the following proposition:

\begin{proposition}[18.2.8]
\label{IV.18.2.8}
Let $f:X\to Y$ be a morphism \'etale, separated and of finite type, and for every $y\in Y$, let $n(y)$ be the geometric number of points of $f^{-1}(y)$.
Then the function $y\leadsto n(y)$ is lower semi-continuous on $Y$.
For it to be continuous at a point $y$ \textup{(and hence \emph{constant} in a neighbourhood of $y$)}, it is necessary and sufficient that there exists an open neighbourhood $U$ of $y$ such that the restriction $f^{-1}(U)\to U$ of $f$ is a \emph{finite} (\'etale) morphism.
\end{proposition}

\begin{proof}
Since $f$ is quasi-finite \sref{IV.17.6.1} and locally of finite presentation, it amounts to saying that $f$ is finite or that $f$ is proper \sref{IV.8.11.1}:
moreover, each fibre $f^{-1}(y)$ is geometrically reduced over $\boldsymbol{k}(y)$.
The conclusions therefore follow from \sref{IV.15.5.9}, (i) and (ii) and from the fact that $f$ is flat.
\end{proof}

\begin{corollary}[18.2.9]
\label{IV.18.2.9}
Let $Y$ be a connected prescheme, $f:X\to Y$ a morphism \'etale, separated and of finite type.
For $f$ to be finite \textup{(in other words, for $X$ to be an \emph{\'etale covering} of $Y$)}, it is necessary and sufficient that all the fibres of $f$ have the same geometric number of points.
\end{corollary}

\begin{remarks}[18.2.10]
\label{IV.18.2.10}
\begin{enumerate}
  \item[(i)] The example of the ``affine line with a doubled point'' \sref[I]{I.5.5.11} shows that a morphism \'etale of finite type of Noetherian preschemes may not be separated;
for this example, the first assertion of \sref{IV.18.2.8} is no longer exact.

  \item[(ii)] For a separated morphism of finite type and \'etale $f:X\to Y$ to make $X$ a \emph{locally trivial} covering, it is necessary and sufficient that for every $x\in X$ there exist an open neighbourhood $U$ of $y=f(x)$ and a $U$-section $g$ of $f^{-1}(U)$ such that $g(y)=x$.
Indeed, the condition is obviously necessary;
the fact that it is also sufficient follows from the fact that every fibre $f^{-1}(y)$ is finite (\sref{IV.17.6.1} and \sref[I]{I.6.2.2}), from the characterisation of sections of a $Y$-scheme \'etale \sref{IV.17.9.3} and from Proposition~\sref{IV.18.2.8}.
\end{enumerate}
\end{remarks}


\subsection{Finite algebras and \'etale algebras.}
\label{subsection:IV.18.3}

\begin{proposition}[18.3.1]
\label{IV.18.3.1}
Let $A$ be a ring, $B$ an $A$-algebra of finite presentation.
\begin{enumerate}
  \item[\textup{(i)}] For $B$ to be an unramified $A$-algebra, it is necessary and sufficient that $B$ be a \emph{projective} $(B\otimes_A B)$-module.
  \item[\textup{(ii)}] Suppose moreover that $B$ is a finite $A$-algebra.  For $B$ to be an \'etale $A$-algebra, it is necessary and sufficient that $B$ be a projective $A$-module and a projective $(B\otimes_A B)$-module.
\end{enumerate}
\end{proposition}

\begin{proof}
\oldpage[IV-4]{115}
The $(B\otimes_A B)$-module structure on $B$ is that coming from the canonical ring homomorphism $B\otimes_A B\to B$, which is surjective and has kernel $\mathfrak{J}_{B/A}$ \sref[0]{0.20.4.1}.

(i) To say that the morphism $\operatorname{Spec}(B)\to\operatorname{Spec}(A)$ is locally of finite presentation is equivalent to saying that $B$ is an $A$-algebra of finite presentation \sref{IV.1.4.6}.  To say that $B$ is an unramified $A$-algebra means \sref{IV.17.4.2} that
\[
 \operatorname{Spec}\bigl((B\otimes_A B)/\mathfrak{J}_{B/A}\bigr)
\]
is the subprescheme induced on an open and closed subset of $\operatorname{Spec}(B\otimes_A B)$.  This holds if and only if $\mathfrak{J}_{B/A}$ is a \emph{direct-factor ideal} of $B\otimes_A B$ (Bourbaki, \emph{Alg.\ comm.}, ch.~II, \textsection~4, no.~3, prop.~15);
but this amounts to saying that the $(B\otimes_A B)$-module quotient $(B\otimes_A B)/\mathfrak{J}_{B/A}$ is projective (Bourbaki, \emph{Alg.}, ch.~II, $3^{\text{e}}$ \'ed., \textsection~2, no.~2, prop.~4).

(ii) Recall that a flat $A$-module of finite presentation is projective, and conversely (Bourbaki, \emph{Alg.\ comm.}, ch.~II, \textsection~5, no.~2, cor.~2 du th.~1), so the assertion (ii) follows from that of (i) and from \sref{IV.17.6.2}.
\end{proof}

\begin{proposition}[18.3.2]
\label{IV.18.3.2}
Let $A$ be a ring, $\mathfrak{J}$ an ideal of $A$ such that, for the $\mathfrak{J}$-preadic topology, $A$ is separated and complete;
set $A_0=A/\mathfrak{J}$.
Then the functor
\[
B\leadsto B\otimes_A A_0
\]
is an equivalence of the category of finite \'etale $A$-algebras, to the category of finite \'etale $A_0$-algebras.
\end{proposition}

\begin{proof}
We will first prove the following lemma:

\begin{lemma}[18.3.2.1]
\label{IV.18.3.2.1}
Let $A$ be a ring, $\mathfrak{J}$ an ideal of $A$ such that, for the $\mathfrak{J}$-preadic topology, $A$ is separated and complete.
\begin{enumerate}
  \item[\textup{(i)}] Every projective $A$-module of finite type $M$ is separated and complete for the $\mathfrak{J}$-preadic topology, and hence is the projective limit of the projective $(A/\mathfrak{J}^{n+1})$-modules $M/\mathfrak{J}^{n+1}M$.
  \item[\textup{(ii)}] Conversely, set $A_n=A/\mathfrak{J}^{n+1}$, and let $(M_n)$ be a projective system of $A_n$-modules such that, for every $n$, the homomorphism $M_{n+1}\otimes_{A_{n+1}}A_n\to M_n$ induced by the transition homomorphism $M_{n+1}\to M_n$ is bijective.
Suppose moreover that the $M_n$ are projective and $M_0$ of finite type.
Then $M=\varprojlim M_n$ is a projective $A$-module of finite type such that the canonical homomorphism $M\otimes_A A_0\to M_0$ is bijective.
\end{enumerate}
\end{lemma}

\begin{proof}
(i) There is a free $A$-module of finite type $L$ such that $M$ is isomorphic to a direct factor of $L$;
as $L$ is separated for the $\mathfrak{J}$-preadic topology, so is every sub-module $N$ of $L$ since $\mathfrak{J}^{n+1}N\subset\mathfrak{J}^{n+1}L$;
in particular $M$ is separated for this topology, and since the surjective homomorphism $f:L\to M$ is continuous for the $\mathfrak{J}$-preadic topology, its kernel $N$ is closed for the topology induced by that of $L$;
since $L$ is complete and $f$ a strict morphism, we conclude that $M$ is complete (Bourbaki, \emph{Top.\ g\'en.}, ch.~IX, $2^{\text{e}}$ \'ed., \textsection~3, no.~1, prop.~4).

(ii) It follows from Nakayama's lemma that, if $M_0$ is generated by a finite family $(x_{i0})$ of $r$ elements and if, for every $n\geq 1$, $x_{in}$ is an element of $M_n$ whose image in $M_{n-1}$ is $x_{i,n-1}$, then $(x_{in})$ $(1\leq i\leq r)$ is a system of generators of $M_n$ (Bourbaki, \emph{Alg.\ comm.}, ch.~II, \textsection~3, cor.~2 to prop.~4).  For every $n$, put $L_n=A_n^r$.  If $(e_{in})_{1\leq i\leq r}$ is the canonical basis of $L_n$, let $u_n:L_n\to M_n$ be the $A$-linear map such that $u_n(e_{in})=x_{in}$ for every $i$.
By hypothesis, we have a \emph{split} exact sequence
\[
0\to N_n\xrightarrow{v_n} L_n\xrightarrow{u_n} M_n\to 0
\]
\oldpage[IV-4]{116}
and since $L_n=L_{n+1}/\mathfrak{J}^{n+1}L_{n+1}$ and $M_n=M_{n+1}/\mathfrak{J}^{n+1}M_{n+1}$, the vertical arrows in the commutative diagram
\[
\xymatrix{
0\ar[r] & N_{n+1}\ar[r]^{v_{n+1}}\ar[d] & L_{n+1}\ar[r]^{u_{n+1}}\ar[d] & M_{n+1}\ar[r]\ar[d] & 0\\
0\ar[r] & N_n\ar[r]_{v_n} & L_n\ar[r]_{u_n} & M_n\ar[r] & 0
}
\]
are all three \emph{surjective}.
Now we have $M=\varprojlim M_n$ and $L=A^r=\varprojlim L_n$;
setting $N=\varprojlim N_n$, $u=\varprojlim u_n$, $v=\varprojlim v_n$, by \sref[0\textsubscript{\mdseries III}]{0.13.2.2} the exact sequence
\[
\label{IV.18.3.2.2}
0\to N\xrightarrow{v} L\xrightarrow{u} M\to 0
\tag{18.3.2.2}
\]
is exact.
Moreover, since for every $n$, $v_n$ is left-invertible and $M_n$ is a projective $A_n$-module, by virtue of \sref[0]{0.19.1.8} the exact sequence \sref{IV.18.3.2.2} is split, which proves the lemma.
\end{proof}

This being established, let us first show that the functor of \sref{IV.18.3.2} is \emph{fully faithful}.
Set, as in the lemma, $A_n=A/\mathfrak{J}^{n+1}$;
let $B$, $C$ be two finite \'etale $A$-algebras, and set, for every $n$, $B_n=B\otimes_A A_n$, $C_n=C\otimes_A A_n$;
by virtue of \sref{IV.18.3.1} and \sref{IV.18.3.2.1}, $B$ and $C$ are separated and complete for the $\mathfrak{J}$-preadic topology, and we have $B=\varprojlim B_n$, $C=\varprojlim C_n$;
moreover, every $A$-algebra homomorphism $u:B\to C$ is continuous for the $\mathfrak{J}$-preadic topologies, and hence gives a projective system of $A_n$-algebra homomorphisms $u_n=u\otimes 1:B_n\to C_n$, of which it is the projective limit;
the converse being evident, we have a canonical bijection
\[
\operatorname{Hom}_{A\text{-}\mathrm{alg.}}(B,C)\xrightarrow{\sim}\varprojlim\operatorname{Hom}_{A_n\text{-}\mathrm{alg.}}(B_n,C_n).
\]
But since $B$ and $C$ are finite \emph{\'etale} $A$-algebras, it follows also from \sref{IV.18.1.2} that the canonical map
\[
\operatorname{Hom}_{A_{n+1}\text{-}\mathrm{alg.}}(B_{n+1},C_{n+1})\to\operatorname{Hom}_{A_n\text{-}\mathrm{alg.}}(B_n,C_n)
\]
is \emph{bijective} for $n\geq 0$, which proves that the canonical map $\operatorname{Hom}_{A\text{-}\mathrm{alg.}}(B,C)\to\operatorname{Hom}_{A_0\text{-}\mathrm{alg.}}(B_0,C_0)$ is bijective.

To complete the proof of \sref{IV.18.3.2}, it suffices to show that for every finite \'etale $A_0$-algebra $B_0$, there exists a finite \'etale $A$-algebra $B$ and an $A_0$-isomorphism $B_0\xrightarrow{\sim}B\otimes_A A_0$.
Now, by \sref{IV.18.1.2}, there is a projective system $(B_n)$ such that $B_n$ is an \'etale $A_n$-algebra and that the homomorphisms $B_{n+1}\otimes_{A_{n+1}}A_n\to B_n$ are bijective.
It follows from \sref{IV.18.3.1} and \sref{IV.18.3.2.1} that the $A$-algebra $B=\varprojlim B_n$ is a projective $A$-module of finite type and that $B_0$ is isomorphic to $B\otimes_A A_0$.
To prove that $B$ is a finite \'etale $A$-algebra, it suffices, by \sref{IV.18.2.5}, to show that for every maximal ideal $\mathfrak{m}$ of $A$, $B_\mathfrak{m}/\mathfrak{m}B_\mathfrak{m}$ is a separable $(A/\mathfrak{m})$-algebra.
Now, since $\mathfrak{J}$ is contained in the radical of $A$ \sref[0\textsubscript{\mdseries I}]{0.7.1.10},
\oldpage[IV-4]{117}
we have $\mathfrak{J}\subset\mathfrak{m}$, so $\mathfrak{m}_0=\mathfrak{m}/\mathfrak{J}$, we have $A_0/\mathfrak{m}_0=A/\mathfrak{m}$ and $B_\mathfrak{m}/\mathfrak{m}B_\mathfrak{m}=(B_0)_{\mathfrak{m}_0}/\mathfrak{m}_0(B_0)_{\mathfrak{m}_0}$;
the conclusion therefore follows from the fact that $B_0$ is a finite \'etale $A_0$-algebra \sref{IV.18.2.5}.
\end{proof}

\begin{env}[18.3.3]
\label{IV.18.3.3}
\emph{Example.} ---
Proposition \sref{IV.18.3.2} applies in particular when $A$ is a \emph{local} ring, separated and complete, $\mathfrak{J}$ being the maximal ideal of $A$, so that $A_0$ is a field and the finite \'etale $A_0$-algebras are precisely the finite separable $A_0$-algebras, hence the direct products of finite separable field extensions of $A_0$.
In particular, if the field $A_0$ is \emph{separably closed}, these extensions are all identical to $A_0$, and hence every \'etale covering of $\operatorname{Spec}(A)$ is \emph{trivial} \sref{IV.18.2.7} by virtue of \sref{IV.18.3.2}.
\end{env}

\begin{theorem}[18.3.4]
\label{IV.18.3.4}
Let $A$ be a Noetherian ring, $\mathfrak{J}$ an ideal of $A$ such that $A$ is separated and complete for the $\mathfrak{J}$-preadic topology, $A_0=A/\mathfrak{J}$.
Let $S=\operatorname{Spec}(A)$, $S_0=\operatorname{Spec}(A_0)$.
Let $X$ be a proper $S$-scheme, and set $X_0=X\times_S S_0$.
Then the functor
\[
Z\leadsto Z\times_X X_0
\]
from the category of finite \'etale $X$-schemes over $X$, to the category of finite \'etale $X_0$-schemes over $X_0$, is an equivalence of categories.
\end{theorem}

\begin{proof}
Let us first show that this functor is \emph{fully faithful}.
Let $Z'$ and $Z''$ be two finite \'etale $X$-schemes over $X$.
Set $S_n=\operatorname{Spec}(A/\mathfrak{J}^{n+1})$, $X_n=X\times_S S_n$, $Z'_n=Z'\times_S S_n$, $Z''_n=Z''\times_S S_n$.
It follows from \sref[III]{III.5.4.1} that we have a canonical bijection $\operatorname{Hom}_X(Z',Z'')\xrightarrow{\sim}\varprojlim\operatorname{Hom}_{X_n}(Z'_n,Z''_n)$.
Now, by virtue of \sref{IV.18.1.2}, the canonical map $\operatorname{Hom}_{X_{n+1}}(Z'_{n+1},Z''_{n+1})\to\operatorname{Hom}_{X_n}(Z'_n,Z''_n)$ is bijective, which proves our assertion.

It remains to show that if $\sh{B}_0$ is a finite \'etale $\sh{O}_{X_0}$-algebra, there exists a finite \'etale $\sh{O}_X$-algebra $\sh{B}$ and an isomorphism $\sh{B}_0\xrightarrow{\sim}\sh{B}\otimes_{\sh{O}_X}\sh{O}_{X_0}$.
It follows from \sref{IV.18.1.2} that there is a projective system $(\sh{B}_n)$, where $\sh{B}_n$ is a finite \'etale $\sh{O}_{X_n}$-algebra, and the second comparison theorem \sref[III]{III.5.1.4} proves the existence of a coherent $\sh{O}_X$-Module $\sh{B}$ and of a projective system of isomorphisms $\sh{B}_n\xrightarrow{\sim}\sh{B}\otimes_{\sh{O}_X}\sh{O}_{X_n}$.
The datum of an algebra structure on an $\sh{O}_X$-Module $\sh{F}$ is equivalent to that of a unit homomorphism $\sh{O}_X\to\sh{F}$ and a multiplication homomorphism $\sh{F}\otimes_{\sh{O}_X}\sh{F}\to\sh{F}$ making certain diagrams involving only tensor powers of $\sh{F}$ commute;
by \sref[III]{III.5.1.3}, \sref[I]{I.10.11.4} and \sref[I]{I.10.11.7}, $\sh{B}$ is therefore endowed in a natural way with a structure of $\sh{O}_X$-algebra for which the isomorphism $\sh{B}_0\xrightarrow{\sim}\sh{B}\otimes_{\sh{O}_X}\sh{O}_{X_0}$ is an isomorphism of algebras.
Moreover, $\sh{B}$ is a \emph{locally free} $\sh{O}_X$-Module;
this also follows from \sref[III]{III.5.1.3}, \sref[I]{I.10.11.4}, \sref[I]{I.10.11.7} and from the fact that, in the category of coherent $\sh{O}_X$-Modules, the locally free $\sh{O}_X$-Modules $\sh{F}$ can be characterized as those for which the functors $\sh{G}\leadsto\shHom_{\sh{O}_X}(\sh{F},\sh{G})$ are \emph{exact}.
Finally, to see that $\sh{B}$ is a finite \'etale $\sh{O}_X$-Algebra, it suffices \sref{IV.18.2.5} to show that for every \emph{closed} point $x\in X$, $\sh{B}_x\otimes_{\sh{O}_{X,x}}\boldsymbol{k}(x)$ is a separable $\boldsymbol{k}(x)$-algebra.
But since the structural morphism $f:X\to S$ is \emph{proper}, $f(x)$ is a closed point of $S$, so it belongs to $S_0$, since $\mathfrak{J}$ is contained in the radical of $A$ \sref[0\textsubscript{\mdseries I}]{0.7.1.10};
the conclusion therefore follows from the fact that $X_0=f^{-1}(S_0)$ and that $\sh{B}_0$ is a finite \'etale $\sh{O}_{X_0}$-algebra.
\end{proof}


\subsection{Local structure of unramified morphisms and of \'etale morphisms.}
\label{subsection:IV.18.4}

\oldpage[IV-4]{118}

\begin{lemma}[18.4.1]
\label{IV.18.4.1}
Let $A$ be a ring, $B$ a finite monogenic $A$-algebra, $u$ a generator of the $A$-algebra $B$, $F\in A[T]$ a polynomial such that $F(u)=0$, $F'$ the derivative polynomial;
set $u'=F'(u)$.
Then the annihilator ideal of $\Omega^1_{B/A}$ in $B$ contains $u'B$;
it is equal to $u'B$ if the ideal $\mathfrak{J}$ of polynomials $G\in A[T]$ such that $G(u)=0$ is generated by $F$, in other words if the canonical surjection $\varphi:A[T]/F A[T]\to B$, which sends the class of $T$ to $u$, is bijective.
\end{lemma}

\begin{proof}
Set $C=A[T]$, so that $B=C/\mathfrak{J}$.
We have the exact sequence \sref[0]{0.20.5.12.1}
\[
\mathfrak{J}/\mathfrak{J}^2\to\Omega^1_{C/A}\otimes_C B\to\Omega^1_{B/A}\to 0
\]
and $\Omega^1_{B/A}$ therefore identifies with the quotient $B/\mathfrak{J}'$, where $\mathfrak{J}'$ is the ideal generated by the elements $G'(u)$, where $G$ ranges over a system of generators of the ideal $\mathfrak{J}$ \sref[0]{0.20.5.13};
whence the lemma immediately follows.
\end{proof}

\begin{proposition}[18.4.2]
\label{IV.18.4.2}
Let the notation be those of \sref{IV.18.4.1}, let $\mathfrak{q}$ be a prime ideal of $B$.
Then:
\begin{enumerate}
  \item[\textup{(i)}] If $\mathfrak{q}$ does not contain $u'$, $B_\mathfrak{q}$ is a formally unramified $A_\mathfrak{p}$-algebra \textup{(}$\mathfrak{p}$ being the inverse image of $\mathfrak{q}$ in $A$\textup{)};
in other terms, $\operatorname{Spec}(B_{u'})$ is formally unramified over $\operatorname{Spec}(A)$.
  \item[\textup{(ii)}] Suppose moreover that $F$ is monic and generates $\mathfrak{J}$.
Then, for $\operatorname{Spec}(B)$ to be \'etale over $\operatorname{Spec}(A)$ at the point $\mathfrak{q}$, it is necessary and sufficient that $u'\notin\mathfrak{q}$.
\end{enumerate}
\end{proposition}

\begin{proof}
The hypothesis that $u'\notin\mathfrak{q}$ implies that $\Omega^1_{B_\mathfrak{q}/A_\mathfrak{p}}=0$ \sref[0]{0.20.5.9}, hence (i) follows from \sref{IV.17.2.1}.
Moreover, under the hypotheses of (ii), $B$ is a free $A$-module by virtue of the Euclidean division;
as the annihilator $\mathfrak{J}'$ of $\Omega^1_{B/A}$ is then equal to $u'B$ by virtue of \sref{IV.18.4.1} and $\Omega^1_{B/A}$ is a $B$-module of finite presentation \sref{IV.16.4.22}, the annihilator of $\Omega^1_{B_\mathfrak{q}/A_\mathfrak{p}}$ is equal to $u'B_\mathfrak{q}$ (Bourbaki, \emph{Alg.\ comm.}, ch.~II, \textsection~2, no.~4, formule (g)), and (ii) follows from (i) and from the implication $\textit{c)}\Rightarrow\textit{a)}$ in \sref{IV.17.6.1}.
\end{proof}

\begin{corollary}[18.4.3]
\label{IV.18.4.3}
With the notation of \sref{IV.18.4.2}, suppose that $F$ is monic and generates $\mathfrak{J}$.
Then, for $B$ to be an \'etale $A$-algebra, it is necessary and sufficient that $u'$ be invertible in $B$ \textup{(or, what amounts to the same, that the ideal of $A[T]$ generated by $F$ and $F'$ be equal to $A[T]$)}.
\end{corollary}

\begin{proof}
In view of \sref{IV.18.4.2},~(ii), saying that $\operatorname{Spec}(B)$ is \'etale over $\operatorname{Spec}(A)$ means that $u'$ does not belong to any prime ideal of $B$, i.e.\ that it is invertible in $B$.
\end{proof}

We say that a monic polynomial $F\in A[T]$ such that the ideal of $A[T]$ generated by $F$ and $F'$ is equal to $A[T]$ is \emph{separable};
it is immediate that this definition coincides with the usual definition (Bourbaki, \emph{Alg.}, ch.~V, \textsection~7, no.~6) when $A$ is a field.

\begin{lemma}[18.4.4]
\label{IV.18.4.4}
Let $A$ be a local ring, $B$ a finite monogenic $A$-algebra, $u$ a generator of the $A$-algebra $B$.
Let $\mathfrak{n}_i$ \textup{(}$1\leq i\leq r$\textup{)} be the maximal ideals of $B$ such that $\operatorname{Spec}(B)$ is formally unramified over $\operatorname{Spec}(A)$ at the points $\mathfrak{n}_i$.
Then there exists a monic polynomial $F\in A[T]$ such that $F(u)=0$ and $F'(u)\notin\mathfrak{n}_i$ for every $i$ \textup{(}$1\leq i\leq r$\textup{)}.
Moreover, if $k$ is the
\oldpage[IV-4]{119}
residue field of $A$ and $f\in k[T]$ is the minimal polynomial of the image of $u$ in $B\otimes_A k$, one can choose $F$ so that its canonical image in $k[T]$ is $f$.
\end{lemma}

\begin{proof}
The maximal ideal $\mathfrak{m}$ of $A$ is the inverse image of each $\mathfrak{n}_i$ \sref[II]{II.6.1.10}.
Set $L=B\otimes_A k$, a finite $k$-algebra, and let $\xi$ be the image of $u$ in $L$; if $n$ is the rank of $L$ over $k$, the minimal polynomial $f\in k[T]$ of $\xi$ over $k$ has degree $n$, and $L$ is isomorphic to $k[T]/f k[T]$.
If $\mathfrak{n}'_i=\mathfrak{n}_i/\mathfrak{m}B$, the hypothesis implies that $\operatorname{Spec}(L)$ is \'etale over $k$ at the points $\mathfrak{n}'_i$ ($1\leq i\leq r$), by \sref{IV.17.6.1},~\emph{c)}; hence $f'(\xi)\notin\mathfrak{n}'_i$ by \sref{IV.18.4.2},~(ii).
Now $\mathfrak{m}B$ is contained in the radical of $B$; since $1,\xi,\ldots,\xi^{n-1}$ form a basis of $L$ over $k$, Nakayama's lemma shows that $1,u,\ldots,u^{n-1}$ generate the $A$-module $B$.
Consequently there exists a monic polynomial $F\in A[T]$ of degree $n$ such that $F(u)=0$; since $\xi$ is a root of the canonical image of $F$ in $k[T]$, this image is necessarily $f$.
The image of $F'(u)$ in $L$ is then $f'(\xi)$, and therefore $F'(u)\notin\mathfrak{n}_i$ for every $i$.
\end{proof}

\begin{proposition}[18.4.5]
\label{IV.18.4.5}
Let $A$ be a local ring, $k$ its residue field, and $B$ a finite $A$-algebra \textup{(respectively,} a finite $A$-algebra of finite presentation\textup{)}.
Suppose moreover that the residue field $k$ is \emph{infinite}, or else that $B$ is a local ring.
Let $n$ be the rank of $L=B\otimes_A k$ over $k$.
For $B$ to be a formally unramified $A$-algebra \textup{(resp.}\ \'etale\textup{)}, it is necessary and sufficient that there exist a monic separable polynomial $F\in A[T]$ such that $B$ is isomorphic to a quotient of $A[T]/F.A[T]$ \textup{(resp.}\ isomorphic to $A[T]/F.A[T]$\textup{)}.
Moreover, one can suppose that $F$ has degree $n$ \textup{(respectively,} $F$ necessarily has degree $n$\textup{)}.
\end{proposition}

\begin{proof}
The conditions are sufficient by virtue of \sref{IV.18.4.2}, without assuming either that $k$ is infinite or that $B$ is local.
To see that the conditions are necessary, note that if $B$ is a formally unramified $A$-algebra, $L$ is a finite separable $k$-algebra, hence a direct product of finitely many finite separable field extensions $k_j$ of $k$ ($1\leq j\leq r$), each generated by an element $\bar{\xi}_j$ with minimal polynomial $f_j$ (Bourbaki, \emph{Alg.}, ch.~V, \textsection~11, no.~4, prop.~4).
Let us show that by virtue of the hypotheses made on $k$ or $B$, there exists an element $\bar{\xi}$ of $L$ generating the $k$-algebra $L$.
This is immediate if $B$ is local, since then $r=1$.
Otherwise, $k$ being assumed infinite, one can suppose that the irreducible polynomials $f_j\in k[T]$ are all \emph{distinct} (by replacing each $\bar{\xi}_j$ by $\bar{\xi}_j+a_j$, for a suitable element $a_j\in k$);
if we set $f=f_1 f_2\ldots f_r$, it is clear that $L$ is isomorphic to $k[T]/f k[T]$ in both cases, hence generated by an element $\bar{\xi}$ of minimal polynomial $f\in k[T]$ of degree $n$.
Choose $u\in B$ whose image in $L$ is $\bar{\xi}$.
Since $\mathfrak{m}B$ is contained in the radical of $B$ and $1,\bar{\xi},\ldots,\bar{\xi}^{n-1}$ form a basis of $L$ over $k$, Nakayama's lemma shows that $1,u,\ldots,u^{n-1}$ generate the $A$-module $B$; hence there exists a monic polynomial $F\in A[T]$ of degree $n$ such that $F(u)=0$;
moreover, since $\bar{\xi}$ is a root of the canonical image of $F$ in $k[T]$, this image is necessarily equal to $f$.
Thus $u$ generates the $A$-algebra $B$, which is consequently a quotient of $A[T]/F A[T]$.
Moreover, $B$ is semilocal and, at each of its maximal ideals $\mathfrak{n}_i$, \sref{IV.18.4.4} gives $F'(u)\notin\mathfrak{n}_i$; hence $F'(u)$ is invertible in $B$, and $F$ is separable.
Finally, if $B$ is an \'etale $A$-algebra, then $B$, being a flat $A$-module of finite presentation \sref{IV.1.4.7}, is free, and $1,u,\ldots,u^{n-1}$ form a basis of the $A$-module $B$ (Bourbaki, \emph{Alg.\ comm.}, ch.~II, \textsection~3, no.~2, prop.~5); in other words,
\oldpage[IV-4]{120}
the $A$-algebra $B$ is isomorphic to $A[T]/F A[T]$; and any other monic polynomial $G\in A[T]$ such that $B\simeq A[T]/G A[T]$ necessarily has degree $n$.
\end{proof}

\begin{theorem}[18.4.6]
\label{IV.18.4.6}
(Chevalley). ---
\begin{enumerate}
  \item[\textup{(i)}] Let $f:X\to Y$ be a morphism locally of finite type, $x$ a point of $X$, $y=f(x)$, and set $A=\sh{O}_{Y,y}$.
For $\sh{O}_{X,x}$ to be a formally \'etale \textup{(respectively,} formally unramified\textup{)} $A$-algebra, it is necessary and sufficient that there exist a monic polynomial $F\in A[T]$ and a maximal ideal $\mathfrak{n}$ of $B=A[T]/F A[T]$ \textup{(respectively,} of a quotient algebra $B$ of $A[T]/F A[T]$\textup{)} such that $\sh{O}_{X,x}$ is $A$-isomorphic to $B_\mathfrak{n}$ and, if $u$ is the image of $T$ in $B$, one has $F'(u)\notin\mathfrak{n}$.
  \item[\textup{(ii)}] Suppose in addition that $f$ is locally of finite presentation.
Then, for $f$ to be \'etale at the point $x$, it is necessary and sufficient that we can moreover take $B=A[T]/F.A[T]$.
\end{enumerate}
\end{theorem}

\begin{proof}
The conditions are sufficient by virtue of \sref{IV.18.4.2}.
To see that they are necessary, we may plainly reduce to the case where $X$ and $Y$ are affine and, in view of \sref{IV.17.4.1.2}, to the case where $f$ is formally unramified and quasi-finite.
It then follows from \sref{IV.8.12.8} that there exist a finite $A$-algebra $C$ and a maximal ideal $\mathfrak{r}$ of $C$ (necessarily lying above the maximal ideal $\mathfrak{m}$ of $A$) such that $\sh{O}_{X,x}$ is $A$-isomorphic to $C_\mathfrak{r}$.
Moreover, by \sref{IV.17.4.1.2}, the residue field $C/\mathfrak{r}=C_\mathfrak{r}/\mathfrak{r}C_\mathfrak{r}$ is a finite separable extension of $k=A/\mathfrak{m}$, and hence is of the form $k[v]$, where $v$ is separable over $k$.
Let $\mathfrak{r}_i$ ($1\leq i\leq h$) be the maximal ideals of the semilocal ring $C$ \emph{other than} $\mathfrak{r}$.
There exists an element $u\in C$ belonging to every $\mathfrak{r}_i$ and whose image in $C/\mathfrak{r}$ is $v$ (Bourbaki, \emph{Alg.\ comm.}, ch.~II, \textsection~1, no.~2, prop.~5).
We shall show that the $A$-subalgebra $B=A[u]$ of $C$ and the ideal $\mathfrak{n}=\mathfrak{r}\cap B$ (necessarily maximal, since $C$ is finite over $B$) answer the question.

For the case where $\sh{O}_{X,x}$ is a formally unramified $A$-algebra, it suffices to prove that $B_\mathfrak{n}$ is \emph{isomorphic to} $C_\mathfrak{r}$; indeed, $B_\mathfrak{n}$ will then be formally unramified over $A$, and the existence of a polynomial $F\in A[T]$ having the asserted properties will follow from \sref{IV.18.4.4}.
Since $f$ is formally unramified, $C/\mathfrak{r}$ is isomorphic to the $A$-algebra $C_\mathfrak{r}/\mathfrak{m}C_\mathfrak{r}$ \sref{IV.17.4.1.2}.
We are thus reduced to proving the following lemma.

\begin{lemma}[18.4.6.1]
\label{IV.18.4.6.1}
Let $A$ be a local ring with maximal ideal $\mathfrak{m}$, let $C$ be a finite $A$-algebra, and let $\mathfrak{r}$ be a maximal ideal of $C$.
Let $u\in C$ belong to every maximal ideal of $C$ distinct from $\mathfrak{r}$, but not to $\mathfrak{r}$, and suppose that $C_\mathfrak{r}/\mathfrak{m}C_\mathfrak{r}$ is a monogenic algebra over $k=A/\mathfrak{m}$, generated by the image of $u$.
Set $B=A[u]$ and $\mathfrak{n}=\mathfrak{r}\cap B$.
Then the canonical homomorphism $B_\mathfrak{n}\to C_\mathfrak{r}$ is an isomorphism.
\end{lemma}

\begin{proof}
Set $R=B\setmin\mathfrak{n}$ and $S=C\setmin\mathfrak{r}$, so that $B_\mathfrak{n}=R^{-1}B$ and $C_\mathfrak{r}=S^{-1}C$; the canonical homomorphism $B_\mathfrak{n}\to C_\mathfrak{r}$ factors as
\[
R^{-1}B\xrightarrow{g}R^{-1}C\xrightarrow{h}S^{-1}C.
\]
It suffices to prove that both homomorphisms are bijective.
First, $h:R^{-1}C\to S^{-1}C=C_\mathfrak{r}$ is bijective.
Indeed, it suffices to see that the images in $R^{-1}C$ of the elements of $S$ are invertible, or equivalently that every maximal ideal $\mathfrak{p}$ of $R^{-1}C$ has inverse image in $C$ disjoint from $S$, and hence equal to $\mathfrak{r}$.
\oldpage[IV-4]{121}
Since $R^{-1}C$ is finite over the local ring $R^{-1}B=B_\mathfrak{n}$, the inverse image of $\mathfrak{p}$ in $R^{-1}B$ is its unique maximal ideal $\mathfrak{n}B_\mathfrak{n}$; if $\mathfrak{q}$ is the inverse image of $\mathfrak{p}$ in $C$, then $\mathfrak{q}\cap B=\mathfrak{n}$ and, since $\mathfrak{n}$ is maximal in $B$ and $C$ is finite over $B$, $\mathfrak{q}$ is one of the maximal ideals of $C$.
Furthermore $u\notin\mathfrak{q}$, since $u\notin\mathfrak{r}$ and $u\in B$; the hypothesis therefore forces $\mathfrak{q}=\mathfrak{r}$.

Since $B\subset C$, the homomorphism $g:R^{-1}B\to R^{-1}C$ is injective \sref[0\textsubscript{\mdseries I}]{0.1.3.2}.
To prove that it is surjective, observe that $R^{-1}C$ is an $R^{-1}B$-module of finite type and that $\mathfrak{m}R^{-1}B$ is contained in the maximal ideal of the local ring $R^{-1}B$.
By Nakayama's lemma, it suffices to prove that
\[
R^{-1}B/\mathfrak{m}R^{-1}B\longrightarrow R^{-1}C/\mathfrak{m}R^{-1}C
\]
is surjective.
But, by the first part of the proof, $R^{-1}C/\mathfrak{m}R^{-1}C$ identifies with $C_\mathfrak{r}/\mathfrak{m}C_\mathfrak{r}$; by hypothesis this $k$-algebra is generated by the image of $u$, and hence it is the image of $R^{-1}B/\mathfrak{m}R^{-1}B$.
\end{proof}

Consider next the case where $f$ is \'etale at $x$.
Replacing $X$ by a neighbourhood of $x$, we may suppose that $X$ is a neighbourhood of $\mathfrak{n}$ in $\operatorname{Spec}(B)$ (1.7.2).
Set $B'=A[T]/F A[T]$ and let $\mathfrak{n}'$ be the inverse image of $\mathfrak{n}$ in $B'$.
Since the image of $F'(T)$ in $B'$ does not belong to $\mathfrak{n}'$, the morphism $\operatorname{Spec}(B')\to\operatorname{Spec}(A)$ is \'etale at $\mathfrak{n}'$ by \sref{IV.18.4.2},~(ii).
As $\operatorname{Spec}(B)\to\operatorname{Spec}(A)$ is \'etale at $\mathfrak{n}$, \sref{IV.17.3.4} shows that $\operatorname{Spec}(B)\to\operatorname{Spec}(B')$ is \'etale at $\mathfrak{n}$; but this morphism is an immersion, so it can be \'etale at a point only if it is a local isomorphism there \sref{IV.17.9.1}.
Thus $B_\mathfrak{n}$ and $B'_{\mathfrak{n}'}$ are isomorphic.

Finally, suppose that $\sh{O}_{X,x}$ is a formally \'etale $A$-algebra.
With the preceding notation, \sref{IV.17.1.5} shows that $B_\mathfrak{n}$ is a formally \'etale $B'_{\mathfrak{n}'}$-algebra; since the homomorphism $B'_{\mathfrak{n}'}\to B_\mathfrak{n}$ is surjective, this is possible only if it is \emph{bijective} \sref[0]{0.19.10.3},~(i).
This completes the proof of \sref{IV.18.4.6}.
\end{proof}

\begin{corollary}[18.4.7]
\label{IV.18.4.7}
Let $f:X\to Y$ be a morphism locally of finite type, $x$ a point of $X$.
For $f$ to be formally unramified at the point $x$, it is necessary and sufficient that there exist an open neighbourhood $U$ of $x$ such that $f|U$ factorises as $U\xrightarrow{j}X'\xrightarrow{h}Y$, where $h$ is an \'etale morphism and $j$ is a closed immersion.
\end{corollary}

\begin{proof}
We can obviously reduce to the case where $Y=\operatorname{Spec}(R)$ is affine and $f$ is of finite type.
If $A=\sh{O}_{Y,y}$, with $y=f(x)$, the condition for $f$ to be formally unramified at the point $x$ is equivalent, by virtue of \sref{IV.17.4.1.2}, to saying that $\sh{O}_{X,x}$ is a formally unramified $A$-algebra.
If it is so, we can apply \sref{IV.18.4.6}, (i);
replacing $Y$ as needed by an affine neighbourhood of $y$, we may suppose (with the notation of \sref{IV.18.4.6}) that the polynomial $F$ is the image in $A[T]$ of a monic polynomial $G\in R[T]$.
We then set $X'=\operatorname{Spec}(R[T]/G.R[T])$;
let $x'$ be the image of the point $\mathfrak{n}$ of $\operatorname{Spec}(B)$ by the morphism corresponding to the composite homomorphism $R[T]/G.R[T]\to A[T]/F.A[T]\to B$.
It follows from \sref{IV.18.4.2} that the canonical morphism $h:X'\to Y$ is \'etale at the point $x'$, hence, restricting $X'$ and $Y$ to open neighbourhoods of $x'$ and $y$ if necessary, we may suppose that $h$ is \'etale.
On the other hand, by virtue of \sref[I]{I.6.5.1},~(ii) and (1.7.2), the local homomorphism $\varphi:\sh{O}_{X',x'}\to\sh{O}_{X,x}$ corresponds to a morphism
\oldpage[IV-4]{122}
$j:U\to X'$ from an open neighbourhood $U$ of $x$ in $X$;
in applying \sref[I]{I.6.5.1},~(i), suppose that $h\circ j=f|U$, hence $j$ is of finite type \sref[I]{I.6.3.4},~(v);
finally, the homomorphism $\varphi$ is surjective by \sref{IV.18.4.6},~(i), hence it follows from \sref[I]{I.6.5.4}, (i), that we can, restricting further $U$ and $X'$, suppose that $j$ is a closed immersion.
This proves the necessity of the condition stated;
the sufficiency is immediate \sref{IV.17.1.3}, (i) and (ii).
\end{proof}

\begin{corollary}[18.4.8]
\label{IV.18.4.8}
Let $f:X\to Y$ be a morphism locally of finite presentation.
For $f$ to be unramified \textup{(respectively,} \'etale\textup{)}, it is necessary and sufficient that there exist a family of flat morphisms $g_\alpha:Y'_\alpha\to Y$ and, for each $\alpha$, an open subset $U_\alpha$ of $X'_\alpha=X\times_Y Y'_\alpha$, such that, if $g'_\alpha:X'_\alpha\to X$ and $f'_\alpha:X'_\alpha\to Y'_\alpha$ are the canonical projections, the $g'_\alpha(U_\alpha)$ cover $X$, and each composite $U_\alpha\to X'_\alpha\xrightarrow{f'_\alpha}Y'_\alpha$ is a closed immersion \textup{(respectively,} an open immersion\textup{)}.
One can moreover take the $g_\alpha$ to be \'etale.
\end{corollary}

\begin{proof}
The necessity of the condition for \'etale morphisms is immediate: take a single $Y'_\alpha=X$, let $g_\alpha=f$, and take for $U_\alpha\subset X\times_Y X$ the diagonal.
When $f$ is unramified, necessity follows from \sref{IV.18.4.7}: choose an open covering $(V_\alpha)$ of $X$ such that, for each $\alpha$, $f|V_\alpha$ factors as $V_\alpha\xrightarrow{j_\alpha}Y'_\alpha\xrightarrow{g_\alpha}Y$, where $j_\alpha$ is a closed immersion and $g_\alpha$ is \'etale.
Then $j_\alpha:V_\alpha\to Y'_\alpha$ factors as $V_\alpha\xrightarrow{s_\alpha}X'_\alpha\xrightarrow{f'_\alpha}Y'_\alpha$, where $s_\alpha$ is a $V_\alpha$-section of $X'_\alpha$; since $g'_\alpha:X'_\alpha\to X$ is \'etale, $s_\alpha$ is an open immersion \sref{IV.17.4.1}, and it suffices to take $U_\alpha=s_\alpha(V_\alpha)$.

The sufficiency of the conditions follows from \sref{IV.17.7.1};
one concludes at each point $x$ of $g'_\alpha(U_\alpha)$, hence in all of $X$ since the $g'_\alpha(U_\alpha)$ cover $X$.
\end{proof}

\begin{proposition}[18.4.9]
\label{IV.18.4.9}
Let $S$ be a prescheme, $f:X\to S$ a morphism, $h:Y\to S$ a morphism locally of finite presentation, $g:X\to Y$ an $S$-morphism, $x$ a point of $X$, and $y=g(x)$.
The following conditions are equivalent:
\begin{enumerate}[label=\emph{\alph*)}]
  \item $h$ is \'etale at the point $y$ and $g$ is flat at the point $x$.
  \item $h$ is unramified at the point $y$ and $f$ is flat at the point $x$.
\end{enumerate}
\end{proposition}

\begin{proof}
Since $f=h\circ g$, condition \emph{a)} implies that $f$ is flat at $x$ \sref{IV.2.1.6}; moreover, an \'etale morphism is unramified, so \emph{a)} implies \emph{b)}.

To prove that \emph{b)} implies \emph{a)}, we may first suppose that $h$ is unramified, after replacing $Y$ by an open neighbourhood of $y$.
Then, after replacing $S$ by an open neighbourhood of $s=h(y)=f(x)$, we may suppose that there exist an \'etale morphism $u:S'\to S$, a point $y'$ of $Y'=Y_{(S')}$ above $y$, and an open neighbourhood $V$ of $y'$ in $Y'$ such that, if $h'=h_{(S')}:Y'\to S'$, the restriction $h'|V$ is a closed immersion \sref{IV.18.4.8}.
If one proves that $h'$ is \'etale at the point $y'$, it will follow that $h$ is \'etale at the point $y$ \sref{IV.17.7.1}, (ii);
moreover, $f'=f_{(S')}:X_{(S')}\to S'$ is flat at every point $x'$ above $x$.
As the projections $v:Y'\to Y$, $w:X'\to X$ are \'etale morphisms (hence flat), if one proves that $g'=g_{(S')}:X_{(S')}\to Y_{(S')}$ is flat at the point $x'$, it will follow from \sref{IV.2.2.11}, (iv) that $g$ is flat at the point $x$.
We may therefore reduce to the case where $h$ is a \emph{closed immersion}
\oldpage[IV-4]{123}
of finite presentation, with $f$ flat at $x$.
Let $\sh{J}$ be the quasi-coherent ideal of finite type in $\sh{O}_S$ defining the closed subprescheme $Y$ of $S$.
The hypothesis that $f$ is flat at $x$ implies that the homomorphism $\sh{O}_{S,s}\to\sh{O}_{X,x}$ is injective \sref[0\textsubscript{\mdseries I}]{0.6.5.1};
\emph{a fortiori} the homomorphism $\sh{O}_{S,s}\to\sh{O}_{Y,y}$ is injective, which means that $\sh{J}_s=0$, since it is the kernel of the preceding homomorphism.
Since $\sh{J}$ is of finite type, there exists an open neighbourhood $U$ of $s$ in $S$ such that $\sh{J}|U=0$ \sref[0\textsubscript{\mdseries I}]{0.5.2.2}.
We can hence suppose that $h$ is an open immersion, and it is clear that $g$ is flat at the point $x$.
\end{proof}

\begin{proposition}[18.4.10]
\label{IV.18.4.10}
Let us denote by $\boldsymbol{P}(f,x)$ a property satisfying the following conditions:
\begin{enumerate}
  \item[$1^\circ$] For every morphism $f:X\to Y$ and every local isomorphism $h:Y\to Z$, $\boldsymbol{P}(f,x)$ is equivalent to $\boldsymbol{P}(h\circ f,x)$ for $x\in X$.
  \item[$2^\circ$] For every morphism $f:X\to Y$, every \'etale morphism $g:Y'\to Y$, every point $x\in X$, if we set $X'=X\times_Y Y'$, $f'=f_{(Y')}:X'\to Y'$ and if $x'\in X'$ is above $x$, the properties $\boldsymbol{P}(f,x)$ and $\boldsymbol{P}(f',x')$ are equivalent \textup{(``invariance by \'etale base change'')}.
\end{enumerate}
Let then $S$ be a prescheme, $f:X\to S$ and $h:Y\to S$ two morphisms, $g:X\to Y$ an $S$-morphism, $x$ a point of $X$, $y=g(x)$, and suppose $h$ \'etale at the point $y$.
Then the properties $\boldsymbol{P}(f,x)$ and $\boldsymbol{P}(g,x)$ are equivalent.
\end{proposition}

\begin{proof}
With the same notation as in the proof of \sref{IV.18.4.9}, and after replacing $S$ (respectively, $Y$) by a neighbourhood of $s=h(y)$ (respectively, of $y$), \sref{IV.18.4.8} gives an \'etale morphism $u:S'\to S$ such that $h'$ is an open immersion.
If $x'\in X'$ lies above $x$, the hypothesis makes $\boldsymbol{P}(f',x')$ (respectively, $\boldsymbol{P}(g',x')$) equivalent to $\boldsymbol{P}(f,x)$ (respectively, $\boldsymbol{P}(g,x)$).
We may therefore reduce to the case where $h$ is an open immersion; condition $1^\circ$ then makes $\boldsymbol{P}(g,x)$ equivalent to $\boldsymbol{P}(h\circ g,x)=\boldsymbol{P}(f,x)$.
\end{proof}

\begin{env}[18.4.11]
\label{IV.18.4.11}
\emph{Examples.} ---
One can take for the property $\boldsymbol{P}(f,x)$ any one of the following, taking into account \sref{IV.17.7.4},~(ii):
\begin{enumerate}
  \item[(i)] $f$ is flat at the point $x$ \sref{IV.2.2.11}, (iv);
  \item[(ii)] $f$ is locally of finite presentation and has codepth $\leq n$ at the point $x$ (\sref{IV.6.8.1} and \sref{IV.6.7.8});
  \item[(iii)] $f$ is locally of finite presentation and Cohen--Macaulay at the point $x$ (\sref{IV.6.8.1} and \sref{IV.6.7.8});
  \item[(iv)] $f$ is locally of finite presentation and possesses the property $(S_n)$ at the point $x$ (\sref{IV.6.8.1} and \sref{IV.6.7.8});
  \item[(v)] $f$ is locally of finite presentation and possesses the property $(R_n)$ at the point $x$ (\sref{IV.6.8.1} and \sref{IV.6.7.8});
  \item[(vi)] $f$ is locally of finite presentation and normal at the point $x$ (\sref{IV.6.8.1} and \sref{IV.6.7.8});
  \item[(vii)] $f$ is locally of finite presentation and reduced at the point $x$ (\sref{IV.6.8.1} and \sref{IV.6.7.8});
  \item[(viii)] $f$ is unramified at the point $x$ \sref{IV.17.7.4};
  \item[(ix)] $f$ is smooth at the point $x$ \sref{IV.17.7.4};
  \item[(x)] $f$ is \'etale at the point $x$ \sref{IV.17.7.4}.
\end{enumerate}
\end{env}

\oldpage[IV-4]{124}
\begin{corollary}[18.4.12]
\label{IV.18.4.12}
\begin{enumerate}
  \item[\textup{(i)}] Let $f:X\to Y$ be a morphism locally of finite type, $x$ a point of $X$.
If $f$ is flat and formally unramified at the point $x$, then, in every factorisation $f|U:U\xrightarrow{j}X'\xrightarrow{h}Y$, where $U$ is an open neighbourhood of $x$, $j$ a closed immersion and $h$ an \'etale morphism, the homomorphism $\sh{O}_{X',j(x)}\to\sh{O}_{X,x}$ corresponding to $j$ is bijective \textup{(in particular $\sh{O}_{X,x}$ is an \emph{essentially \'etale} $\sh{O}_{Y,f(x)}$-algebra \sref{IV.18.6.1})}.
  \item[\textup{(ii)}] For a morphism $f:X\to Y$ to be \'etale, it is necessary and sufficient that it be locally of finite type, formally unramified and flat.
\end{enumerate}
\end{corollary}

\begin{proof}
(i) Since the homomorphism $\sh{O}_{X',j(x)}\to\sh{O}_{X,x}$ is surjective, it suffices to prove that it is injective, and for that it suffices to prove that $\sh{O}_{X,x}$ is an $\sh{O}_{X',j(x)}$-module \emph{faithfully flat}, or merely \emph{flat} (\sref[0\textsubscript{\mdseries I}]{0.6.5.1} and \sref{IV.6.6.2});
in other words, it amounts to showing that $j$ is a flat morphism at the point $x$;
but since $h\circ j=f$ is flat at the point $x$ by hypothesis and $h$ is \'etale, this follows from \sref{IV.18.4.10} and \sref{IV.18.4.11}, (i).

(ii) It remains only to prove the sufficiency of the stated conditions.
For every $x\in X$, there is an open neighbourhood $U$ of $x$ on which $f|U$ admits a factorization having the properties in (i).
Since $f$ is flat and formally unramified at \emph{every} point of $U$, part (i) applies not only at $x$ but at every $z\in U$.
Thus, if $\sh{J}$ is the quasi-coherent ideal of $\sh{O}_{X'}$ defining the closed subprescheme $U$ associated with $j$, then $\sh{J}_{j(z)}=0$ for every $z\in U$; consequently $j$ is an open immersion, and $f$ is \'etale at every point of $U$, hence at every point of $X$.
\end{proof}

\begin{corollary}[18.4.13]
\label{IV.18.4.13}
Let $f:X\to Y$ be a morphism locally of finite type, $x$ a point of $X$, $y=f(x)$.
Suppose that $y$ admits an open neighbourhood which is a reduced prescheme and has only a finite number of irreducible components.
Then, for $f$ to be \'etale at the point $x$, it is necessary and sufficient that $f$ be flat and formally unramified at the point $x$.
\end{corollary}

\begin{proof}
There is only the sufficiency of the condition to prove.
The question being local on $X$ and $Y$, one can suppose \sref{IV.18.4.7} that $f$ factorises as $X\xrightarrow{j}X'\xrightarrow{h}Y$ where $h$ is \'etale and $j$ a closed immersion.
Then $X'$ is reduced \sref{IV.17.5.7} and, replacing $X'$ by an open neighbourhood of $j(x)$, one can suppose that $X'$ has only a finite number of irreducible components.
Then the maximal points of $X'$ are above the maximal points of $Y$ \sref{IV.2.3.4} and since they are finite, their number is finite.
It all amounts to showing, with the notations of the proof of \sref{IV.18.4.12}, that $\sh{J}_{x'}=0$ for all points $x'$ in a neighbourhood of $j(x)$ in $X'$, knowing that $\sh{J}_{j(x)}=0$.
Now, replacing $X'$ by an affine neighbourhood of $j(x)$, we can suppose that all the irreducible components of $X'$ contain $j(x)$;
if $X'=\operatorname{Spec}(A')$, and if $\mathfrak{p}'$ is the prime ideal of $A'$ corresponding to the point $j(x)$, the morphism $\operatorname{Spec}(A'_{\mathfrak{p}'})\to\operatorname{Spec}(A')$ is dominant, hence the corresponding homomorphism $A'\to A'_{\mathfrak{p}'}$ is \emph{injective} since $A'$ is reduced \sref[I]{I.1.2.7}.
If $\sh{J}=\widetilde{\mathfrak{J}}$, where $\mathfrak{J}$ is an ideal of $A'$, then $\mathfrak{J}$ identifies with a submodule of $A'_{\mathfrak{p}'}$; since $\mathfrak{J}_{\mathfrak{p}'}=0$ by hypothesis, it follows that $\mathfrak{J}=0$.
\end{proof}

\begin{corollary}[18.4.14]
\label{IV.18.4.14}
Let $A$ be a local ring with maximal ideal $\mathfrak{m}$, residue field $k$, $B$ a finite $A$-algebra.
\oldpage[IV-4]{125}
\begin{enumerate}
  \item[\textup{(i)}] For $B$ to be a formally unramified $A$-algebra, it is necessary and sufficient that $B\otimes_A k$ be a $k$-algebra \'etale.
  \item[\textup{(ii)}] For $B$ to be a finite \'etale $A$-algebra, it is necessary and sufficient that $B\otimes_A k$ be an \'etale $k$-algebra and that $B$ be a flat $A$-module \textup{(which is equivalent to saying \sref[0\textsubscript{\mdseries I}]{0.6.3.3} that for every maximal ideal $\mathfrak{n}$ of $B$ (necessarily above $\mathfrak{m}$), $B_\mathfrak{n}$ is a flat $A$-module).}
\end{enumerate}
\end{corollary}

\begin{proof}
Only sufficiency needs proof.
It is clear that (ii) follows from (i) and \sref{IV.18.4.12},~(ii), taking into account \sref{IV.17.1.2},~(i).
To prove (i), let us note that if $B\otimes_A k$ is an \'etale $k$-algebra, we have $\Omega^1_{B\otimes_A k/k}=0$ \sref{IV.17.2.1}.
Since $\Omega^1_{B\otimes_A k/k}=\Omega^1_{B/A}\otimes_B(B\otimes_A k)$ \sref[0]{0.20.5.5} and since $B$ is an $A$-algebra of finite type, $\Omega^1_{B/A}$ is a $B$-module of finite type \sref[0]{0.20.4.7}.
But since $B$ is a finite $A$-algebra, $\mathfrak{m}B$ is contained in the radical of $B$ (Bourbaki, \emph{Alg.\ comm.}, ch.~V, \textsection~2, no.~1, prop.~1), hence the lemma of Nakayama proves that $\Omega^1_{B/A}=0$, and hence $B$ is a formally unramified $A$-algebra \sref{IV.17.2.1}.
\end{proof}


\subsection{Henselian local rings.}
\label{subsection:IV.18.5}
\footnote{The notion of a Henselian local ring is due to Azumaya, and that of Henselization to Nagata, to whom the principal results of this theory are also due.}

\begin{env}[18.5.1]
\label{IV.18.5.1}
Let $X$ be a prescheme and $\sh{E}$ a \emph{locally free} $\sh{O}_X$-Module of finite rank; its dual $\check{\sh{E}}=\shHom_{\sh{O}_X}(\sh{E},\sh{O}_X)$ is therefore a locally free $\sh{O}_X$-Module whose rank at every point of $X$ equals that of $\sh{E}$ at that point, and the canonical homomorphism $\sh{E}\to\shHom_{\sh{O}_X}(\check{\sh{E}},\sh{O}_X)=(\check{\sh{E}})^\vee$ is \emph{bijective}.
For every morphism $X'\to X$, set $\sh{E}_{(X')}=\sh{E}\otimes_{\sh{O}_X}\sh{O}_{X'}$, which is a locally free $\sh{O}_{X'}$-Module, and consider the \emph{set} $\Gamma(X',\sh{E}_{(X')})$ of the sections of this $\sh{O}_{X'}$-Module over $X'$.
We shall see that this defines a \emph{contravariant representable functor}
\[
\label{IV.18.5.1.1}
{}^t\! V_{\sh{E}}:X'\leadsto\Gamma(X',\sh{E}_{(X')})
\tag{18.5.1.1}
\]
from the category of $X$-\emph{preschemes} to the category of \emph{sets} \sref[0\textsubscript{\mdseries III}]{0.8.1.8}.

First of all, this is indeed a well-defined functor, for if $f:X''\to X'$ is an $X$-morphism of $X$-preschemes, we have $\sh{E}_{(X'')}=f^*(\sh{E}_{(X')})$, whence \sref[0\textsubscript{\mdseries I}]{0.4.4.3.2} a map $\Gamma(X',\sh{E}_{(X')})\to\Gamma(X',f_*(\sh{E}_{(X'')}))=\Gamma(X'',\sh{E}_{(X'')})$ which completes the definition of the functor ${}^t\! V$.
Let us then show that the $X$-prescheme $\mathbf{V}(\check{\sh{E}})$ \sref[II]{II.1.7.8} \emph{represents} the functor ${}^t\! V$.
It is indeed immediate that $(\sh{E}_{(X')})^\vee=(\check{\sh{E}})_{(X')}$, hence $\mathbf{V}(\check{\sh{E}})\times_X X'=\mathbf{V}((\sh{E}_{(X')})^\vee)$; taking \sref[I]{I.3.3.14} into account, we are reduced to defining a bijection $\Gamma(\mathbf{V}(\check{\sh{E}})/X)\xrightarrow{\sim}\Gamma(X,\sh{E})$ and to verifying that the bijection $\Gamma(\mathbf{V}(\check{\sh{E}}_{(X')})/X')\xrightarrow{\sim}\Gamma(X',\sh{E}_{(X')})$ is functorial in $X'$.
Now, there is a canonical bijection from $\Gamma(\mathbf{V}(\check{\sh{E}})/X)$ to $\operatorname{Hom}_{\sh{O}_X}(\check{\sh{E}},\sh{O}_X)$ \sref[II]{II.1.7.8}, and transposition $u\leadsto{}^t u$ is a canonical bijection from $\operatorname{Hom}_{\sh{O}_X}(\check{\sh{E}},\sh{O}_X)$ to $\operatorname{Hom}_{\sh{O}_X}(\sh{O}_X,\sh{E})=\Gamma(X,\sh{E})$, because of the identification of $(\check{\sh{E}})^\vee$ with $\sh{E}$.
The verification of the functorial property is immediate.
\end{env}

\oldpage[IV-4]{126}
Observe also that, in accordance with the general theory \sref[0\textsubscript{\mdseries III}]{0.8.1.6}, the identity automorphism of $\mathbf{V}(\check{\sh{E}})$ corresponds canonically to a section $c$ of $\sh{E}_{(\mathbf{V}(\check{\sh{E}}))}$ over $\mathbf{V}(\check{\sh{E}})$, that is \sref[II]{II.1.4.1}, to a homomorphism of $\mathbf{S}_{\sh{O}_X}(\check{\sh{E}})$-Modules $u:\mathbf{S}_{\sh{O}_X}(\check{\sh{E}})\to\mathbf{S}_{\sh{O}_X}(\check{\sh{E}})\otimes_{\sh{O}_X}\sh{E}$; for every affine open subset $W$ of $X$, if one sets $\Gamma(W,\sh{O}_X)=A$, if one identifies $\Gamma(W,\mathbf{S}_{\sh{O}_X}(\check{\sh{E}}))$ with a polynomial algebra $C=A[T_1,\ldots,T_n]$, so that the $T_i$ form a basis of $\Gamma(W,\check{\sh{E}})$, and if finally one denotes by $(e_i)$ the basis dual to $(T_i)$ in $\Gamma(W,\sh{E})=\Gamma(W,\check{\sh{E}})^\vee$, one sees immediately that $u$ corresponds to the homomorphism of $C$-modules such that $u(1)=\sum_{i=1}^n T_i\otimes e_i$.

If $X'=\operatorname{Spec}(A)$ is affine and such that $\sh{E}_{(X')}$ is isomorphic to $\sh{O}_{X'}^n$, $\Gamma(X',\sh{E}_{(X')})$ is a free $A$-module of rank $n$; figuratively, one may say that the object $\mathbf{V}(\check{\sh{E}})$ represents ``the set of points of the twisted affine space over $X$ defined by $\sh{E}$''.

\begin{env}[18.5.2]
\label{IV.18.5.2}
Recall \sref[II]{II.1.7.8} that by definition $\mathbf{V}(\check{\sh{E}})=\operatorname{Spec}(\mathbf{S}_{\sh{O}_X}(\check{\sh{E}}))$.
Let $\sh{J}$ be a quasi-coherent ideal of $\mathbf{S}_{\sh{O}_X}(\check{\sh{E}})$, so that $\operatorname{Spec}(\mathbf{S}_{\sh{O}_X}(\check{\sh{E}})/\sh{J})$ is a closed subprescheme of $\mathbf{V}(\check{\sh{E}})$; we shall interpret it as \emph{representing a functor} from the category of $X$-preschemes to that of sets.
A section $u\in\Gamma(X,\sh{E})$ identifies canonically with an $\sh{O}_X$-homomorphism $u:\sh{O}_X\to\sh{E}$, to which correspond by transposition an $\sh{O}_X$-homomorphism ${}^t u:\check{\sh{E}}\to\sh{O}_X$ and hence an $\sh{O}_X$-algebra homomorphism $v:\mathbf{S}_{\sh{O}_X}(\check{\sh{E}})\to\sh{O}_X$.
Let $\mathrm{Al}(X,\sh{E},\sh{J})$ denote the set of $u\in\Gamma(X,\sh{E})$ such that $\sh{J}$ is \emph{contained in the kernel of} $v$;
it follows immediately that
\[
X'\leadsto\mathrm{Al}(X',\sh{E}_{(X')},\sh{J}\otimes_{\sh{O}_X}\sh{O}_{X'})
\]
is represented by $\operatorname{Spec}(\mathbf{S}_{\sh{O}_X}(\check{\sh{E}})/\sh{J})$.
If $X'=\operatorname{Spec}(A)$ is affine, $\sh{E}_{(X')}\simeq\sh{O}_{X'}^n$, and $\sh{J}\otimes_{\sh{O}_X}\sh{O}_{X'}=\widetilde{\mathfrak{J}}$ for an ideal $\mathfrak{J}$ of $A[T_1,\ldots,T_n]$, then this set identifies with the subset of $A^n$ consisting of the points $(t_1,\ldots,t_n)$ such that $P(t_1,\ldots,t_n)=0$ for every $P\in\mathfrak{J}$.
Thus $\operatorname{Spec}(\mathbf{S}_{\sh{O}_X}(\check{\sh{E}})/\sh{J})$ represents ``the algebraic subset of the twisted affine space defined by $\sh{E}$ consisting of the points annihilating $\sh{J}$''.
We also denote this $X$-prescheme by $\mathbf{Al}(\sh{E},\sh{J})$.
If $\sh{J}$ is an ideal of finite type in $\mathbf{S}_{\sh{O}_X}(\check{\sh{E}})$, then $\mathbf{Al}(\sh{E},\sh{J})$ is an $X$-prescheme of finite presentation, since $\mathbf{S}_{\sh{O}_X}(\check{\sh{E}})$ is an $\sh{O}_X$-algebra of finite presentation.
\end{env}

\begin{lemma}[18.5.3]
\label{IV.18.5.3}
Let $S$ be a prescheme, $f:X\to S$ a finite and locally free morphism \sref{IV.18.2.7}.
Consider the contravariant functor from the category of $S$-preschemes to the category of sets
\[
\label{IV.18.5.3.1}
S'\leadsto\operatorname{Of}(X\times_S S')
\tag{18.5.3.1}
\]
where $\operatorname{Of}(X\times_S S')$ denotes the set of subsets that are both open and closed of the underlying space of $X\times_S S'$.
Then this functor is representable by an $S$-prescheme $\mathbf{Of}(X)$, which is affine, \'etale and of finite presentation over $S$.
\end{lemma}

\begin{proof}
\oldpage[IV-4]{127}
By hypothesis $X=\operatorname{Spec}(\sh{B})$, where $\sh{B}=f_*(\sh{O}_X)$ is a finite and locally free $\sh{O}_S$-algebra.
Set $X'=X\times_S S'$, $f'=f_{(S')}$, $\sh{B}'=\sh{B}\otimes_{\sh{O}_S}\sh{O}_{S'}=f'_*(\sh{O}_{X'})$, so that $X'=\operatorname{Spec}(\sh{B}')$;
there is then a canonical bijection, functorial in $S'$, from the set $\operatorname{Of}(X')$ to the set $\operatorname{Id}(\sh{B}')$ of \emph{idempotents} of the ring $\Gamma(S',\sh{B}')=\Gamma(X',\sh{O}_{X'})$.
Indeed, by virtue of the equivalence of the category of affine $S'$-schemes and the opposite category of quasi-coherent $\sh{O}_{S'}$-algebras (\sref[II]{II.1.2.7} and \sref[II]{II.1.3.1}), there is a canonical bijective correspondence between decompositions of $X'$ as a \emph{disjoint union} $X'_1\sqcup X'_2$ of two induced open subpreschemes and decompositions of $\sh{B}'$ as a \emph{direct sum} of two ideals $\sh{B}'_1$, $\sh{B}'_2$;
these latter in turn form a set in canonical bijective correspondence with $\operatorname{Id}(\sh{B}')$.
It suffices therefore to prove the lemma for the functor $S'\leadsto\operatorname{Id}(\sh{B}')$.

For this, we shall show that there exists an ideal $\sh{J}$ of finite type in $\mathbf{S}_{\sh{O}_S}(\check{\sh{B}})$ such that $\operatorname{Id}(\sh{B}')$ is of the form $\mathrm{Al}(S',\sh{B}',\sh{J}\otimes_{\sh{O}_S}\sh{O}_{S'})$ \sref{IV.18.5.2}.
Since $\sh{B}$ is an $\sh{O}_S$-algebra, its inverse image $\sh{B}_{\mathbf{V}(\check{\sh{B}})}$ is an $\sh{O}_{\mathbf{V}(\check{\sh{B}})}$-algebra, and we may form in this algebra the square $c^2$ of the canonical section $c$ \sref{IV.18.5.1};
it corresponds canonically to a homomorphism $u^{(2)}:\mathbf{S}_{\sh{O}_S}(\check{\sh{B}})\to\mathbf{S}_{\sh{O}_S}(\check{\sh{B}})\otimes_{\sh{O}_S}\sh{B}$, and one verifies immediately that for every affine open subset $W$ of $S$, with the notations of \sref{IV.18.5.1}, $u^{(2)}$ corresponds to the $C$-module homomorphism such that $u^{(2)}(1)=\sum_k(\sum_{i,j}c_{ijk}T_iT_j)\otimes e_k$, where $(c_{ijk})$ is the table of multiplication of the algebra $\Gamma(W,\sh{B})$.
Let $\sh{J}$ be the ideal of $\mathbf{S}_{\sh{O}_S}(\check{\sh{B}})$ locally generated by the coefficients of $(u\setmin u^{(2)})(1)$ with respect to a basis of $\sh{B}$; it answers the question.
Indeed, $\Gamma(W,\sh{J})$ is generated by the polynomials $P_k(T_1,\ldots,T_n)=T_k\setmin\sum_{i,j}c_{ijk}T_iT_j$, and the idempotents of $\Gamma(W,\sh{B})$ are precisely the elements $\sum_i t_i e_i$ such that $(t_1,\ldots,t_n)$ annihilates every $P_k$ ($1\leq k\leq n$).
We deduce from \sref{IV.18.5.2} that the $S$-prescheme affine $\mathbf{Of}(X)=\mathbf{Al}(\sh{B},\sh{J})$ does represent the functor \sref{IV.18.5.3.1}.
It remains to show that $\mathbf{Of}(X)$ is \'etale over $S$, or equivalently, that it is \emph{formally \'etale} over $S$.
But if $S'$ is an $S$-prescheme, $S_0'$ a closed sub-prescheme of $S'$ defined by a locally nilpotent Ideal (and having the same underlying space as $S'$), it is clear that $X'=X\times_S S'$ and $X_0'=X\times_S S_0'$ have the same underlying space, hence the canonical map $\operatorname{Of}(X')\to\operatorname{Of}(X_0')$ is \emph{bijective}, which completes the proof \sref{IV.17.1.1}.
\end{proof}

\begin{proposition}[18.5.4]
\label{IV.18.5.4}
Let $S$ be a prescheme, $S_0$ a closed sub-prescheme of $S$;
consider the following properties:
\begin{enumerate}[label=\emph{\alph*)}]
  \item For every \textup{finite} morphism $g:S'\to S$, the canonical map
\[
\label{IV.18.5.4.1}
\operatorname{Of}(S')\to\operatorname{Of}(S'\times_S S_0)
\tag{18.5.4.1}
  \]
is bijective \textup{(cf.\ \sref{IV.18.5.3})}.
  \item[\textup{a$'$)}] For every \textup{finite and locally free} morphism $g:S'\to S$, the map \sref{IV.18.5.4.1} is bijective.
\oldpage[IV-4]{128}
  \item For every \textup{\'etale and separated} morphism $g:S'\to S$, the canonical map
\[
\label{IV.18.5.4.2}
\Gamma(S'/S)\to\Gamma(S'\times_S S_0/S_0)
\tag{18.5.4.2}
\]
is bijective.
\end{enumerate}
Condition \textup{b)} implies \textup{a$'$)};
if moreover $S$ is quasi-compact and quasi-separated, condition \textup{a)} implies \textup{b)}.
\end{proposition}

\begin{proof}
Let us first prove that \emph{b)} implies \emph{a$'$)}.
Suppose therefore \emph{b)} satisfied, and let $g:S'\to S$ be a finite and locally free morphism;
set $S_0'=S'\times_S S_0$, so that $g_0=g_{(S_0)}:S_0'\to S_0$ is finite and locally free.
Then it follows from \sref{IV.18.5.3} that $P=\mathbf{Of}(S')$ is an $S$-prescheme \emph{\'etale and separated};
moreover, by the definition of the functor $\operatorname{Of}$, we see immediately that if we set $P_0=\mathbf{Of}(S_0')$ (for the category of $S_0$-preschemes), we have $P_0=P\times_S S_0$.
This being so, by definition we have the commutative diagram
\[
\xymatrix{
\Gamma(P/S)\ar[r]\ar[d]^{\wr} & \Gamma(P_0/S_0)\ar[d]^{\wr}\\
\operatorname{Of}(S')\ar[r] & \operatorname{Of}(S_0')
}
\]
where the vertical arrows are the canonical bijections.
Since hypothesis \emph{b)}, applied to the morphism $P\to S$, implies that the first row is a bijection, the same holds for the second, which establishes our assertion.

Before proving that \emph{a)} implies \emph{b)} when $S$ is quasi-compact and quasi-separated, we will establish the following lemma.

\begin{lemma}[18.5.4.3]
\label{IV.18.5.4.3}
If $S$ and $S_0$ satisfy condition \textup{a)} of \sref{IV.18.5.4}, then, for every \textup{finite} morphism $g:S'\to S$, $S'$ is the unique neighbourhood of $S_0'=g^{-1}(S_0)=S'\times_S S_0$ in $S'$.
\end{lemma}

\begin{proof}
Indeed, it amounts to saying that if $T'$ is a closed subset of $S'$ such that $T'\cap S_0'=\varnothing$, then $T'=\varnothing$.
Now, if we also denote by $T'$ the closed sub-prescheme of $S'$ having $T'$ as underlying space, the composite morphism $h:T'\to S'\xrightarrow{g}S$ is finite and $h^{-1}(S_0)$ is empty;
condition \emph{a)}, applied to the morphism $h$, implies that $T'$ is necessarily empty.
\end{proof}

This lemma being established, let us first prove that, under hypothesis \emph{a)}, the map \sref{IV.18.5.4.2} is \emph{injective}.
Indeed, if $u'$, $u''$ are two $S$-sections of $S'$, the fact that the morphism $S'\to S$ is unramified implies that the prescheme of coincidences of $u'$ and $u''$ is induced on an open subset $U$ of $S$ \sref{IV.17.4.6}.
If the restrictions to $S_0$ of $u'$ and $u''$ are equal, then the fact that $u'$ and $u''$ are open immersions \sref{IV.17.4.1} implies that $U$ contains $S_0$, and hence is equal to $S$ by virtue of Lemma \sref{IV.18.5.4.3}, applied to the case $S'=S$.

\oldpage[IV-4]{129}
It remains to show that, under hypothesis \emph{a)}, the map \sref{IV.18.5.4.2} is \emph{surjective} (when $S$ is quasi-compact and quasi-separated).
Let therefore $u_0:S_0\to S_0'$ be an $S_0$-section of $S_0'$;
$u_0(S_0)$ being quasi-compact in $S'$, can be covered by a finite number of open affine subsets $V_i$ such that the restriction $g|V_i$ is a morphism of finite type;
if $V$ is the union of the $V_i$, we conclude that $g|V:V\to S$ is a morphism \emph{of finite presentation}, being separated and locally of finite presentation by hypothesis \sref{IV.1.6.1}.
Replacing $S'$ by $V$, one can hence suppose that $g$ is of finite presentation.
Since $S$ is quasi-compact and quasi-separated, and $g$ is quasi-finite and separated \sref{IV.17.6.1}, it follows from the ``Main theorem'' \sref{IV.8.12.6} that $g$ factorises as $S'\xrightarrow{j}S''\xrightarrow{f}S$, where $j$ is an open immersion and $f$ a \emph{finite} morphism.
Set $S_0''=S''\times_S S_0$, $j_0=j_{(S_0)}:S_0'\to S_0''$, which is an open immersion, and $f_0=f_{(S_0)}:S_0''\to S_0$, which is a finite morphism.
Then $u_0$ is also an $S_0$-section of $S_0''$.
Since $g_0:S_0'\to S_0$ is \'etale, $u_0$ is an open immersion of $S_0$ into $S_0'$ \sref{IV.17.4.1}, so $u_0(S_0)$ is open in $S_0'$, and \emph{a fortiori} in $S_0''$;
on the other hand, since $f_0$ is finite and hence separated, $u_0$ is a closed immersion of $S_0$ into $S_0''$ \sref[I]{I.5.4.6}.
Thus $X_0=u_0(S_0)$ is \emph{both open and closed} in $S_0''$.
By virtue of hypothesis \emph{a)}, there exists an open and closed subset $X$ of $S''$ such that $X\cap S_0''=X_0$.
Let us first show that the morphism $f:S''\to S$ is \emph{\'etale at the points of $X$}:
indeed, the set $U$ of points of $X$ where $f|X$ is \'etale is open and contains $X_0\subset S_0'$ by hypothesis.
But the lemma \sref{IV.18.5.4.3} applied to the finite morphism $f|X$ proves that $U=X$.
On the other hand, $S'\cap X$ is open in $X$ and contains $X_0$ by hypothesis, so the same reasoning proves that $S'\cap X=X$, i.e.\ $X\subset S'$.
It remains to show that, for every $s\in S$, the geometric number $n(s)$ of points of $X\cap f^{-1}(s)$ is equal to $1$;
it will then follow that $f|X$ is radicial and surjective, and, since $f|X$ is \'etale, we shall have proved \sref{IV.17.9.1} that $f|X$ is an isomorphism from the open subset $X\subset S'$ onto $S$, whose inverse isomorphism $u$ will be the desired $S$-section extending $u_0$.
But since $f|X$ is \'etale and \emph{finite}, $s\leadsto n(s)$ is continuous on $S$ \sref{IV.18.2.8}, and since $X\cap S_0'=X_0$, we have $n(s)=1$ on $S_0$;
the set of points $s\in S$ such that $n(s)=1$, being open in $S$ and containing $S_0$, is equal to $S$ by \sref{IV.18.5.4.3}.
\end{proof}

\begin{remark}[18.5.4.4]
\label{IV.18.5.4.4}
One can show that the statement \sref{IV.18.5.4} remains valid when, in condition \emph{b)}, one supposes merely the morphism $g$ \'etale (but not necessarily separated) \cite[exp.~XII]{IV-43}.
\end{remark}

\begin{definition}[18.5.5]
\label{IV.18.5.5}
We say that a prescheme $S$ and a closed sub-prescheme $S_0$ of $S$ form a \emph{Henselian couple} if they satisfy condition \textup{a)} of \sref{IV.18.5.4}.
\end{definition}

Taking into account \sref[I]{I.5.1.8}, it amounts to the same to say that $(S,S_0)$ is a Henselian couple or that $(S_{\mathrm{red}},(S_0)_{\mathrm{red}})$ is one.

\begin{proposition}[18.5.6]
\label{IV.18.5.6}
\begin{enumerate}
  \item[\textup{(i)}] If $(S,S_0)$ is a Henselian couple, then, for every finite morphism $f:S'\to S$, if $S_0'$ is the closed sub-prescheme $f^{-1}(S_0)$ of $S'$, the couple $(S',S_0')$ is Henselian.
  \item[\textup{(ii)}] Let $S=\coprod_\alpha S^{(\alpha)}$ be a sum of preschemes, $S_0$ a closed sub-prescheme of $S$, the sum of sub-preschemes $S_0^{(\alpha)}$ of the $S^{(\alpha)}$.
For the couple $(S,S_0)$ to be Henselian, it is necessary and sufficient that each of the couples $(S^{(\alpha)},S_0^{(\alpha)})$ be Henselian.
\end{enumerate}
\end{proposition}

\begin{proof}
Assertion (i) is an immediate consequence of the definition, since for every finite morphism $g:S''\to S'$, $f\circ g:S''\to S$ is a finite morphism.
Similarly, under the conditions of (ii), for a morphism $g:S'\to S$ to be finite, it is necessary and sufficient that each \oldpage[IV-4]{130}of its restrictions $g^{(\alpha)}:S'^{(\alpha)}=g^{-1}(S^{(\alpha)})\to S^{(\alpha)}$ be finite.
If we set $S_0'=g^{-1}(S_0)$ and $S_0'^{(\alpha)}=(g^{(\alpha)})^{-1}(S_0^{(\alpha)})$, the open-and-closed subsets $U$ (resp.\ $U_0$) of $S'$ (resp.\ $S_0'$) correspond bijectively to the families $(U^{(\alpha)})$ (resp.\ $(U_0^{(\alpha)})$), where $U^{(\alpha)}$ (resp.\ $U_0^{(\alpha)}$) is an open-and-closed subset of $S'^{(\alpha)}$ (resp.\ $S_0'^{(\alpha)}$); assertion (ii) follows.
\end{proof}

\begin{remark}[18.5.7]
\label{IV.18.5.7}
Let $S=\operatorname{Spec}(A)$ be an affine scheme, $S_0$ a closed sub-prescheme defined by an ideal $\mathfrak{J}$ of $A$.
Then, if the couple $(S,S_0)$ is Henselian, the ideal $\mathfrak{J}$ is necessarily contained in the \emph{radical} of $A$.
Indeed, if $\mathfrak{m}$ is a maximal ideal of $A$, $\mathfrak{m}$ must belong to $V(\mathfrak{J})=S_0$, by virtue of \sref{IV.18.5.4.3}, whence the conclusion.
In particular, suppose that $S_0$ is \emph{reduced to a point}, then $\mathfrak{J}$ must be the radical of $A$, so $A$ must be a \emph{local} ring, and $S_0$ the unique closed point of $\operatorname{Spec}(A)$.
\end{remark}

\begin{definition}[18.5.8]
\label{IV.18.5.8}
We say that a ring $A$ is Henselian if it is semi-local and if, denoting by $\mathfrak{r}$ the radical of $A$, the couple $(\operatorname{Spec}(A),\operatorname{Spec}(A/\mathfrak{r}))$ is Henselian.
We call a Henselian local scheme a scheme isomorphic to the spectrum of a Henselian local ring.
\end{definition}

\begin{proposition}[18.5.9]
\label{IV.18.5.9}
\begin{enumerate}
  \item[\textup{(i)}] For a semi-local ring $A$ to be Henselian, it is necessary and sufficient that it be a direct product of Henselian local rings.
  \item[\textup{(ii)}] For a local ring $A$ to be Henselian, it is necessary and sufficient that every finite $A$-algebra $B$ be isomorphic to a product of \textup{local} rings.
\end{enumerate}
\end{proposition}

\begin{proof}
(i) Indeed, the definition, applied to $S=\operatorname{Spec}(A)$ and $S_0=\operatorname{Spec}(A/\mathfrak{r})$, shows that $S_0$ is a finite discrete set, closed in $S$, and $S$ is the union of a finite number of open and closed subsets $S_i$ ($1\leq i\leq n$), pairwise disjoint, each containing exactly one of the maximal ideals $\mathfrak{m}_i$ of $A$;
the conclusion follows from \sref{IV.18.5.6}, (ii) and the remark \sref{IV.18.5.7}.

(ii) Every finite morphism $S'\to\operatorname{Spec}(A)$ is of the form $\operatorname{Spec}(B)\to\operatorname{Spec}(A)$, where $B$ is a finite $A$-algebra.
If $k$ is the residue field of $A$, $\operatorname{Spec}(B\otimes_A k)$ is a finite Artinian spectrum, hence finite and discrete.
To say that the couple $(\operatorname{Spec}(A),\operatorname{Spec}(k))$ is Henselian means that $B$ is a direct product of rings $A_i$ such that $\operatorname{Spec}(A_i\otimes_A k)$ is reduced to a point, i.e.\ that each $A_i$ (which is a finite $A$-algebra) must have only a single maximal ideal (Bourbaki, \emph{Alg.\ comm.}, ch.~V, \textsection~2, no.~1, prop.~1).
\end{proof}

The study of Henselian rings is thus essentially reduced to that of Henselian \emph{local} rings.

\begin{proposition}[18.5.10]
\label{IV.18.5.10}
If $A$ is a Henselian ring, every finite $A$-algebra $B$ is a Henselian ring \textup{(hence} a direct product of Henselian local rings \textup{\sref{IV.18.5.9})}.
\end{proposition}

\begin{proof}
Indeed, if $\mathfrak{r}'$ is the radical of $B$, the inverse image of $\mathfrak{r}'$ in $A$ is the radical $\mathfrak{r}$ of $A$, and every prime ideal of $B$ lying over a maximal ideal of $A$ is a maximal ideal of $B$;
hence the set $V(\mathfrak{r}')$ in $\operatorname{Spec}(B)$ is the inverse image of the set $V(\mathfrak{r})$ in $\operatorname{Spec}(A)$.
The proposition is then a consequence of \sref{IV.18.5.6}, (i).
\end{proof}

\begin{theorem}[18.5.11]
\label{IV.18.5.11}
Let $A$ be a local ring, $\mathfrak{m}$ its maximal ideal.
The following conditions are equivalent:
\begin{enumerate}[label=\emph{\alph*)}]
  \item $A$ is Henselian, in other words, every finite $A$-algebra $B$ is isomorphic to a product of local rings.
\oldpage[IV-4]{131}
  \item[\textup{a$'$)}] Condition \textup{a)} is satisfied for all $A$-algebras $B$ of the form $A[T]/F.A[T]$, where $F\in A[T]$ is a monic polynomial.
  \item Let $S=\operatorname{Spec}(A)$, $S_0=\operatorname{Spec}(\boldsymbol{k})$.
For every \'etale morphism $g:S'\to S$, if we set $S_0'=S'\otimes_A k$, every $S_0$-section $u_0$ of $S_0'$ is the restriction of an $S$-section $u$ of $S'$.
  \item For every morphism locally of finite type $f:X\to S$, separated, and every point $x\in X$ such that $f(x)$ is equal to the closed point $s$ of $S$ and that $f$ is quasi-finite at the point $x$ \textup{(\textbf{Err}$_{\mathrm{III}}$, 20)}, $X$ is the sum of two preschemes $X'$, $X''$ such that $X'=\operatorname{Spec}(\sh{O}_{X,x})$ and that $f|X':X'\to S$ is a \textup{finite} morphism.
  \item[\textup{c$'$)}] For every morphism locally of finite type $f:X\to S$, and every point $x\in X$ such that $f$ be quasi-finite at the point $x$ and $f(x)$ equal to the closed point $s$ of $S$, $\sh{O}_{X,x}$ is a finite algebra over $\sh{O}_{S,s}=A$.
  \item[\textup{c$''$)}] For every morphism locally of finite presentation $f:X\to S$, and every point $x\in X$ such that $f$ be quasi-finite at the point $x$ and $f(x)$ equal to the closed point $s$ of $S$, $\sh{O}_{X,x}$ is a finite algebra and of finite presentation over $\sh{O}_{S,s}=A$.
\end{enumerate}
\end{theorem}

\begin{proof}
Let us first note that \emph{c$'$)} (resp.\ $c'')$) is equivalent to the same condition where moreover $f$ is supposed \emph{separated}, the question being local on $X$.
Similarly, condition \emph{b)} is equivalent to the same condition where moreover $g$ is supposed \emph{separated}:
indeed, it suffices to apply the latter to the restriction of $g$ to an open neighbourhood affine of the point $u_0(S_0)$ in $S'$.
Let us henceforth restrict to the case where, in \emph{b)}, \emph{c$'$)} and \emph{c$''$)}, the morphisms are separated.

The fact that \emph{a)} implies \emph{b)} and that \emph{b)} implies \emph{a$'$)} is a particular case of \sref{IV.18.5.4}.
Let us moreover show that \emph{a$'$)} implies \emph{a)}.
It is a matter of proving that if $e_0$ is an idempotent of $C=B\otimes_A k$, there is an idempotent $e\in B$ whose image is $e_0$ in $C$.
If $b$ is an element of $B$ whose image in $C$ is $e_0$, the sub-$A$-algebra $A[b]=B'$ of $B$ is finite, and the canonical image $C'$ of $B'\otimes_A k$ in $C$ contains $e_0$.
Now, $B'\otimes_A k$ is a finite $k$-algebra, hence a direct product of finite local $k$-algebras, and consequently $e_0$ is the image in $C'$ of an idempotent $e_0'$ of $B'\otimes_A k$.
We are thus reduced to the case where $B$ is monogenic, and hence by $a')$, a quotient of an algebra of the form $A[T]/F.A[T]$;
now, by \emph{a$'$)}, $A[T]/F.A[T]$ is a direct product of local rings, hence so is each of its quotient algebras, which proves the existence of the idempotent $e$.

It is immediate that \emph{c)} implies \emph{a)}, as one sees by induction on the number of maximal ideals of the semi-local ring $B$.
To see that \emph{a)} implies \emph{c)}, we may, by \sref{IV.13.1.4}, restrict to the case where $f$ is affine and quasi-finite.
Then, by the ``Main theorem'' \sref{IV.8.12.8}, $f$ can be written as the composite
\[
X\xrightarrow{j}Y\xrightarrow{g}S,
\]
where $g$ is finite and $j$ is an open immersion.
Since $Y=\operatorname{Spec}(B)$, where $B$ is a finite $A$-algebra, condition \emph{a)} implies that $B$ is a direct product of local rings, which are evidently finite $A$-algebras, and $\sh{O}_{X,x}$ identifies with one of these local rings because $f(x)=s$;
moreover, every open subset of $Y$ containing $x$ necessarily contains $\operatorname{Spec}(\sh{O}_{X,x})$.

It is trivial that \emph{c)} implies \emph{c$'$)}, by virtue of the remark at the beginning.
Let us prove that \emph{c$'$)} implies \emph{c$''$)}.
Suppose \emph{c$'$)} satisfied and let us prove that, under the hypotheses \oldpage[IV-4]{132}of \emph{c$''$)}, the subset $Z=\operatorname{Spec}(\sh{O}_{X,x})$ identifies with an \emph{open and closed} subset of $X$, which will establish \emph{c$''$)} \sref[I]{I.2.4.2}.
First, the composite morphism $Z\xrightarrow{j}X\xrightarrow{f}S$, where $j$ is the canonical morphism \sref[I]{I.2.4.1}, is finite and of finite presentation by hypothesis;
and since $f$ is separated and locally of finite presentation, $j$ is also finite and of finite presentation (\sref[II]{II.6.1.5} and (1.4.3)), hence $j$ is a closed morphism \sref[II]{II.6.1.10}, which proves that $Z$ is closed in $X$.
It then follows from \sref[I]{I.2.4.2} and \sref[I]{I.4.2.2} that $j$ is a closed immersion.
If $\sh{J}$ is the ideal of $\sh{O}_X$ defining $Z$, then by hypothesis $\sh{J}_x=0$, and hence $\sh{J}_z=0$ at every point $z$ of some open neighbourhood $V$ of $x$ in $X$, since $\sh{J}$ is a quasi-coherent $\sh{O}_X$-Module of finite type (\sref{IV.1.4.7} and \sref[0\textsubscript{\mdseries I}]{0.5.2.2}).
Such a neighbourhood contains $Z$, so $Z$ is open in $X$, since $\sh{J}|V=0$.

Finally, \emph{c$''$)} implies \emph{a$'$)}:
indeed, if $B=A[T]/F.A[T]$, the morphism $X=\operatorname{Spec}(B)\to\operatorname{Spec}(A)=Y$ is finite and of finite presentation, and the preceding argument shows that $B$ is a direct product of finite $A$-algebras $\sh{O}_{X,x_i}$, where the $x_i$ are the points of the fibre of the closed point of $Y$.
\end{proof}

\begin{corollary}[18.5.12]
\label{IV.18.5.12}
Let $A$ be a semi-local ring, $\mathfrak{r}$ its radical;
set $S=\operatorname{Spec}(A)$, $S_0=\operatorname{Spec}(A/\mathfrak{r})$.
For $A$ to be Henselian, it is necessary and sufficient that, for every \textup{finite} morphism $f:X\to S$ and every \textup{\'etale and separated} morphism $g:Y\to S$, if we set $X_0=X\times_S S_0$ \textup{and} $Y_0=Y\times_S S_0$, the canonical map
\[
\operatorname{Hom}_S(X,Y)\to\operatorname{Hom}_{S_0}(X_0,Y_0)
\]
is bijective.
\end{corollary}

\begin{proof}
The sufficiency of the condition follows from the equivalence of \emph{a)} and \emph{b)} in \sref{IV.18.5.11}, by applying this condition to the case where $f=1_S$.
To see that the condition is necessary, note that if $A$ is Henselian, the couple $(X,X_0)$ is Henselian by \sref{IV.18.5.6}, (i);
moreover $\operatorname{Hom}_S(X,Y)=\Gamma(X\times_S Y/X)$ and $\operatorname{Hom}_{S_0}(X_0,Y_0)=\Gamma(X_0\times_{S_0}Y_0/X_0)$;
since $X\times_S Y$ is \'etale and separated over $X$, the conclusion follows from the fact that \emph{a)} implies \emph{b)} in \sref{IV.18.5.4}.
\end{proof}

\begin{remark}[18.5.13]
\label{IV.18.5.13}
The equivalent conditions \emph{a)}, \emph{a$'$)}, \emph{b)} of Theorem~\sref{IV.18.5.11} are also equivalent to the following (``Hensel's lemma''):
\begin{enumerate}
  \item[\textup{a$''$)}] \itshape For every monic polynomial $F\in A[T]$, with canonical image $F_0\in k[T]$, and every decomposition $F_0=G_0 H_0$ of $F_0$ into a product of two monic polynomials that are \textup{coprime} $G_0$, $H_0$ in $k[T]$, there exists a unique couple $(G,H)$ of monic polynomials of $A[T]$ possessing the following properties:
$G_0$ and $H_0$ are the canonical images respectively of $G$ and $H$, we have $F=GH$, and the ideal of $A[T]$ generated by $G$ and $H$ is equal to $A[T]$.
\end{enumerate}
\end{remark}

\begin{lemma}[18.5.13.1]
\label{IV.18.5.13.1}
Let $A$ be a local ring with residue field $k$, $F\in A[T]$ a monic polynomial, $B$ the $A$-algebra $A[T]/F.A[T]$.
There exists a canonical bijective correspondence between the decompositions of $B$ into the direct product of two quotient $A$-algebras $B'$, $B''$ and the decompositions $F=GH$ of $F$ into a product of two monic polynomials $G$, $H$ of $A[T]$, such that the ideal generated by $G$ and $H$ is equal to $A[T]$;
the quotient algebras $B'$, $B''$ corresponding to such a couple of polynomials $G$, $H$ are respectively $A[T]/H.A[T]$ and $A[T]/G.A[T]$.
\end{lemma}

\begin{proof}
\oldpage[IV-4]{133}
If $F=GH$ and $G$ and $H$ generate the ideal $A[T]$, there are two polynomials $P$, $Q$ of $A[T]$ such that $1=PG+QH$.
We deduce that the intersection of the principal ideals $\mathfrak{a}=G.A[T]$ and $\mathfrak{b}=H.A[T]$ is equal to $\mathfrak{c}=F.A[T]$:
indeed, if $R\in G.A[T]\cap H.A[T]$, we can write $R=PRG+QRH$;
since $RG$ (resp.\ $RH$) is a multiple of $F$ since $R$ is a multiple of $H$ (resp.\ $G$), therefore $R\in F.A[T]$.
Since $A[T]=\mathfrak{a}+\mathfrak{b}$, $A[T]/\mathfrak{c}$ is the direct sum of ideals $\mathfrak{a}/(\mathfrak{a}\cap\mathfrak{b})$ and $\mathfrak{b}/(\mathfrak{a}\cap\mathfrak{b})$, canonically isomorphic respectively to $A[T]/\mathfrak{b}$ and $A[T]/\mathfrak{a}$.

Conversely, suppose given a decomposition of $B$ into a direct product of two $A$-algebras $B'$, $B''$, which identify canonically with two ideals $e'B$, $e''B$ of $B$, corresponding to a decomposition $1=e'+e''$ into orthogonal idempotents $e'$, $e''$ of $B$.
Set further $B_0=B\otimes_A k=k[T]/F_0.k[T]$, where $F_0$ is the canonical image of $F$ in $k[T]$, of the same degree $n$ as $F$;
if $e_0'$, $e_0''$ are the canonical images of $e'$, $e''$ in $B_0$, these are two orthogonal idempotents such that $1=e_0'+e_0''$, and $B_0$ is the direct product of $B_0'=e_0'B_0$ and $B_0''=e_0''B_0$.
Let $t$ and $t_0$ be the canonical images of $T$ in $B$ and $B_0$;
$B_0'$ (resp.\ $B_0''$) is a $k$-algebra generated by $t_0'=e_0't_0$ (resp.\ $t_0''=e_0''t_0$), and it admits a basis of the form $\{e_0',t_0',t_0'^2,\ldots,t_0'^{s-1}\}$ (resp.\ $\{e_0'',t_0'',t_0''^2,\ldots,t_0''^{r-1}\}$) with $r+s=n$.
On the other hand, $B$ being a free $A$-module (with basis $\{1,t,\ldots,t^{n-1}\}$), $B'$ and $B''$ are projective $A$-modules, hence free since $A$ is a local ring (Bourbaki, \emph{Alg.\ comm.}, ch.~II, \textsection~5, no.~3, cor.\ of prop.~5);
it follows from the above and from Bourbaki, \emph{Alg.\ comm.}, ch.~II, \textsection~3, no.~3, prop.~5, that if we set $t'=e't$, $t''=e''t$, $\{e',t',\ldots,t'^{s-1}\}$ (resp.\ $\{e'',t'',\ldots,t''^{r-1}\}$) is a basis of the $A$-module $B'$ (resp.\ $B''$).
There is hence a monic polynomial $H$ (resp.\ $G$) of degree $s$ (resp.\ $r$) in $A[T]$ such that $e'H(t')=0$ and $e''G(t'')=0$;
as $t^h=t'^h+t''^h$ for every integer $h\geq 1$, with $t'^h=e't'^h$ and $t''^h=e''t''^h$, we also have $G(t)=e'G(t')$ and $H(t)=e''H(t'')$, whence $G(t)H(t)=0$;
we conclude that the polynomial $G(T)H(T)$ is divisible by $F(T)$;
but since these two monic polynomials have the same degree, we have $GH=F$.
Moreover, $B'$ (resp.\ $B''$) is isomorphic to $A[T]/H.A[T]$ (resp.\ $A[T]/G.A[T]$).
Finally, there are two polynomials $R$, $S$ in $A[T]$ such that $e'=R(t)$ and $e''=S(t)$;
as $R(t)=e'R(t')+e''R(t'')$, we necessarily have $e''R(t'')=0$ and similarly $e'S(t')=0$, so that, by definition of $G$ and $H$, $R=PG$ and $S=QH$, where $P$, $Q$ belong to $A[T]$;
the relation $1=R(t)+S(t)$ in $B$ gives $PG+QH=1+LF$ for some polynomial $L\in A[T]$, and since $F=GH$, this proves that the ideal generated by $G$ and $H$ is equal to $A[T]$, and completes the proof of the lemma.
\end{proof}

This lemma being established, it suffices to apply it, to the ring $A$ on the one hand, to the field $k$ on the other, to see immediately that conditions \emph{a$'$)} and \emph{a$''$)} are equivalent.

\begin{proposition}[18.5.14]
\label{IV.18.5.14}
Every semi-local ring $A$, separated and complete for the $\mathfrak{r}$-preadic topology \textup{(}where $\mathfrak{r}$ is the radical of $A$\textup{)} is Henselian.
\end{proposition}

\begin{proof}
Indeed, $A$ is a direct product of separated and complete local rings (Bourbaki, \emph{Alg.\ comm.}, ch.~III, \textsection~2, no.~13, cor.\ of prop.~19), so we are reduced to the case where $A$ is a local ring.
Let us verify criterion \emph{a$'$)} of \sref{IV.18.5.11}.
Since $B$ is a free $A$-module of finite type, it is evidently separated and complete for the $\mathfrak{r}$-preadic topology, \oldpage[IV-4]{134}which is also the $\mathfrak{s}$-preadic topology, where $\mathfrak{s}$ is the radical of the semi-local ring $B$, because $B/\mathfrak{r}B$ is an Artinian ring with radical $\mathfrak{s}/\mathfrak{r}B$.
It is then known (Bourbaki, \emph{Alg.\ comm.}, ch.~III, \textsection~2, no.~13, cor.\ of prop.~19) that $B$ is a direct product of local rings.
\end{proof}

\begin{proposition}[18.5.15]
\label{IV.18.5.15}
Let $A$ be a Henselian ring, $\mathfrak{r}$ its radical.
Then the functor $B\leadsto B/\mathfrak{r}B$ is an equivalence of the category of finite \'etale $A$-algebras with the category of finite \'etale $(A/\mathfrak{r})$-algebras.
\end{proposition}

\begin{proof}
The fact that the functor is fully faithful is a particular case of \sref{IV.18.5.12}.
To show that this functor is an equivalence of categories, we may reduce to the case where $A$ is local;
it suffices then to apply \sref{IV.18.1.1} (for \'etale morphisms), with $S=\operatorname{Spec}(A)$ and $S_0=\operatorname{Spec}(A/\mathfrak{r})$, reduced to a single point.
\end{proof}

\begin{remarks}[18.5.16]
\label{IV.18.5.16}
\begin{enumerate}
  \item[(i)] We do not know if Proposition~\sref{IV.18.5.15} generalises to a Henselian couple $(S,S_0)$, even when $S$ is affine and Noetherian.
  \item[(ii)] Let $A$ be a \emph{Noetherian} ring, $\mathfrak{J}$ an ideal of $A$ such that $A$ is \emph{separated and complete} for the $\mathfrak{J}$-preadic topology.
Then the couple $(\operatorname{Spec}(A),\operatorname{Spec}(A/\mathfrak{J}))$ is Henselian:
indeed, for every finite $A$-algebra $B$, $B$ is separated and complete for the $\mathfrak{J}$-preadic topology \sref[0\textsubscript{\mdseries I}]{0.7.3.6}.
Replacing $B$ by $A$ and $\mathfrak{J}B$ by $\mathfrak{J}$, everything amounts to seeing that the map that to every idempotent of $A$ associates its class mod $\mathfrak{J}$, is \emph{bijective}.
Now $A=\varprojlim A/\mathfrak{J}^n$.
Write $\operatorname{Idem}(A)$ for the set of idempotents of $A$, and, for every ring homomorphism $\varphi:A\to B$, let $\operatorname{Idem}(\varphi)$ denote the map $\operatorname{Idem}(A)\to\operatorname{Idem}(B)$ induced by $\varphi$.
It follows from the definition of the projective limit that
\[
\operatorname{Idem}(A)=\varprojlim\operatorname{Idem}(A/\mathfrak{J}^n),
\]
for the transition maps $\psi_{nm}:\operatorname{Idem}(A/\mathfrak{J}^m)\to\operatorname{Idem}(A/\mathfrak{J}^n)$ induced by the canonical maps $A/\mathfrak{J}^m\to A/\mathfrak{J}^n$.
But since $\operatorname{Spec}(A/\mathfrak{J}^n)\to\operatorname{Spec}(A/\mathfrak{J}^m)$ is a homeomorphism, the $\psi_{nm}$ are bijections (as seen in the proof of \sref{IV.18.5.3});
this proves our assertion.
\end{enumerate}
\end{remarks}

\begin{theorem}[18.5.17]
\label{IV.18.5.17}
Let $A$ be a Henselian local ring, $S=\operatorname{Spec}(A)$, $s$ the closed point of $S$.
For every smooth morphism $f:X\to S$, the canonical map
\[
\Gamma(X/S)\to\Gamma(X_s/\boldsymbol{k}(s))
\]
\textup{(}where $X_s=f^{-1}(s)$\textup{)} is surjective.
\end{theorem}

\begin{proof}
The datum of a $\boldsymbol{k}(s)$-section of $X_s$ is equivalent to that of a point $x\in X$ above $s$ rational over $\boldsymbol{k}(s)$, and it is a matter of proving that there exists an $S$-section $u:S\to X$ such that $u(s)=x$.
Taking into account \sref{IV.17.16.3},~(i), we can suppose that $f$ is \emph{\'etale}.
Then the conclusion follows from the criterion \sref{IV.18.5.11}, \emph{b)}.
(The reader will note that by virtue of this criterion, the validity of \sref{IV.18.5.17} for a given local ring $A$ is necessary and sufficient for $A$ to be Henselian.)
\end{proof}

\begin{remark}[18.5.18]
\label{IV.18.5.18}
The conditions of \sref{IV.18.5.11} are also equivalent to the following:
\begin{enumerate}
  \item[\textup{d)}] \itshape For every morphism locally of finite type $f:X\to S$ and every point $x\in X$ such that $f(x)$ be equal to the closed point $s$ of $S$ and that $\sh{O}_{X,x}/\mathfrak{m}_s\sh{O}_{X,x}$ be a field canonically isomorphic to $\boldsymbol{k}(s)=k$, there exists an open neighbourhood $U$ of $x$ in $X$ such that $f|U$ is a closed immersion.
\end{enumerate}
\end{remark}

\begin{proof}
\oldpage[IV-4]{135}
It is immediate that \emph{d)} implies condition \emph{b)} of \sref{IV.18.5.11} since a closed immersion \'etale is an open immersion \sref{IV.18.9.1}.
Conversely, suppose the conditions of \sref{IV.18.5.11} satisfied, and let us prove \emph{d)}.
The hypothesis of \emph{d)} implies that $f$ is quasi-finite at the point $x$ (\textbf{Err}$_{\mathrm{III}}$, 20), by virtue of \sref{IV.13.1.4};
hence, by condition \emph{c)} of \sref{IV.18.5.11}, $\sh{O}_{X,x}$ is a finite $A$-algebra and $\operatorname{Spec}(\sh{O}_{X,x})$ is an open neighbourhood $U$ of $x$ in $X$.
Moreover, since $\sh{O}_{X,x}/\mathfrak{m}_s\sh{O}_{X,x}$ is isomorphic to $k=A/\mathfrak{m}_s$, the Nakayama lemma proves that the homomorphism $A\to\sh{O}_{X,x}$ is surjective, hence $f|U$ is a closed immersion.
\end{proof}

\begin{proposition}[18.5.19]
\label{IV.18.5.19}
Let $A$ be a Noetherian Henselian local ring, $S=\operatorname{Spec}(A)$, $s$ the closed point of $S$, and $f:X\to S$ a proper morphism;
set $X_s=f^{-1}(s)$.
Then the map $Y\leadsto Y\cap X_s$ is a bijection from the set of connected components of $X$ to the set of connected components of $X_s$.
\end{proposition}

\begin{proof}
By considering a closed sub-prescheme of $X$ having for underlying space a connected component of $X$, we are reduced to showing that if $X$ is \emph{connected and non-empty}, then $X_s$ is \emph{connected and non-empty}.
The fact that $X_s$ is non-empty follows from \sref[II]{II.7.2.1};
to show that $X_s$ is connected, we argue by contradiction, considering the Stein factorisation $f:X\xrightarrow{f'}S'\xrightarrow{g}S$ of the proper morphism $f$ \sref[III]{III.4.3.3};
by hypothesis the finite discrete set $g^{-1}(s)$ would contain at least two points.
Since $A$ is Henselian and $g$ separated and of finite type, it would then follow from \sref{IV.18.5.11}, \emph{c)} that $S'$ would be the sum of two non-empty preschemes $S_1'$, $S_2'$ (since the intersection of any one of them with $g^{-1}(s)$ is reduced to a single point).
Since $f'$ is surjective \sref[III]{III.4.3.1}, we would conclude that $X$ is a sum of two non-empty preschemes, contradicting the hypothesis.
\end{proof}
\subsection{Henselization.}
\label{subsection:IV.18.6}

\begin{env}[18.6.1]
\label{IV.18.6.1}
Given a local ring $A$, we shall say that a local $A$-algebra $B$ is \emph{essentially \'etale} if there exist an \'etale $A$-algebra $C$ and a prime ideal $\mathfrak{n}$ of $C$ such that $B$ is $A$-isomorphic to $C_{\mathfrak{n}}$ and the composite homomorphism $A\to C\to C_{\mathfrak{n}}$ is local (in other words, $\mathfrak{n}$ lies above the maximal ideal $\mathfrak{m}$ of $A$).
This implies \sref{IV.17.6.1} that $\mathfrak{n}$ is the only prime ideal of $B$ lying above $\mathfrak{m}$, since $B$ is a local ring;
if $B$, $B'$ are two local essentially \'etale $A$-algebras, every $A$-homomorphism $B\to B'$ is therefore local.

If $A'$ is a local essentially \'etale $A$-algebra and $A''$ a local essentially \'etale $A'$-algebra, then $A''$ is a local essentially \'etale $A$-algebra.
Indeed, by hypothesis $A'=B_{\mathfrak n}$ where $B$ is an \'etale $A$-algebra and $\mathfrak{n}$ a prime ideal of $B$ above the maximal ideal of $A$, and $A''=B'_{\mathfrak{n}'}$, where $B'$ is an \'etale $A'$-algebra and $\mathfrak{n}'$ a prime ideal of $B'$ above $\mathfrak{n}B_{\mathfrak{n}}$.
If we set $S=B\setmin\mathfrak{n}$, so that $A'=S^{-1}B$, then $B'$ is of the form $S^{-1}C$, where $C$ is a $B$-algebra of finite presentation.
Consequently $C$ is an $A$-algebra of finite presentation and $A''$ is of the form $C_{\mathfrak{r}}$, where $\mathfrak{r}$ is a prime ideal of $C$ above the maximal ideal of $A$.
Since $A'$ is formally \'etale over $A$ and $A''$ is formally \'etale over $A'$, $A''$ is formally \oldpage[IV-4]{136}\'etale over $A$ (for the discrete topologies \sref{IV.17.1.3});
the morphism $\operatorname{Spec}(C)\to\operatorname{Spec}(A)$ is therefore \'etale at the point $\mathfrak{r}$ \sref{IV.17.6.1}, and hence there exists an element $g\in C$ such that $C_g$ is an \'etale $A$-algebra, which proves that $A''$ is an essentially \'etale $A$-algebra.

Given a local ring $A$, there exists a \emph{set} $\mathfrak{E}$ of local essentially \'etale $A$-algebras such that every local essentially \'etale $A$-algebra is $A$-isomorphic to an algebra belonging to $\mathfrak{E}$.
It suffices indeed to observe that there exists a set $\mathfrak{F}$ of $A$-algebras of finite type such that every $A$-algebra of finite type is isomorphic to an algebra belonging to $\mathfrak{F}$;
one may take $\mathfrak{F}$ to be the set of quotients of the polynomial algebras $A[T_1,\ldots,T_n]$ ($n\in\mathbf{N}$).

In what follows, if $A$ and $B$ are two local rings, we denote by $\Hom.\operatorname{loc}(A,B)$ the set of \emph{local} homomorphisms from $A$ to $B$.
\end{env}

\begin{lemma}[18.6.2]
\label{IV.18.6.2}
Let $A$, $A'$ be two local rings, $k$, $k'$ their residue fields,
$\vphi:A\to A'$ a local homomorphism making $A'$ an $A$-algebra essentially \'{e}tale \sref{IV.18.6.1}
and such that the corresponding homomorphism $k\to k'$ is bijective.
Then, for every local \textup{Henselian} ring $B$, the canonical map
\[
  \Hom(\vphi,\,1_{B}):\Hom.\operatorname{loc}(A',B)\to\Hom.\operatorname{loc}(A,B)
\]
is bijective.
\end{lemma}

\begin{proof}
Let $S=\Spec(A)$, $Y=\Spec(B)$, and let $s$ and $y$ be the closed points of $S$ and $Y$ respectively.
By hypothesis, $A'$ is isomorphic to a local ring $\sh{O}_{X,x}$ of an $S$-scheme $X$ \'{e}tale over $S$, at a point $x\in X$ above $s$.
Suppose given a local homomorphism $\psi:A\to B$, making $Y$ an $S$-scheme;
it amounts to showing that there exists \emph{one and only one} $S$-morphism $f:Y\to X$ such that $f(y)=x$.
Let $X'=X\times_{S}Y$, and note that since $\boldsymbol{k}(x)=\boldsymbol{k}(s)$, there exists a unique point $x'\in X'$ above $x$ and $y$, and that $\boldsymbol{k}(x')=\boldsymbol{k}(y)$.
It remains to show that there exists a unique $Y$-section $f'$ of $X'$ such that $f'(y)=x'$.
Now, the morphism $g:X'\to Y$ is \'{e}tale and separated, and the fibre $X'_0=g^{-1}(y)$ has at the point $x'$ the local ring $\boldsymbol{k}(x')=\boldsymbol{k}(y)$.
If we set $Y_0=\Spec(\boldsymbol{k}(y))$, there exists a unique $Y_0$-section $f'_0$ of $X'_0$ such that $f'_0(y)=x'$, and the conclusion follows from the fact that $B$ is assumed Henselian and from \sref{IV.18.5.11},~\emph{b)}.
\end{proof}

We say that a local $A$-algebra $A'$ satisfying the conditions of \sref{IV.18.6.2} is \emph{strictly essentially \'{e}tale}.
Note that the criterion \sref{IV.18.5.11},~\emph{b)} means that, \emph{for $A$ to be Henselian, it is necessary and sufficient that every strictly essentially \'{e}tale $A$-algebra be $A$-isomorphic to} $A$.

\begin{lemma}[18.6.3]
\label{IV.18.6.3}
Let $A$ be a local ring, $A_1$, $A_2$ two $A$-algebras that are local and strictly essentially \'{e}tale.
\begin{enumerate}[label=\emph{\roman*)}]
  \item There exists at most one $A$-homomorphism (necessarily local) from $A_1$ to $A_2$.
  \item There exists a local $A$-algebra $A_3$ that is strictly essentially \'{e}tale, and two $A$-homomorphisms $A_1\to A_3$, $A_2\to A_3$.
\end{enumerate}
\end{lemma}

\begin{proof}
Let $S=\Spec(A)$;
by hypothesis there are two $S$-schemes \'{e}tale over $S$, $X_1$, $X_2$, and two points $x_1\in X_1$, $x_2\in X_2$ above the closed point $s$ of $S$ and such that $A_1=\sh{O}_{X_1,x_1}$, $A_2=\sh{O}_{X_2,x_2}$;
Let $X_3=X_1\times_{S}X_2$;
the hypotheses $\boldsymbol{k}(x_1)=\boldsymbol{k}(x_2)=\boldsymbol{k}(s)$ \sref[I]{I.3.4.9} imply that there exists a unique point $x_3\in X_3$ above $x_1$ and $x_2$, and that $\boldsymbol{k}(x_3)=\boldsymbol{k}(s)$.
Furthermore, $X_3$ is \'{e}tale over $S$ \sref{IV.17.3.3}, so $A_3=\sh{O}_{X_3,x_3}$ satisfies the conditions of (ii).
On the other hand, \oldpage[IV-4]{137}we have seen that every $A$-homomorphism from $A_1$ to $A_2$ is necessarily local;
it corresponds to an $S$-morphism $f$ from $X_2'=\operatorname{Spec}(A_2)$ to $X_1$ such that $f(x_2)=x_1$, or, equivalently, on setting $X_3'=X_1\times_S X_2'$, to an $X_2'$-section $f'$ of $X_3'$ such that $f'(x_2)=x_3$.
Since $\boldsymbol{k}(x_3)=\boldsymbol{k}(x_2)$ and $X_2'$ is connected, the uniqueness of $f'$ follows from \sref{IV.17.4.9}, whence (i).
\end{proof}

\begin{env}[18.6.4]
\label{IV.18.6.4}
Let us denote by $\mathfrak{S}$ the subset of the set $\mathfrak{E}$ defined in \sref{IV.18.6.1} formed by the $A$-algebras that are strictly essentially \'{e}tale belonging to $\mathfrak{E}$.
It follows from \sref{IV.18.6.3} that the relation ``there exists an $A$-homomorphism from $A_1$ to $A_2$'' is a \emph{preorder} relation on $\mathfrak{S}$, and makes $\mathfrak{S}$ a \emph{filtered increasing} set.
Indexing $\mathfrak{S}$ by itself by means of the identity map, it follows from \sref{IV.18.6.3},~(i) that if $\lambda\leq\mu$ in $\mathfrak{S}$, there exists a \emph{unique} $\vphi_{\mu\lambda}:A_\lambda\to A_\mu$ and $(A_\lambda,\vphi_{\mu\lambda})$ is clearly an \emph{inductive system} of local $A$-algebras, the $\vphi_{\mu\lambda}$ being local homomorphisms.
Furthermore, it follows from \sref{IV.17.3.5} that for $\lambda\leq\mu$, $A_\mu$ is an $A_\lambda$-algebra that is strictly essentially \'{e}tale.
\end{env}

\begin{definition}[18.6.5]
\label{IV.18.6.5}
The Henselization of a local ring $A$ is the inductive limit of the inductive system $(A_\lambda,\vphi_{\mu\lambda})$, denoted ${}^h\!A$, defined in \sref{IV.18.6.4}.
\end{definition}

This definition does not depend on the choice of $\mathfrak{E}$; if $\mathfrak{E}'$ is another set having the same property as $\mathfrak{E}$, and $\mathfrak{S}'$ the subset of $\mathfrak{E}'$ consisting of the strictly essentially \'{e}tale $A$-algebras, it follows from \sref{IV.18.6.3},~(ii) that, for the preorder relations considered, $\mathfrak{S}$ and $\mathfrak{S}'$ are \emph{cofinal} in $\mathfrak{S}''=\mathfrak{S}\cup\mathfrak{S}'$, so they give the same inductive limit up to an isomorphism.
We shall moreover see below \sref{IV.18.6.6},~(i) and~(ii) that ${}^h\!A$ and the canonical homomorphism $A\to{}^h\!A$ are the solution of a universal problem, and are therefore determined up to a unique isomorphism.

One will observe that if $A$ is a field, then evidently ${}^h\!A=A$.
In the general case, ${}^h\!A$ is also the Henselization of all the $A$-algebras that are strictly essentially \'{e}tale $A_\lambda$ ($\lambda\in\mathfrak{S}$), since if $B$ is an $A_\lambda$-algebra that is strictly essentially \'{e}tale, $B$ is also an $A$-algebra that is strictly essentially \'{e}tale \sref{IV.18.6.1}, so (up to an $A$-isomorphism) one of the $A_\mu$ for $\mu\geq\lambda$, and the $A_\mu$ such that $\mu\geq\lambda$ and that $A_\mu$ is an $A_\lambda$-algebra that is strictly essentially \'{e}tale form a cofinal set in the preordered set of the $A_\lambda$, by virtue of \sref{IV.18.6.1} and \sref{IV.18.6.3}.

\begin{theorem}[18.6.6]
\label{IV.18.6.6}
Let $A$ be a local ring, ${}^h\!A$ its Henselization.
\begin{enumerate}[label=\emph{\roman*)}]
  \item ${}^h\!A$ is a Henselian local ring and the structural homomorphism $A\to{}^h\!A$ is local.
  \item For every Henselian local ring $B$, the canonical map
  \[
    \Hom.\operatorname{loc}({}^h\!A,B)\to\Hom.\operatorname{loc}(A,B)
  \]
  is bijective.
  \item ${}^h\!A$ is a faithfully flat $A$-module, and if $\mathfrak{m}$ is the maximal ideal of $A$, $\mathfrak{m}.{}^h\!A$ is the maximal ideal of ${}^h\!A$, and the homomorphism $A/\mathfrak{m}\to{}^h\!A/\mathfrak{m}.{}^h\!A$ of residue fields is bijective.
  \item If $\hat{A}$ and $({}^h\!A)\hat{\phantom{x}}$ are the separated completions of the local rings $A$ and ${}^h\!A$, the homomorphism $\hat{A}\to({}^h\!A)\hat{\phantom{x}}$ deduced from the structural homomorphism $A\to{}^h\!A$ by completion is bijective.
  \item For ${}^h\!A$ to be Noetherian, it is necessary and sufficient that $A$ be so.
  \item If $A$ is Henselian, the canonical homomorphism $A\to{}^h\!A$ is bijective.
\end{enumerate}
\end{theorem}

\begin{proof}
\oldpage[IV-4]{138}
Let $\mathfrak{m}_\lambda$ be the maximal ideal of $A_\lambda$;
it follows from the fact that, in \sref{IV.17.6.1}, \emph{a)} implies \emph{c$'$)}, that for $\lambda\leq\mu$ one has $\mathfrak{m}_\mu=\mathfrak{m}_\lambda A_\mu$, that the homomorphism $A_\lambda/\mathfrak{m}_\lambda\to A_\mu/\mathfrak{m}_\mu$ is bijective, and that $A_\mu$ is a flat $A_\lambda$-module.
The fact that ${}^h\!A$ is local, that the homomorphism $A\to{}^h\!A$ is local, assertion~(iii) and the sufficiency of~(v) follow therefore from \sref[0\textsubscript{\mdseries III}]{0.10.3.1.3};
the necessity of~(v) follows from the fact that ${}^h\!A$ is a faithfully flat $A$-module \sref[0\textsubscript{\mdseries I}]{0.6.5.2}.

To prove that ${}^h\!A$ is Henselian, let us apply the criterion \sref{IV.18.5.11},~\emph{b)}.
Let $S=\Spec({}^h\!A)$, $S_0=\Spec({}^h\!A/\mathfrak{m}.{}^h\!A)$, let $g:S'\to S$ be an \'{e}tale morphism, set $S_0'=g^{-1}(S_0)$, and let $f_0:S_0\to S_0'$ be an $S_0$-section.
reasoning as in \sref{IV.18.5.4}, one can suppose that $g$ is of finite presentation.
It follows then from \sref{IV.8.8.2} and \sref{IV.17.7.5} that there exist an index $\lambda\in\mathfrak{S}$, an \'etale morphism $g_\lambda:S'^{(\lambda)}\to S^{(\lambda)}=\Spec(A_\lambda)$ and, on setting $S_0^{(\lambda)}=\Spec(A_\lambda/\mathfrak{m}A_\lambda)$ and $S_0'^{(\lambda)}=g_\lambda^{-1}(S_0^{(\lambda)})$, an $S_0^{(\lambda)}$-section $f_0^{(\lambda)}:S_0^{(\lambda)}\to S_0'^{(\lambda)}$, such that
\[
S'=S'^{(\lambda)}\times_{S^{(\lambda)}}S,\qquad
g=g_\lambda\times 1,\qquad
f_0=f_0^{(\lambda)}\times 1.
\]
Let $s$ be the closed point of $S$, $x=f_0(s)$, let $x_\lambda$ be the projection of $x$ in $S'^{(\lambda)}$, and let $s_\lambda$ be the closed point of $S^{(\lambda)}$.
Since $x_\lambda$ lies above $s_\lambda$, the local ring $C_\lambda$ of $S'^{(\lambda)}$ at $x_\lambda$ is an essentially \'etale $A_\lambda$-algebra;
moreover, since $f_0^{(\lambda)}(s_\lambda)=x_\lambda$, we have $\boldsymbol{k}(x_\lambda)=\boldsymbol{k}(s_\lambda)$, in other words $C_\lambda$ is a local strictly essentially \'etale $A_\lambda$-algebra, and consequently, by \sref{IV.18.6.1}, it is $A_\lambda$-isomorphic to an $A$-algebra $A_\mu$ with $\mu\geq\lambda$.
There is therefore an $A_\lambda$-homomorphism $C_\lambda\to{}^h\!A$, that is, an $S^{(\lambda)}$-morphism $h:S\to S'^{(\lambda)}$ such that $h(s)=x_\lambda$, and hence an $S$-section $f$ of $S'$ such that $f(s)=x$;
this completes the proof of (i).

To prove~(ii), it suffices to note that $\Hom.\operatorname{loc}({}^h\!A,B)\cong\varprojlim\Hom.\operatorname{loc}(A_\lambda,B)$ and that by \sref{IV.18.6.2} the canonical homomorphisms $\Hom.\operatorname{loc}(A_\lambda,B)\leftarrow\Hom.\operatorname{loc}(A,B)$ are bijective.

To prove~(iv), note that for every $\lambda$ and every integer $n>0$, $\mathfrak{m}_\lambda^n=\mathfrak{m}^nA_\lambda$, and
\[
(\mathfrak{m}.{}^h\!A)^n=\mathfrak{m}^n.{}^h\!A=\varinjlim\mathfrak{m}_\lambda^n.
\]
By exactness of the functor $\varinjlim$ in the category of $A$-modules,
\[
{}^h\!A/(\mathfrak{m}.{}^h\!A)^n=\varinjlim(A_\lambda/\mathfrak{m}_\lambda^n),
\]
and it therefore suffices to show that, for every $n$ and every index $\lambda$, the homomorphism $A/\mathfrak{m}^n\to A_\lambda/\mathfrak{m}_\lambda^n$ is bijective.
Now, this holds by hypothesis for $n=1$; on the other hand, since $A_\lambda$ is a flat $A$-module, one has
\[
  \mathfrak{m}_\lambda^n/\mathfrak{m}_\lambda^{n+1}=(\mathfrak{m}^n/\mathfrak{m}^{n+1})\otimes_{A}A_\lambda
  =(\mathfrak{m}^n/\mathfrak{m}^{n+1})\otimes_{A/\mathfrak{m}}(A_\lambda/\mathfrak{m}A_\lambda)
\]
and as $A/\mathfrak{m}\to A_\lambda/\mathfrak{m}A_\lambda=A_\lambda/\mathfrak{m}_\lambda$ is bijective, the homomorphism $\mathfrak{m}^n/\mathfrak{m}^{n+1}\to\mathfrak{m}_\lambda^n/\mathfrak{m}_\lambda^{n+1}$ is also bijective;
the conclusion then follows from Bourbaki, \emph{Alg. comm.}, chap.~III, \S~2, no~8, cor.~3 of th.~1.

Finally, if $A$ is Henselian, it follows from the remark preceding \sref{IV.18.6.3} that the canonical homomorphisms $A\to A_\lambda$ are bijective, which proves~(vi) by definition of ${}^h\!A$.
\end{proof}

\begin{env}[18.6.7]
\label{IV.18.6.7}
Let $A$ be a \emph{semi-local} ring, $\mathfrak{m}_i$ ($1\leq i\leq r$) its maximal ideals.
The \emph{Henselization} of $A$ is the ring denoted again ${}^h\!A$, defined as the \emph{product} $\prod_i{}^h(A_{\mathfrak{m}_i})$ of the Henselizations of the local rings $A_{\mathfrak{m}_i}$.
It is a faithfully flat $A$-module and a semi-local $A$-algebra, whose maximal ideals are the $\mathfrak{m}_i.{}^h\!A$ by virtue of \sref{IV.18.6.6},~(iii);
furthermore, if $\mathfrak{r}=\bigcap\mathfrak{m}_i$ is the radical of $A$, it follows from the above that $\mathfrak{r}.{}^h\!A$
\oldpage[IV-4]{139}
is the radical of ${}^h\!A$ and that the canonical map $A/\mathfrak{r}\to{}^h\!A/\mathfrak{r}.{}^h\!A$ is bijective.
As the separated completion $\hat{A}$ of $A$ for the $\mathfrak{r}$-adic topology is the product of the separated completions $\hat{A}_{\mathfrak{m}_i}$ (Bourbaki, \emph{Alg. comm.}, chap.~III, \S~2, no~13, prop.~18), the canonical homomorphism $\hat{A}\to({}^h\!A)\hat{\phantom{x}}$ is bijective by \sref{IV.18.6.6},~(iv), and it is clear by \sref{IV.18.6.6},~(v) that for ${}^h\!A$ to be Noetherian, it is necessary and sufficient that $A$ be so.

To obtain the analogue of the universal property \sref{IV.18.6.6},~(ii), when $A$ and $B$ are two \emph{semi-local} rings of radicals $\mathfrak{r}$, $\mathfrak{s}$, let us agree to call a \emph{semi-local} homomorphism $\vphi$ from $A$ to $B$ any homomorphism such that $\vphi(\mathfrak{r})\subset\mathfrak{s}$.
In the particular case where $B$ is a \emph{product} of local rings $B_j$ ($1\leq j\leq n$), the data of such a homomorphism amounts to the data of its projections $\vphi_j:A\to B_j$, which are homomorphisms subject only to the condition that $\vphi_j(\mathfrak{r})\subset\mathfrak{n}_j$, where $\mathfrak{n}_j$ is the maximal ideal of $B_j$.
Moreover, since $\mathfrak{m}_i$ ($1\leq i\leq m$) are the maximal ideals of $A$, the prime ideal $\vphi_j^{-1}(\mathfrak{n}_j)$ must contain one of the $\mathfrak{m}_i$, since it must contain their intersection $\mathfrak{r}$, hence it is equal to one of the $\mathfrak{m}_i$ and $\vphi_j$ factorises as $A\to A_{\mathfrak{m}_i}\xrightarrow{\psi_{ij}}B_j$, where $\psi_{ij}$ is a \emph{local} homomorphism.
If $\Hom.\operatorname{sloc}(A,B)$ denotes the set of semi-local homomorphisms from $A$ to $B$, one can thus identify this set with $\prod_j\left(\bigcup_i\Hom.\operatorname{loc}(A_{\mathfrak{m}_i},B_j)\right)$;
in the case considered, since a semi-local Henselian ring is a direct product of Henselian local rings, one sees that the universal property \sref{IV.18.6.6},~(ii) is still valid for a semi-local ring $A$, provided one replaces ``local homomorphisms'' by ``semi-local homomorphisms'' \sref{IV.18.6.7}.
\end{env}

\begin{proposition}[18.6.8]
\label{IV.18.6.8}
Let $A$ be a semi-local ring, and let $B$ be a semi-local $A$-algebra integral over $A$.
Then $B\otimes_{A}({}^h\!A)$ is a semi-local $B$-algebra isomorphic to ${}^h\!B$.
\end{proposition}

\begin{proof}
Since $B\otimes_{A}({}^h\!A)$ is integral over ${}^h\!A$, each of its maximal ideals lies above a maximal ideal of ${}^h\!A$, hence above a maximal ideal of $A$.  Since $B$ is integral over $A$, the prime ideals of $B$ lying above a maximal ideal of $A$ are maximal ideals of $B$.  Thus the projections in $\Spec({}^h\!A)$ and $\Spec(B)$ of a maximal ideal of $B\otimes_{A}({}^h\!A)$ are both maximal ideals.  Taking account of \sref{IV.18.6.6},~(iii), and \sref[I]{I.3.4.9}, one concludes that $B\otimes_{A}({}^h\!A)$ is semi-local.

Furthermore, for every Henselian local ring $C$, $\Hom.\operatorname{sloc}(B\otimes_{A}({}^h\!A),C)$ is in bijection with the set of pairs $(\vphi,\psi)$ of semi-local homomorphisms $\vphi:B\to C$ and $\psi:{}^h\!A\to C$ such that the composites
\[
  A\longrightarrow B\xrightarrow{\vphi}C,
  \qquad
  A\longrightarrow{}^h\!A\xrightarrow{\psi}C
\]
are equal.  By the bijection between $\Hom.\operatorname{sloc}(A,C)$ and $\Hom.\operatorname{sloc}({}^h\!A,C)$, for every $\vphi\in\Hom.\operatorname{sloc}(B,C)$ there exists a unique $\psi$ having this property.  Hence the map
\[
  \Hom.\operatorname{sloc}(B\otimes_{A}({}^h\!A),C)
  \longrightarrow \Hom.\operatorname{sloc}(B,C)
\]
is bijective, which proves the proposition by the uniqueness of the solution of a universal problem.
\end{proof}

\begin{theorem}[18.6.9]
\label{IV.18.6.9}
Let $A$ be a semi-local ring.
\begin{enumerate}[label=\emph{\roman*)}]
  \item For ${}^h\!A$ to be reduced \textup{(resp.} normal\textup{)}, it is necessary and sufficient that $A$ be so.
  \item Suppose $A$ Noetherian.
  Then, for every prime ideal $\mathfrak{p}$ of $A$, the ring $({}^h\!A)_{\mathfrak{p}}/\mathfrak{p}.({}^h\!A)_{\mathfrak{p}}$ is a finite direct product of separable algebraic extensions of $\boldsymbol{k}(\mathfrak{p})$ \textup{(which implies that the fibres of the canonical morphism $\Spec({}^h\!A)\to\Spec(A)$ are \emph{geometrically regular})}.
\end{enumerate}
\end{theorem}

\begin{proof}
\oldpage[IV-4]{140}
(i) As ${}^h\!A$ is a faithfully flat $A$-module, it follows from \sref{IV.2.1.13} that if ${}^h\!A$ is reduced (resp. normal), so is $A$.
To prove the converse, one can restrict to the case where $A$ is local, so that, with the notations of \sref{IV.18.6.4}, one has ${}^h\!A=\varinjlim A_\alpha$.
Now, it follows from the definition of the $A_\alpha$ and from \sref{IV.17.5.7} that if $A$ is reduced (resp. integral and integrally closed), so is each of the $A_\alpha$;
furthermore, the homomorphisms $A_\alpha\to A_\beta$ for $\alpha\leq\beta$ are injective by virtue of \sref[0\textsubscript{\mdseries I}]{0.6.5.1} since $A_\beta$ is a faithfully flat $A_\alpha$-module;
hence the morphisms $\Spec(A_\beta)\to\Spec(A_\alpha)$ are dominant, and one concludes from \sref{IV.5.13.2} (resp. \sref{IV.5.13.4}) that ${}^h\!A$ is reduced (resp. integral and integrally closed).

(ii) One can restrict to the case where $A$ is local.
By virtue of the fact that the functor $\varinjlim$ commutes with tensor products, the fibre of the morphism $\Spec({}^h\!A)\to\Spec(A)$ at a point $\mathfrak{p}$ is the projective limit of the fibres of the morphisms $\Spec(A_\alpha)\to\Spec(A)$ at this point.
As ${}^h\!A$ is Noetherian, one is reduced to proving the following lemma:
\end{proof}

\begin{lemma}[18.6.9.1]
\label{IV.18.6.9.1}
Let $(B_\alpha)$ be an inductive system of Artinian rings, $B=\varinjlim B_\alpha$ its inductive limit.
If $B$ is Noetherian, then $B$ is Artinian;
if moreover, for $\alpha\leq\beta$, the homomorphisms $\vphi_{\beta\alpha}:B_\alpha\to B_\beta$ are injective, then there exists $\lambda$ such that, for $\alpha\geq\lambda$, the number of local components $B_\alpha^{(i)}$ of $B_\alpha$ is constant, that $\vphi_{\beta\alpha}(B_\alpha^{(i)})\subset B_\beta^{(i)}$ for every $i$ and that the $B^{(i)}=\varinjlim B_\alpha^{(i)}$ are the local components of $B$.
\end{lemma}

\begin{proof}
Indeed $Y_\alpha=\Spec(B_\alpha)$, $Y=\Spec(B)$, which, as a topological space, is equal to $\varprojlim Y_\alpha$ \sref{IV.8.2.10}.
The spaces $Y_\alpha$ are finite and discrete;
furthermore, if the $\vphi_{\beta\alpha}$ are injective, the canonical morphisms $Y_\beta\to Y_\alpha$ are dominant, hence surjective.
The lemma then follows from the following purely topological lemma:
\end{proof}

\begin{lemma}[18.6.9.2]
\label{IV.18.6.9.2}
Let $(Y_\alpha,\psi_{\alpha\beta})$ be a projective system of finite discrete topological spaces.
If $Y=\varprojlim Y_\alpha$ is Noetherian, $Y$ is finite and discrete.
If moreover the $\psi_{\alpha\beta}$ are surjective, then there exists $\lambda$ such that for $\alpha\geq\lambda$, the number of elements of $Y_\alpha$ is constant and equal to the number of elements of $Y$.
\end{lemma}

\begin{proof}
Indeed, since $Y$ is compact and Noetherian, every subset of $Y$ is compact, hence closed; thus $Y$ is discrete, and therefore finite since it is compact.
If moreover the $\psi_{\alpha\beta}$ are surjective, the same holds for $\psi_\alpha:Y\to Y_\alpha$, so $\operatorname{Card}(Y)\geq\operatorname{Card}(Y_\alpha)$, and as $\operatorname{Card}(Y_\alpha)$ is an increasing function of $\alpha$, there exists $\lambda$ such that for $\alpha\geq\lambda$, $\operatorname{Card}(Y_\alpha)=\operatorname{Card}(Y_\lambda)$;
the $\psi_{\alpha\beta}$ are then bijective for $\alpha\geq\lambda$ and one therefore has $\operatorname{Card}(Y)=\operatorname{Card}(Y_\lambda)$.
\end{proof}

\begin{corollary}[18.6.10]
\label{IV.18.6.10}
Let $A$ be a local Noetherian ring.
For ${}^h\!A$ to possess one of the following properties:
\begin{enumerate}[label=\emph{\alph*)}]
  \item being a Cohen-Macaulay ring \sref[0]{0.16.5.3};
  \item satisfying the property $(S_n)$ \sref{IV.5.7.3};
  \item being regular;
  \item satisfying the property $(R_n)$ \sref{IV.5.8.2};
\end{enumerate}
it is necessary and sufficient that $A$ possess the same property.
\end{corollary}

\begin{proof}
\oldpage[IV-4]{141}
It follows immediately from the fact that the fibres of the morphism $\Spec({}^h\!A)\to\Spec(A)$ are geometrically regular \sref{IV.18.6.9} and from \sref{IV.6.4.1},~\sref{IV.6.5.1} and \sref{IV.6.5.3}.
\end{proof}

\begin{corollary}[18.6.11]
\label{IV.18.6.11}
Let $A$ be a semi-local Noetherian ring.
For a prime ideal $\mathfrak{p}$ of ${}^h\!A$ to belong to $\operatorname{Ass}({}^h\!A)$, it is necessary and sufficient that $\mathfrak{p}\cap A$ belong to $\operatorname{Ass}(A)$.
\end{corollary}

\begin{proof}
Taking account of the description of the fibres of the morphism $\Spec({}^h\!A)\to\Spec(A)$, this also follows immediately from \sref{IV.3.3.1}.
\end{proof}

\begin{proposition}[18.6.12]
\label{IV.18.6.12}
Let $A$ be a local ring;
the following properties are equivalent:
\begin{enumerate}[label=\emph{\alph*)}]
  \item $A$ is unibranch \sref[0]{0.23.2.1} \textup{(resp.} reduced and unibranch\textup{)}.
  \item For every $A$-algebra $A'$ that is local and strictly essentially \'{e}tale, $\Spec(A')$ is irreducible \textup{(resp.} integral\textup{)}.
  \item $\Spec({}^h\!A)$ is irreducible \textup{(resp.} integral\textup{)}.
\end{enumerate}
\end{proposition}

\begin{proof}
First note that ${}^h(A_{\mathrm{red}})=({}^h\!A)_{\mathrm{red}}$: indeed, ${}^h(A_{\mathrm{red}})=({}^h\!A)\otimes_{A}A_{\mathrm{red}}$ by \sref{IV.18.6.8}, and since ${}^h(A_{\mathrm{red}})$ is reduced by \sref{IV.18.6.9}, it is equal to $({}^h\!A)_{\mathrm{red}}$.
This allows us, in the proof, to consider only the case where $A$ is \emph{reduced}; then so are $A'$ \sref{IV.17.5.7} and ${}^h\!A$ \sref{IV.18.6.9}.

With the notations of \sref{IV.18.6.4}, condition \emph{b)} means that all the $A_\alpha$ are integral; one therefore concludes that ${}^h\!A=\varinjlim A_\alpha$ is integral \sref{IV.5.13.3}, so \emph{b)} implies \emph{c)}.

It is clear that if ${}^h\!A$ is integral, so are each of the $A_\alpha\subset{}^h\!A$, hence \emph{c)} implies \emph{b)}.

Let us show that \emph{c)} implies \emph{a)}.
Note first that $A\subset{}^h\!A$ and that ${}^h\!A$ is an integral domain;
let $K$ and $L$ be the fields of fractions of $A$ and ${}^h\!A$ respectively.
It suffices to show that for every \emph{finite} $A$-algebra $B$ contained in $K$, the ring $B$ is local.  Now $B$ is semi-local, and its Henselization is $B\otimes_{A}({}^h\!A)$ by \sref{IV.18.6.8}.  By flatness, $B\otimes_{A}({}^h\!A)$ identifies with a subring of $K\otimes_{A}({}^h\!A)$, and the latter with a subring of $L$; hence ${}^h\!B$ is integral.  But since ${}^h\!B$ is a semi-local Henselian ring, it is a finite direct product of local rings \sref{IV.18.5.9}; it can therefore be integral only if it is local, which indeed implies that $B$ is local.

Let us finally prove that \emph{a)} implies \emph{c)}.
Let $C$ be the integral closure of the integral ring $A$.
As $C$ is a local ring by hypothesis, $C\otimes_{A}({}^h\!A)$ is the Henselization of $C$ \sref{IV.18.6.8}, hence is integral and integrally closed \sref{IV.18.6.9}.
Since $A$ is a subring of $C$, flatness identifies ${}^h\!A$ with a subring of ${}^h\!C$; it is therefore integral as well, which completes the proof.
\end{proof}

\begin{corollary}[18.6.13]
\label{IV.18.6.13}
Let $A$ be a local Henselian ring.
Then, for $A$ to be unibranch, it is necessary and sufficient that $A_{\mathrm{red}}$ be integral.
\end{corollary}

\begin{proposition}[18.6.14]
\label{IV.18.6.14}
\begin{enumerate}[label=\emph{\roman*)}]
  \item Every filtered inductive limit of Henselian local rings (where the transition homomorphisms are local) is a Henselian local ring.
  \item The functor $A\leadsto{}^h\!A$ in the category of local rings (where the morphisms are local homomorphisms) commutes with filtered inductive limits.
  \item Every Henselian local ring $A$ is an inductive limit of an inductive filtered system of Henselian and Noetherian local rings $A_\lambda$ (the transition homomorphisms $A_\lambda\to A_\mu$ being local).
\end{enumerate}
\end{proposition}

\begin{proof}
(i) Let $(A_\lambda)$ be a filtered inductive system of Henselian local rings, where the transition homomorphisms are local.
Let us show that $A'=\varinjlim A_\lambda$, which is local \sref[0\textsubscript{\mdseries III}]{0.10.3.1.3}, is Henselian.
Now, let $C$ be an \'{e}tale $A'$-algebra, $\mathfrak{n}$ a prime ideal
\oldpage[IV-4]{142}
of $C$ above the maximal ideal $\mathfrak{m}'$ of $A'$, such that $C_{\mathfrak{n}}$ is strictly essentially \'{e}tale over $A'$ \sref{IV.18.6.1}, in other words such that $\boldsymbol{k}(\mathfrak{m}')\to\boldsymbol{k}(\mathfrak{n})$ is an isomorphism. By virtue of \sref{IV.8.8.2},~(ii) and \sref{IV.17.7.8}, there exist an index $\lambda$ and an \'{e}tale $A_\lambda$-algebra $C_\lambda$ such that $C=C_\lambda\otimes_{A_\lambda}A'$ up to an isomorphism.
Furthermore, if $\mathfrak{n}_\lambda$ is the inverse image of $\mathfrak{n}$ in $C_\lambda$, and $\mathfrak{m}_\lambda$ the maximal ideal of $A_\lambda$, one can, by virtue of \sref{IV.8.8.2.4} and of the transitivity of fibres \sref[I]{I.3.6.4}, suppose that $(C_\lambda)_{\mathfrak{n}_\lambda}$ is strictly essentially \'{e}tale over $A_\lambda$;
since $A_\lambda$ is Henselian, $(C_\lambda)_{\mathfrak{n}_\lambda}$ is necessarily isomorphic to $A_\lambda$ \sref{IV.18.5.11},~\emph{b)};
there is therefore a neighbourhood of $\mathfrak{n}_\lambda$ in $\Spec(C_\lambda)$ isomorphic to $\Spec(A_\lambda)$ \sref[I]{I.6.5.4};
one concludes that there is a neighbourhood of $\mathfrak{n}$ in $\Spec(C)$ isomorphic to $\Spec(A')$, hence that $C_\mathfrak{n}$ is isomorphic to $A'$.  This proves that $A'$ is Henselian \sref{IV.18.6.2} and completes the proof.

(ii) Suppose that $B=\varinjlim B_\lambda$, where the $B_\lambda$ are local rings, the transition morphisms $B_\lambda\to B_\mu$ ($\lambda\leq\mu$) being local.
Set $A_\lambda={}^h\!B_\lambda$;
the $A_\lambda$ are Henselian local rings \sref{IV.18.6.6},~(vi), and by virtue of \sref{IV.18.6.6},~(ii), the transition morphisms $B_\lambda\to B_\mu$ uniquely determine local homomorphisms $A_\lambda\to A_\mu$, so that $(A_\lambda)$ is an inductive system.
It amounts to showing that $A'=\varinjlim A_\lambda$ is canonically isomorphic to $A={}^h\!B$.
By virtue of \sref{IV.18.6.6},~(ii) and of the definition of inductive limits, one has, for every Henselian local ring $E$,
\[
  \Hom.\operatorname{loc}(A',E)=\varprojlim\Hom.\operatorname{loc}(A_\lambda,E)=\varprojlim\Hom.\operatorname{loc}(B_\lambda,E)=\Hom.\operatorname{loc}(B,E).
\]
But as $A'$ is Henselian by~(i), this proves our assertion.

(iii) The ring $A$ is a filtered inductive limit of its Noetherian local subrings $B_\lambda$, the transition homomorphisms being local \sref{IV.5.13.3},~(iii).
As the ${}^h\!B_\lambda$ are Henselian and Noetherian local rings \sref{IV.18.6.6},~(v) and that $A={}^h\!A$ \sref{IV.18.6.6},~(vi), it suffices to apply~(ii) to the inductive system $(B_\lambda)$.
\end{proof}

\begin{corollary}[18.6.15]
\label{IV.18.6.15}
Let $A$ be a Henselian local ring, $X$ an $A$-prescheme of finite presentation over $A$.
Then there exists a Noetherian Henselian local ring $A_0$, a local homomorphism $A_0\to A$, an $A_0$-prescheme $X_0$ of finite type over $A_0$ and an $A$-isomorphism $X_0\otimes_{A_0}A\simto X$.
\end{corollary}

\begin{proof}
This follows from \sref{IV.18.6.14} and \sref{IV.8.8.2},~(ii).
\end{proof}


\subsection{Henselization and excellent rings.}
\label{subsection:IV.18.7}

\begin{env}[18.7.1]
\label{IV.18.7.1}
In this paragraph we denote by $\boldsymbol{P}(Z,k)$ a property of the form considered in \sref{IV.7.3.1}, where we further suppose that the property $\boldsymbol{Q}(A,k)$ satisfies the following condition:

For every separable algebraic extension $k'$ of $k$, and every local Noetherian $k'$-algebra $A$, the property $\boldsymbol{Q}(A,k)$ is \emph{equivalent} to $\boldsymbol{Q}(A,k')$.

It is immediate that all the properties $\boldsymbol{P}$ considered in \sref{IV.7.3.8} satisfy the preceding condition, since if $K$ is a finite extension of $k$, $k'\otimes_kK$ is a finite direct product of separable algebraic extensions of $K$.
\end{env}

\oldpage[IV-4]{143}

\begin{proposition}[18.7.2]
\label{IV.18.7.2}
Let $A$ be a local Noetherian ring.
For ${}^h\!A$ to be a $\boldsymbol{P}$-ring \sref{IV.7.3.13}, it is necessary and sufficient that $A$ be so.
\end{proposition}

\begin{proof}
Indeed, by virtue of \sref{IV.18.6.6},~(iv), the completion $\hat{A}$ of $A$ is also the completion of ${}^h\!A$ and one therefore has $A\subset{}^h\!A\subset\hat{A}$.
By \sref{IV.18.6.9}, for every $x\in\Spec(A)$, the fibre at $x$ of the morphism $\Spec({}^h\!A)\to\Spec(A)$ is discrete and finite, and at each point $y_i$ of this fibre, $\boldsymbol{k}(y_i)$ is a separable algebraic extension of $\boldsymbol{k}(x)$.
The formal fibre of $A$ at $x$ is therefore the prescheme \emph{sum} of the formal fibres of ${}^h\!A$ at the points $y_i$.
The conclusion then follows from the hypothesis on $\boldsymbol{P}$ in \sref{IV.18.7.1}.
\end{proof}

\begin{corollary}[18.7.3]
\label{IV.18.7.3}
Let $A$ be a local Noetherian ring.
For ${}^h\!A$ to be universally Japanese, it is necessary and sufficient that $A$ be so.
\end{corollary}

\begin{proof}
Indeed, for a Noetherian local ring $B$, it amounts to the same to say that $B$ is universally Japanese, or that the formal fibres of $B$ are geometrically reduced \sref{IV.7.6.4} and \sref{IV.7.7.1};
the conclusion therefore follows from \sref{IV.18.7.2}.
\end{proof}

\begin{corollary}[18.7.4]
\label{IV.18.7.4}
Let $A$ be a local Noetherian ring.
For the formal fibres of ${}^h\!A$ to be geometrically regular, it is necessary and sufficient that those of $A$ be so.
\end{corollary}

\begin{proof}
This is a particular case of \sref{IV.18.7.2}.
\end{proof}

\begin{proposition}[18.7.5]
\label{IV.18.7.5}
Let $A$ be a local Noetherian ring.
If $A$ is universally catenary \sref{IV.5.6.2} \textup{(resp.} formally catenary \textup{\sref{IV.7.1.9}, resp.} strictly formally catenary \textup{\sref{IV.7.2.6})}, so is ${}^h\!A$.
\end{proposition}

\begin{proof}
With the notations of \sref{IV.18.6.4}, if $A$ is universally catenary, the same holds for the $A_\alpha$, which are $A$-algebras essentially of finite type \sref{IV.5.6.3}.
To prove that in this case ${}^h\!A$ is universally catenary, it will therefore suffice to establish the following lemma:
\end{proof}

\begin{lemma}[18.7.5.1]
\label{IV.18.7.5.1}
Let $(B_\alpha,\vphi_{\beta\alpha})$ be an inductive system of Noetherian rings and suppose that, for $\alpha\leq\beta$, $\vphi_{\beta\alpha}$ makes $B_\beta$ a faithfully flat $B_\alpha$-module and that $B=\varinjlim B_\alpha$ is Noetherian.
Then, if the $B_\alpha$ are catenary \textup{(resp.} universally catenary\textup{)}, the same holds for $B$.
\end{lemma}

\begin{proof}
To say that $B$ is universally catenary is equivalent to saying that $B[T_1,\ldots,T_n]$ is catenary for all $n$ \sref{IV.5.6.2}, and it is immediate that the inductive system $(B_\alpha[T_1,\ldots,T_n])$ satisfies the same conditions as $(B_\alpha)$;
it therefore suffices to prove that if the $B_\alpha$ are catenary, then $B$ is catenary.
If $\mathfrak{p}_0\supset\mathfrak{p}_1\supset\ldots\supset\mathfrak{p}_r$ is a saturated chain of prime ideals, the inverse images $\mathfrak{p}_{0\alpha}\supset\mathfrak{p}_{1\alpha}\supset\ldots\supset\mathfrak{p}_{r\alpha}$ of these ideals in $B_\alpha$ form a chain of prime ideals, and it suffices to show that for $\alpha$ sufficiently large, this chain is \emph{saturated}.
Now, each $\mathfrak{p}_i$ is an ideal of finite type, hence there exists $\lambda$ such that, for $\alpha\geq\lambda$, one has $\mathfrak{p}_i=\mathfrak{p}_{i\alpha}.B$.
One therefore has $1=\codim(V(\mathfrak{p}_i),V(\mathfrak{p}_{i+1}))=\codim(V(\mathfrak{p}_{i\alpha}),V(\mathfrak{p}_{i+1,\alpha}))$ since $B$ is a flat $B_\alpha$-module \sref{IV.6.1.4} and \sref[0\textsubscript{\mdseries I}]{0.6.2.3}, which completes the proof of the lemma.

Suppose now $A$ formally catenary.
Let $\mathfrak{p}$ be a prime ideal of ${}^h\!A$, and let $\mathfrak{q}$ be its inverse image in $A$, so that $\mathfrak{q}.{}^h\!A\subset\mathfrak{p}$ and consequently ${}^h\!A/\mathfrak{p}$ is a quotient ring of ${}^h\!A/\mathfrak{q}.{}^h\!A$.  It therefore suffices to prove that ${}^h\!A/\mathfrak{q}.{}^h\!A$ is formally catenary for every prime ideal $\mathfrak{q}$ of $A$ \sref{IV.7.1.9}.  Since ${}^h\!A/\mathfrak{q}.{}^h\!A={}^h(A/\mathfrak{q})$ by \sref{IV.18.6.8}, it suffices \sref{IV.7.1.11} to show that if $A$ is \emph{integral} and formally catenary, ${}^h\!A$ is formally \emph{equidimensional}.
But as the completion of ${}^h\!A$ is equal to $\hat{A}$, this follows from the hypothesis that $A$ is formally catenary.

\oldpage[IV-4]{144}

Finally, suppose that $A$ is strictly formally catenary.  We have just seen that ${}^h\!A$ is formally catenary, and it remains to prove that the fibres of the morphism $\Spec(\hat{A})\to\Spec({}^h\!A)$ satisfy property $(S_1)$ \sref{IV.7.2.5},~\emph{b)}.  This follows from the hypothesis on $A$ and from \sref{IV.18.7.2}, applied to the case where $\boldsymbol{P}(Z,k)$ is property $(S_1)$ for $Z$.
\end{proof}

\begin{corollary}[18.7.6]
\label{IV.18.7.6}
If a local Noetherian ring $A$ is excellent \sref{IV.7.8.2}, so is ${}^h\!A$.
\end{corollary}

\begin{remark}[18.7.7]
\label{IV.18.7.7}
It can happen that ${}^h\!A$ is an excellent local ring without $A$ being universally catenary (nor \emph{a fortiori} excellent).
To see this, let us take the example of $A$ from \sref{IV.5.6.11}, with the same notations.
The ring $E$ constructed in \sref{IV.5.6.11} is excellent when $k_0$ is of characteristic 0 \sref{IV.7.8.3},~(ii) and~(iii);
consider two rings $E_1$, $E_2$ isomorphic to $E$, and ``glue'' $\Spec(E_1)$ and $\Spec(E_2)$ such that if $\mathfrak{m}_i$ and $\mathfrak{m}'_i$ are the maximal ideals of $E_i$ ($i=1,2$), corresponding to $\mathfrak{m}$ and $\mathfrak{m}'$, the point $\mathfrak{m}_1$ is ``glued'' to $\mathfrak{m}'_2$ and the point $\mathfrak{m}_2$ to $\mathfrak{m}'_1$;
to be precise, if $\varepsilon_i$ and $\varepsilon'_i$ are the canonical homomorphisms from $E_i$ onto $\boldsymbol{k}(\mathfrak{m}_i)$ and $\boldsymbol{k}(\mathfrak{m}'_i)$ ($i=1,2$), the scheme obtained is $\Spec(R)$, where $R$ is the subring of $E_1\times E_2$ formed by the pairs $(x_1,x_2)$ such that $\varepsilon'_2(x_2)=\sigma(\varepsilon_1(x_1))$ and $\varepsilon'_1(x_1)=\sigma(\varepsilon_2(x_2))$.
One verifies easily (for example with the help of \sref{IV.17.6.3}) that $R$ is a finite \'{e}tale $C$-algebra, and that there are two maximal ideals $\mathfrak{r}_1$, $\mathfrak{r}_2$ of $R$ above the maximal ideal $\mathfrak{n}$ of $C$, the completions of $R_{\mathfrak{r}_1}$ and $R_{\mathfrak{r}_2}$ being canonically isomorphic to that of $A=C_{\mathfrak{n}}$.  Thus $R_{\mathfrak{r}_1}$ is an $A$-algebra that is strictly essentially \'{e}tale, and consequently its Henselization is equal to ${}^h\!A$.
Furthermore, we saw in \sref{IV.7.8.4},~(ii), that the formal fibres of $A$ are geometrically regular; the same therefore holds for those of ${}^h\!A$ by \sref{IV.18.7.2}.  Finally, $R$ is universally catenary, since its quotients by its two minimal prime ideals are isomorphic to $E$, the conclusion following from \sref{IV.5.6.3},~(iii).  Hence $R_{\mathfrak{r}_1}$ is excellent, and consequently so is ${}^h\!A$ by \sref{IV.18.7.6}, whereas $A$ is not universally catenary.
\end{remark}


\subsection{Strictly local rings and strict Henselization.}
\label{subsection:IV.18.8}

\begin{proposition}[18.8.1]
\label{IV.18.8.1}
Let $A$ be a local ring.
The following conditions are equivalent:
\begin{enumerate}[label=\emph{\alph*)}]
  \item $A$ is a Henselian ring and its residue field $k$ is separably closed.
  \item $A$ is a Henselian ring and every \'{e}tale covering of $S=\Spec(A)$ is trivial.
  \item For every \'{e}tale morphism $f:X\to S$ and every point $x\in X$ such that $f(x)$ is the closed point $s$ of $S$, there exists an $S$-section $u:S\to X$ of $X$ such that $u(s)=x$.
\end{enumerate}
\end{proposition}

\begin{proof}
The hypothesis \emph{c)} implies that for every \'{e}tale morphism $f:X\to S$, the residue fields at the points of $f^{-1}(s)$ are all isomorphic to $k$, and moreover the condition \emph{b)} of \sref{IV.18.5.11} is satisfied.
Hence $A$ is Henselian, and the fact that every \'{e}tale covering of $S$ is trivial follows from \sref{IV.18.2.10},~(ii);
hence \emph{c)} implies \emph{b)}.
As \'{e}tale coverings of $\Spec(k)$ are trivial if and only if $k$ is separably closed, \emph{b)} implies \emph{a)} by virtue of \sref{IV.18.5.11},~\emph{b)}.
Finally, it follows also from \sref{IV.18.5.11},~\emph{b)} that \emph{a)} implies \emph{c)}.
\end{proof}

\begin{definition}[18.8.2]
\label{IV.18.8.2}
A local ring is said to be \emph{strictly local} if it satisfies the equivalent conditions of \sref{IV.18.8.1}.
A scheme isomorphic to the spectrum of a strictly local ring is called a \emph{strictly local} scheme.
\end{definition}

\begin{remark}[18.8.3]
\label{IV.18.8.3}
The conditions of \sref{IV.18.8.1} are also equivalent to the following:
\begin{enumerate}
  \item[\textup{d)}] \itshape For every morphism locally of finite type $f:X\to S$ and every point $x\in X$ such that $f(x)$ is the closed point $s$ of $S$ and that $\sh{O}_{X,x}/\mathfrak{m}_s\sh{O}_{X,x}$ is a field, a finite separable extension of $\boldsymbol{k}(s)=k$, there exists an open neighbourhood $U$ of $x$ in $X$ such that $f|U$ is a closed immersion.
\end{enumerate}
\end{remark}

\oldpage[IV-4]{145}

(One will note that the hypothesis of \emph{d)} is satisfied if $f$ is \emph{formally unramified} at the point $x$ \sref{IV.17.4.1.2}.)

We leave the reader the proof, which is substantially the same as that of \sref{IV.18.5.18}.

\begin{lemma}[18.8.4]
\label{IV.18.8.4}
Let $A$, $A'$ be two local rings, $\vphi:A\to A'$ a local homomorphism making $A'$ an $A$-algebra essentially \'{e}tale \sref{IV.18.6.1}.
Then, for every strictly local ring $B$, every local homomorphism $\psi:A\to B$ and every homomorphism of $\boldsymbol{k}(A)$-algebras $\alpha:\boldsymbol{k}(A')\to\boldsymbol{k}(B)$, there exists one and only one local homomorphism $\psi':A'\to B$ such that $\psi=\psi'\circ\vphi$ and that $\alpha$ is equal to the homomorphism $\overline{\psi'}$ deduced from $\psi'$ by passage to quotients.
\end{lemma}

\begin{proof}
With the notations of \sref{IV.18.6.2}, the homomorphism $\alpha$ corresponds to a $\boldsymbol{k}(s)$-morphism $\Spec(\boldsymbol{k}(y))\to\Spec(\boldsymbol{k}(x))$, hence to a point $x'$ well determined in $X'$ above $y$.
The existence of a $Y$-section of $f':X'\to Y$ taking $y$ to $x'$ follows from \sref{IV.18.8.1},~\emph{c)}, since $f'$ is \'{e}tale and $B$ is strictly local; its uniqueness follows from the fact that $f'$ is separated and from \sref{IV.17.4.9}.
\end{proof}

\begin{lemma}[18.8.5]
\label{IV.18.8.5}
Let $A$ be a local ring, $A_1$, $A_2$ two $A$-algebras that are local and essentially \'{e}tale, $K$ a field that is an extension of $\boldsymbol{k}(A)$.
\begin{enumerate}[label=\emph{\roman*)}]
  \item For every $\boldsymbol{k}(A)$-homomorphism $\gamma:\boldsymbol{k}(A_1)\to\boldsymbol{k}(A_2)$, there exists at most one $A$-homomorphism (necessarily local) $\psi:A_1\to A_2$ such that $\gamma$ is the homomorphism $\overline{\psi}$ deduced from $\psi$ by passage to quotients.
  \item For every pair of local $A$-homomorphisms $\beta_1:A_1\to K$, $\beta_2:A_2\to K$ there exists a local $A$-algebra $A_3$ that is essentially \'{e}tale, two $A$-homomorphisms $\vphi_1:A_1\to A_3$, $\vphi_2:A_2\to A_3$ and an $A$-homomorphism $\beta:A_3\to K$ such that $\beta_1=\beta\circ\vphi_1$ and $\beta_2=\beta\circ\vphi_2$.
\end{enumerate}
\end{lemma}

\begin{proof}
With the notations of \sref{IV.18.6.3}, the homomorphisms $\beta_1$, $\beta_2$ in~(ii) correspond to two $S$-morphisms $\Spec(K)\to X_1$, $\Spec(K)\to X_2$, with respective images $x_1$, $x_2$;
one thus deduces a well-determined $S$-morphism $\Spec(K)\to X_3$ \sref[I]{I.3.2.1} whose image $x_3$ lies above $x_1$ and $x_2$;
the $A$-algebra $A_3=\sh{O}_{X_3,x_3}$ and the homomorphism $A_3\to K$ corresponding to these data then satisfy the conditions of~(ii).
On the other hand, the data of $\gamma$ in~(i) corresponds to an $S$-morphism $\Spec(\boldsymbol{k}(x_2))\to\Spec(\boldsymbol{k}(x_1))$ or also to a point $x'_3$ of $X'_3$ above $x_2$;
the uniqueness of $\psi$ follows from the uniqueness of an $X'_2$-section of $X'_3$ passing through $x'_3$ \sref{IV.17.4.9}.
\end{proof}

\begin{env}[18.8.6]
\label{IV.18.8.6}
Let $A$ be a local ring, $\mathfrak{m}$ its maximal ideal, $k$ its residue field, $\Omega$ a \emph{separable closure} of $k$.
Consider the set $\mathfrak{E}$ of essentially \'{e}tale $A$-algebras defined in \sref{IV.18.6.1}, and denote by $L$ the set of \emph{local} $A$-\emph{homomorphisms} $A'\to\Omega$, where $A'$ ranges over $\mathfrak{E}$;
one will note that such a homomorphism factorises as $\lambda:A'\to\boldsymbol{k}(A')\xrightarrow{\omega_\lambda}\Omega$, where $\omega_\lambda:\boldsymbol{k}(A')\to\Omega$ is a $k$-homomorphism (since $\mathfrak{m}A'$ is the maximal ideal of $A'$);
conversely the data of such a homomorphism $\omega_\lambda$ uniquely determines a homomorphism $\lambda\in L$.
For every $\lambda\in L$, we denote by $A_\lambda$ the essentially \'{e}tale $A$-algebra belonging to $\mathfrak{E}$ on which $\lambda$ is defined.
The set $L$ is \emph{preordered} by the relation ``there exists an $A$-homomorphism $\vphi_{\mu\lambda}:A_\lambda\to A_\mu$ such that $\lambda=\mu\circ\vphi_{\mu\lambda}$'', and it follows from \sref{IV.18.8.5},~(i) that this homomorphism $\vphi_{\mu\lambda}$ is \emph{unique}.
Furthermore, $L$ is \emph{filtered increasing} for the preceding preorder relation (which one also denotes $\lambda\leq\mu$), by virtue
\oldpage[IV-4]{146}
of \sref{IV.18.8.5},~(ii).
It is clear that $(A_\lambda,\vphi_{\mu\lambda})$ is an \emph{inductive system} of local $A$-algebras, the $\vphi_{\mu\lambda}$ being local homomorphisms;
furthermore, it follows from \sref{IV.17.3.5} that for $\lambda\leq\mu$, $A_\mu$ is an $A_\lambda$-algebra that is essentially \'{e}tale.
If $\overline{\vphi}_{\mu\lambda}:\boldsymbol{k}(A_\lambda)\to\boldsymbol{k}(A_\mu)$ is the $k$-homomorphism deduced from $\vphi_{\mu\lambda}$ by passage to quotients, $(\boldsymbol{k}(A_\lambda),\overline{\vphi}_{\mu\lambda})$ is an \emph{inductive system} of separable algebraic extensions of $k$, and the $\omega_\lambda:\boldsymbol{k}(A_\lambda)\to\Omega$ form an inductive system of $k$-homomorphisms.
Furthermore, the inductive limit
\begin{equation}
  \omega:\varinjlim\boldsymbol{k}(A_\lambda)\to\Omega
  \label{IV.18.8.6.1}
\tag{18.8.6.1}
\end{equation}
is a $k$\emph{-isomorphism}.
It suffices indeed to prove that for every finite separable extension $k'$ of $k$, there exists an $A$-algebra $A'$ that is essentially \'{e}tale, such that $k'$ is the residue field of $A'$ and that the homomorphism $k\to k'$ is deduced from $A\to A'$ by passage to quotients.
But this follows from \sref{IV.18.1.1} applied to \'{e}tale morphisms, since $\Spec(A)$ is the unique neighbourhood of its closed point.

Note now that if $\mathfrak{m}_\lambda$ is the maximal ideal of $A_\lambda$, it follows from \sref{IV.17.6.1} that, for $\lambda\leq\mu$, $\mathfrak{m}_\mu=\mathfrak{m}_\lambda A_\mu$ and that $A_\mu$ is a flat $A_\lambda$-module.
Hence it follows from \sref[0\textsubscript{\mdseries III}]{0.10.3.1.3} that the ring $\varinjlim A_\lambda$ is \emph{local}, that the homomorphism $w:A\to\varinjlim A_\lambda$ is local, and that one has a canonical $k$-isomorphism
\[
  \varinjlim\boldsymbol{k}(A_\lambda)\xrightarrow{\sim}\boldsymbol{k}(\varinjlim A_\lambda);
\]
identifying these two fields, one deduces a canonical $k$-isomorphism
\begin{equation}
  \omega^{-1}:\Omega\xrightarrow{\sim}\boldsymbol{k}(\varinjlim A_\lambda).
  \label{IV.18.8.6.2}
\tag{18.8.6.2}
\end{equation}
\end{env}

\begin{definition}[18.8.7]
\label{IV.18.8.7}
Let $A$ be a local ring, $k=\boldsymbol{k}(A)$ its residue field, $i:k\to\Omega$ a homomorphism from $k$ into a separable closure $\Omega$ of $k$;
the \emph{strict Henselization} of $A$ \emph{relative to} $i$, denoted ${}^{hs}\!A_{(i)}$ (or ${}^{hs}\!A$ when this causes no confusion), is the inductive limit of the inductive system $(A_\lambda,\vphi_{\mu\lambda})$ defined in \sref{IV.18.8.6}.
\end{definition}

As one has a $k$-isomorphism $\omega^{-1}:\Omega\xrightarrow{\sim}\boldsymbol{k}({}^{hs}\!A_{(i)})$, $\Omega$ is endowed canonically with the structure of a ${}^{hs}\!A_{(i)}$-algebra by the local homomorphism
\[
  {}^{hs}\!A_{(i)}\to\boldsymbol{k}({}^{hs}\!A_{(i)})\xrightarrow{\overset{\omega}{\sim}}\Omega.
\]

As in \sref{IV.18.6.5} one sees that the definition given in \sref{IV.18.8.7} does not depend on the choice of $\mathfrak{E}$;
we shall moreover see below \sref{IV.18.8.8},~(i) and~(ii) that ${}^{hs}\!A_{(i)}$ is an object representing a functor entirely defined by the data of $A$ and $i$ \sref[0\textsubscript{\mdseries III}]{0.8.1.8}, and is hence defined as such \emph{up to a unique isomorphism}.

\begin{proposition}[18.8.8]
\label{IV.18.8.8}
Let $A$ be a local ring, $k$ its residue field, $i:k\to\Omega$ a homomorphism from $k$ into a separable closure $\Omega$ of $k$, ${}^{hs}\!A_{(i)}$ the strict Henselization of $A$ corresponding to $i$.
\begin{enumerate}[label=\emph{\roman*)}]
  \item ${}^{hs}\!A_{(i)}$ is a strictly local ring \sref{IV.18.8.2} and the structural homomorphism $w:A\to{}^{hs}\!A_{(i)}$ is local.
  \item For every local homomorphism $u:A\to B$, where $B$ is a strictly local ring, and every $k$-homomorphism $\psi:\Omega\to\boldsymbol{k}(B)$, there exists one and only one $A$-homomorphism $v:{}^{hs}\!A_{(i)}\to B$ such that $\psi$ factorises as $\Omega\xrightarrow{\omega^{-1}}\boldsymbol{k}({}^{hs}\!A_{(i)})\xrightarrow{\overline{v}}\boldsymbol{k}(B)$, where $\overline{v}$ is deduced from $v$ by passage to quotients.
  \oldpage[IV-4]{147}
  \item ${}^{hs}\!A_{(i)}$ is a faithfully flat $A$-module, and if $\mathfrak{m}$ is the maximal ideal of $A$, $\mathfrak{m}.{}^{hs}\!A_{(i)}$ is the maximal ideal of ${}^{hs}\!A_{(i)}$, and $\boldsymbol{k}({}^{hs}\!A_{(i)})$ is a separable closure of $k$.
  \item For ${}^{hs}\!A_{(i)}$ to be Noetherian, it is necessary and sufficient that $A$ be so.
  \item If $A$ is strictly local (hence $\Omega=k$), one has ${}^{hs}\!A\cong A$.
  \item If $i':k\to\Omega'$ is a second homomorphism from $k$ into a separable closure of $k$, for every $k$-isomorphism $\sigma:\Omega\xrightarrow{\sim}\Omega'$, the corresponding $A$-homomorphism \textup{(}by~\textup{(ii))} $v_\sigma:{}^{hs}\!A_{(i)}\to{}^{hs}\!A_{(i')}$ is an isomorphism.
\end{enumerate}
\end{proposition}

\begin{proof}
We have already seen above that ${}^{hs}\!A_{(i)}$ is a local ring and that $w$ is a local homomorphism;
the fact that ${}^{hs}\!A_{(i)}$ is Henselian (and hence strictly local) is demonstrated as in \sref{IV.18.6.6}.
The assertions of~(iii), as well as the sufficiency of condition~(iv), also follow from \sref[0\textsubscript{\mdseries III}]{0.10.3.1.3}; the necessity of condition~(iv) follows from the fact that ${}^{hs}\!A_{(i)}$ is a faithfully flat $A$-module \sref[0\textsubscript{\mdseries I}]{0.6.5.2}.

To prove~(ii), remark, with the notations of \sref{IV.18.8.6}, that for every $\lambda\in L$, there exists one and only one local homomorphism $v_\lambda:A_\lambda\to B$, such that the composite $\boldsymbol{k}(A_\lambda)\xrightarrow{\omega_\lambda}\Omega\xrightarrow{\psi}\boldsymbol{k}(B)$ is deduced from $v_\lambda$ by passage to quotients, by virtue of \sref{IV.18.8.4} and the hypothesis that $B$ is strictly local.
The uniqueness of $v_\lambda$ implies moreover that the $v_\lambda$ form an inductive system, whence the existence of the homomorphism $v$ satisfying the conditions of~(ii); its uniqueness follows from \sref{IV.17.4.9}.
The uniqueness of $v$ immediately implies assertion~(vi), by considering the composites $v_\sigma\circ v_{\sigma^{-1}}$ and $v_{\sigma^{-1}}\circ v_\sigma$.
Finally, if $A$ is strictly local, and $A'=B_{\mathfrak{n}}$ is an $A$-algebra local and essentially \'{e}tale, where $B$ is an $A$-algebra \'{e}tale, it follows from \sref{IV.18.8.1} that there exists a $\Spec(A)$-section of $\Spec(B)$ taking the value $\mathfrak{n}$ at the closed point of $\Spec(A)$, hence by \sref{IV.17.4.1} that the homomorphism $A\to A'$ is bijective, which proves~(v).
\end{proof}

\begin{env}[18.8.8.1]
\label{IV.18.8.8.1}
Let $\boldsymbol{C}$ be the category of strictly local rings, with local homomorphisms as morphisms;
for every $B\in\boldsymbol{C}$, denote by $F(B)$ the set of pairs $(u,\psi)$ consisting of a local homomorphism $u:A\to B$ and a $k$-homomorphism $\psi:\Omega\to\boldsymbol{k}(B)$ such that the composite $\boldsymbol{k}(A)=k\xrightarrow{i}\Omega\xrightarrow{\psi}\boldsymbol{k}(B)$ is deduced from $u$ by passage to quotients.
One can say that the object ${}^{hs}\!A_{(i)}$ \emph{represents} the covariant functor $F:\boldsymbol C\to\boldsymbol{\mathit{Ens}}$ \sref[0\textsubscript{\mdseries III}]{0.8.1.11}.

One will remark that by virtue of \sref{IV.18.8.8},~(vi), the strict Henselizations ${}^{hs}\!A_{(i)}$ are all $A$-isomorphic (for the various $k$-homomorphisms $i$ of $k$ into separable closures of $k$), but for given $i$ and $i'$, there is in general an infinity of $A$-isomorphisms ${}^{hs}\!A_{(i)}\xrightarrow{\sim}{}^{hs}\!A_{(i')}$.
To be precise, the group of $A$-automorphisms of ${}^{hs}\!A_{(i)}$ is isomorphic to the Galois group of $\Omega$ over $k$.
\end{env}

\begin{env}[18.8.9]
\label{IV.18.8.9}
Let now $A$ be a \emph{semi-local} ring, $\mathfrak{m}_j$ ($1\leq j\leq r$) its maximal ideals, and for every $j$, consider a homomorphism $i_j:\boldsymbol{k}(A_{\mathfrak{m}_j})\to\Omega_j$ of the residue field into a separable closure.
The \emph{strict Henselization} of $A$ relative to the $i_j$ is the product $\prod_{j=1}^{r}({}^{hs}(A_{\mathfrak{m}_j}))$, denoted ${}^{hs}\!A$ (when there is no risk of confusion).
\oldpage[IV-4]{148}
It is therefore a faithfully flat $A$-module and a semi-local $A$-algebra, whose maximal ideals are the $\mathfrak{m}_j.{}^{hs}\!A$, while its radical is $\mathfrak{r}.{}^{hs}\!A$ (if $\mathfrak{r}$ denotes the radical of $A$).
Finally, the universal property \sref{IV.18.8.8},~(ii) still holds when replacing ``local homomorphism'' by ``semi-local homomorphism'' \sref{IV.18.6.7}, and replacing $\Omega$ by $\prod_j\Omega_j$.
\end{env}

\begin{proposition}[18.8.10]
\label{IV.18.8.10}
Let $A$ be a semi-local ring, $B$ a finite $A$-algebra.
Let $\mathfrak{n}_j$ \textup{($1\leq j\leq r$)} be the maximal ideals of $B$.
Then there is, for every $j$, a maximal ideal $\mathfrak{q}_j$ of $B\otimes_{A}({}^{hs}\!A)$ above $\mathfrak{n}_j$, such that ${}^{hs}\!B$ is isomorphic to the direct product of the local rings $(B\otimes_{A}({}^{hs}\!A))_{\mathfrak{q}_j}$ \textup{($1\leq j\leq r$)}.
\end{proposition}

\begin{proof}
One can evidently restrict to the case where $A$ is local, replacing $A$ by $A_{\mathfrak{m}_j}$, where $\mathfrak{m}_j$ is the maximal ideal of $A$ that is the inverse image of $\mathfrak{n}_j$.
On the other hand, let $C=B_{\mathfrak{n}_j}$;
by virtue of the definition of ${}^{hs}\!B$ \sref{IV.18.8.9}, everything amounts to seeing that ${}^{hs}\!C$ is $C$-isomorphic to $(C\otimes_{A}({}^{hs}\!A))_{\mathfrak{q}}$, where $\mathfrak{q}$ is one of the maximal ideals of $C\otimes_{A}({}^{hs}\!A)$.
Let $k$ and $K$ be the residue fields of $A$ and $C$, $\Omega$ the separable closure of $k$, and let us suppose $K$ and $\Omega$ contained in the same closure of $k$;
as $K$ is a finite extension of $k$, it is immediate that $K\otimes_k\Omega$ is a finite direct product of fields, all isomorphic to the separable closure $K\Omega$ of $K$.
On the other hand, since ${}^{hs}\!A$ is Henselian, $C\otimes_{A}({}^{hs}\!A)$ is a direct product of the local rings $(C\otimes_{A}({}^{hs}\!A))_{\mathfrak{q}_i}$, where $\mathfrak{q}_i$ ($1\leq i\leq s$) ranges over the maximal ideals of $C\otimes_{A}({}^{hs}\!A)$; these local rings are Henselian \sref{IV.18.5.10} and have $K\Omega$ as residue field, hence they are strictly local \sref{IV.18.8.1}.
Let $C'=C\otimes_{A}({}^{hs}\!A)$ and $C'_i=C'_{\mathfrak{q}_i}$, which is therefore a quotient ring of $C'$;
let $f:C\to C'$, $g:{}^{hs}\!A\to C'$, $\vphi_i:C'\to C'_i$ be the canonical maps.
\[
  \xymatrix{
    &&& {}^{hs}\!C \ar@<-0.8ex>[dl]_{w'} \\
    && C'_i \ar@<-0.8ex>[ur]_{w} \\
    C \ar[uurrr]^{v} \ar[r]_{f} & C' \ar[ur] \\
    A \ar[u] \ar[r] & {}^{hs}\!A \ar[u]_{g} \ar[uuurr]_{u}
  }
\]
Let $v:C\to{}^{hs}\!C$ be the canonical local homomorphism;
there exists a unique local homomorphism $u:{}^{hs}\!A\to{}^{hs}\!C$ which, by passage to quotients, gives the canonical injection $\Omega\to K\Omega$ and such that the composite $A\to C\xrightarrow{v}{}^{hs}\!C$ is equal to the composite $A\to{}^{hs}\!A\xrightarrow{u}{}^{hs}\!C$ \sref{IV.18.8.8};
hence there exists a unique local homomorphism $w_0:C'\to{}^{hs}\!C$ such that $w_0\circ g=u$, $w_0\circ f=v$, and since ${}^{hs}\!C$ is a local ring, this homomorphism factorises as $C'\to C'_i\xrightarrow{w}{}^{hs}\!C$ for a well-determined index $i$, $w$ being a local homomorphism.
On the other hand, the ring $C'_i$ being strictly local, there exists a local homomorphism $w':{}^{hs}\!C\to C'_i$ which, by passage to quotients, gives the identity automorphism of $K\Omega$ and which is such that $w'\circ v=\vphi_i\circ f$ \sref{IV.18.8.8}.
One deduces first that $w\circ w'$ is an endomorphism of ${}^{hs}\!C$ which, by passage to quotients, gives the identity automorphism of $K\Omega$, hence
\oldpage[IV-4]{149}
is the identity \sref{IV.18.8.8}.
On the other hand, $w'\circ w\circ\vphi_i\circ f=w'\circ v=\vphi_i\circ f$ and $w'\circ w\circ\vphi_i\circ g=w'\circ u=\vphi_i\circ g$;
whence $w'\circ w\circ\vphi_i=\vphi_i$, which is to say $w'\circ w$ is the identity automorphism of $C'_i$.
\end{proof}

\begin{remark}[18.8.11]
\label{IV.18.8.11}
The proof of \sref{IV.18.8.10} uses the fact that $B$ is \emph{integral} over $A$, but uses the fact that $B$ is a \emph{finite} $A$-algebra only to establish that $K\otimes_k\Omega$ is a finite direct product of fields.
This latter property still holds when $B$ is assumed to be an \emph{integral semi-local} $A$-algebra whose maximal ideals $\mathfrak{n}_j$ have \emph{finite} separable residue-field degrees $[\boldsymbol{k}(\mathfrak{n}_j):k]_s$.
\end{remark}

\begin{proposition}[18.8.12]
\label{IV.18.8.12}
Let $A$ be a semi-local ring.
\begin{enumerate}[label=\emph{\roman*)}]
  \item For ${}^{hs}\!A$ to be reduced \textup{(resp.} normal\textup{)}, it is necessary and sufficient that $A$ be so.
  \item Suppose $A$ Noetherian.
  Then, for every prime ideal $\mathfrak{p}$ of $A$, the ring $({}^{hs}\!A)_{\mathfrak{p}}/\mathfrak{p}.({}^{hs}\!A)_{\mathfrak{p}}$ is a finite direct product of separable algebraic extensions of $\boldsymbol{k}(\mathfrak{p})$ \textup{(which implies that the fibres of the canonical morphism $\Spec({}^{hs}\!A)\to\Spec(A)$ are \emph{geometrically regular} and are \emph{discrete} spaces)}.
\end{enumerate}
\end{proposition}

\begin{corollary}[18.8.13]
\label{IV.18.8.13}
Let $A$ be a local Noetherian ring.
For ${}^{hs}\!A$ to possess one of the following properties:
\begin{enumerate}[label=\emph{\alph*)}]
  \item being a Cohen-Macaulay ring \sref[0]{0.16.5.3};
  \item satisfying the condition $(S_n)$ \sref{IV.5.7.3};
  \item being regular;
  \item satisfying the condition $(R_n)$ \sref{IV.5.8.2};
\end{enumerate}
it is necessary and sufficient that $A$ possess the same property.
\end{corollary}

\begin{corollary}[18.8.14]
\label{IV.18.8.14}
Let $A$ be a semi-local Noetherian ring.
For a prime ideal $\mathfrak{p}$ of ${}^{hs}\!A$ to belong to $\operatorname{Ass}({}^{hs}\!A)$, it is necessary and sufficient that $\mathfrak{p}\cap A$ belong to $\operatorname{Ass}(A)$.
\end{corollary}

\begin{proof}
The proofs are the same as those of \sref{IV.18.6.9},~\sref{IV.18.6.10} and \sref{IV.18.6.11} respectively.
\end{proof}

\begin{proposition}[18.8.15]
\label{IV.18.8.15}
Let $A$ be a local ring;
the following properties are equivalent:
\begin{enumerate}[label=\emph{\alph*)}]
  \item $A$ is geometrically unibranch \textup{(resp.} reduced and geometrically unibranch\textup{)}.
  \item For every $A$-algebra $A'$ that is local and essentially \'{e}tale, $\Spec(A')$ is irreducible \textup{(resp.} integral\textup{)}.
  \item $\Spec({}^{hs}\!A)$ is irreducible \textup{(resp.} integral\textup{)}.
\end{enumerate}
\end{proposition}

\begin{proof}
As in \sref{IV.18.6.12}, one reduces to the case where $A$ is reduced, using the relation ${}^{hs}(A_{\mathrm{red}})=({}^{hs}\!A)\otimes_{A}A_{\mathrm{red}}$ \sref{IV.18.8.11}; the equivalence of \emph{b)} and \emph{c)} is proved as in \sref{IV.18.6.12}.  The same is true of the implication $\textit{a)}\Rightarrow\textit{c)}$, taking account of \sref{IV.18.8.11} and of the definition of a geometrically unibranch integral local ring.
Finally, to prove that \emph{c)} implies \emph{a)}, using the same notations as in \sref{IV.18.6.12}, it suffices to see that ${}^{hs}\!B$ is integral.
But, by \sref{IV.18.8.10}, ${}^{hs}\!B$ is a localization of the semi-local ring $B\otimes_{A}({}^{hs}\!A)$.  By flatness, the latter again identifies with a subring of the field $L$ and is therefore integral; consequently ${}^{hs}\!B$ is integral as well, which completes the proof.
\end{proof}

\begin{corollary}[18.8.16]
\label{IV.18.8.16}
Let $A$ be a local Henselian ring \textup{(resp.} strictly local\textup{)}.
For $A$ to be unibranch \textup{(resp.} geometrically unibranch\textup{)}, it is necessary and sufficient that $\Spec(A)$ be irreducible.
\end{corollary}

\oldpage[IV-4]{150}

\begin{proposition}[18.8.17]
\label{IV.18.8.17}
Let $A$ be a local Noetherian ring.
If $A$ is universally catenary, the same holds for ${}^{hs}\!A$.
\end{proposition}

\begin{proof}
The proof is the same as that of \sref{IV.18.7.5}.
\end{proof}

\begin{proposition}[18.8.18]
\label{IV.18.8.18}
\begin{enumerate}[label=\emph{\roman*)}]
  \item Every filtered inductive limit of strictly local rings (where the transition homomorphisms are local) is a strictly local ring.
  \item The functor $A\leadsto{}^{hs}\!A$ in the category of local rings (where the morphisms are local homomorphisms) commutes with filtered inductive limits.
  \item Every strictly local ring $A$ is an inductive limit of a filtered inductive system of Noetherian strictly local rings (the transition homomorphisms being local).
\end{enumerate}
\end{proposition}

\begin{proof}
The proofs are modelled on those of \sref{IV.18.6.14}.
\end{proof}


\subsection{Formal fibres of Noetherian Henselian rings.}
\label{subsection:IV.18.9}

\begin{theorem}[18.9.1]
\label{IV.18.9.1}
Let $A$ be a local Noetherian Henselian ring, whose formal fibres are geometrically normal.
Then the formal fibres of $A$ are geometrically integral \textup{(hence} geometrically connected\textup{)}.
\end{theorem}

\begin{proof}
The two assertions stated are in fact equivalent, a prescheme that is locally Noetherian, connected and normal being integral.
As every finite integral $A$-algebra is a Henselian local ring \sref{IV.18.5.9} and \sref{IV.18.5.10}, it follows from \sref{IV.7.3.16.2} that one is reduced to proving that if one further supposes $A$ \emph{integral}, then the fibre of the morphism $\Spec(\hat{A})\to\Spec(A)$ at the generic point of $\Spec(A)$ is \emph{integral}, which follows from the following corollary.
\end{proof}

\begin{corollary}[18.9.2]
\label{IV.18.9.2}
Let $A$ be a local Noetherian Henselian integral ring, whose formal fibres are geometrically normal.
Then $\hat{A}$ is integral.
\end{corollary}

\begin{proof}
Indeed, it follows from \sref{IV.18.8.16} that $A$ is unibranch, hence the theorem follows from the hypothesis made on the formal fibres of $A$ and from \sref{IV.7.6.3}.
\end{proof}

\begin{corollary}[18.9.3]
\label{IV.18.9.3}
Under the hypotheses of \sref{IV.18.9.2}, the field of fractions $L$ of $\hat{A}$ is a separable extension of the field of fractions $K$ of $A$, and $K$ is algebraically closed in $L$.
\end{corollary}

\begin{proof}
This follows from the fact that the fibre of the morphism $\Spec(\hat{A})\to\Spec(A)$ at the generic point is geometrically integral and from \sref{IV.4.3.2} and \sref{IV.4.3.5}.
\end{proof}

\begin{corollary}[18.9.4]
\label{IV.18.9.4}
Let $A$ be a local Noetherian ring, Henselian, and whose formal fibres are geometrically normal \textup{(for example a Noetherian Henselian local ring that is} excellent\textup{)}, and let $A'=\hat{A}$.
Let $X$ be an $A$-prescheme, $X'=X\otimes_{A}A'$, $g:X'\to X$ the canonical projection.
Then the map $U\leadsto g^{-1}(U)$ is a bijection from the set of subsets of $X$ that are both open and closed to the set of subsets of $X'$ that are both open and closed.
\end{corollary}

\begin{proof}
As $g$ is a faithfully flat and quasi-compact morphism \sref{IV.2.3.12}, one knows that the topology of $X$ is the \emph{quotient} of that of $X'$ by the equivalence relation defined by $g$.
It therefore amounts to seeing that an open and closed subset $U'$ of $X'$ is \emph{saturated} for this relation.
Now, as the morphism $\Spec(A')\to\Spec(A)$ has its fibres geometrically connected \sref{IV.18.9.1}, the fibres of $g$ are \emph{connected}, whence our assertion immediately follows.

In particular, if $X$ is locally connected (which will be the case if $X$ is locally Noetherian), then, for every connected component $U$ of $X$ (which is open and closed
\oldpage[IV-4]{151}
in $X$), $g^{-1}(U)$ is connected (Bourbaki, \emph{Top. g\'{e}n.}, chap.~I\textsuperscript{er}, 3\textsuperscript{e} \'{e}d., \S~11, no~3, prop.~7).
If one denotes by $\pi_0(X)$ the set of connected components of $X$, the map $\pi_0(X')\to\pi_0(X)$ canonically deduced from $g$ is therefore \emph{bijective}.
\end{proof}

\begin{corollary}[18.9.5]
\label{IV.18.9.5}
Under the hypotheses of \sref{IV.18.9.4}, the functor
\begin{equation}
  Z\leadsto Z\otimes_{A}A'
  \tag{18.9.5.1}
\label{IV.18.9.5.1}
\end{equation}
from the category of \'{e}tale preschemes over $X$, to that of \'{e}tale preschemes over $X'$, is fully faithful.
\end{corollary}

\begin{proof}
Let $Z_1$, $Z_2$ be two \'{e}tale preschemes over $X$, and set $Z'_i=Z_i\otimes_{A}A'$ ($i=1,2$);
it is a matter of proving that every $X'$-morphism $Z'_1\to Z'_2$ comes, by base change, from a unique $X$-morphism $Z_1\to Z_2$.
Suppose first that $Z_2$ is \emph{separated} over $X$;
then, as $\Hom_{X}(Z_1,Z_2)$ identifies with $\Gamma((Z_1\times_{X}Z_2)/Z_1)$, and as $Z_1\times_{X}Z_2$ is \'{e}tale and separated over $Z_1$, $\Hom_{X}(Z_1,Z_2)$ identifies functorially, by virtue of \sref{IV.17.9.3}, with the set of open and closed subsets $U$ of $Z_1\times_{X}Z_2$ such that the restriction to $U$ of the projection $p:Z_1\times_{X}Z_2\to Z_1$ is a surjective and radicial morphism.
The assertion thus follows from \sref{IV.18.9.4} and the fact that $Z'_1\times_{X'}Z'_2=(Z_1\times_{X}Z_2)\otimes_{A}A'$.

Let us pass to the general case.
It will suffice to prove that the graph $\Gamma'$ of an $X'$-morphism $Z'_1\to Z'_2$ which is \emph{open} in $Z'_1\times_{X'}Z'_2$ \sref{IV.17.9.3}, is of the form $\Gamma\otimes_{A}A'$, where $\Gamma$ is induced by an \emph{open} subset of $Z=Z_1\times_{X}Z_2$.
Indeed, the restriction to $\Gamma$ of the projection $p:Z\to Z_1$ is then an isomorphism: after the base change $A\to A'$, it becomes the restriction of $p':Z'\to Z'_1$ to $\Gamma'$, which is an isomorphism, and it suffices to apply \sref{IV.2.7.1},~(viii).
The conclusion then follows from the characterisation of graphs \sref[I]{subsection:I.5.3}.

If $q:Z'=Z\otimes_{A}A'\to Z$ is the canonical projection, it remains to prove that $\Gamma'=q^{-1}(\Gamma)$ for an open subset $\Gamma$ of $Z$.
Since $q$ is faithfully flat and quasi-compact, it suffices \sref{IV.2.3.12} to prove that $\Gamma'=q^{-1}(U)$ for a subset $U\subset Z$.
Let $Z''=Z'\times_{Z}Z'=Z\times_{X}X''$ where $X''=X'\times_{X}X'$, and let $q_1:Z''\to Z'$, $q_2:Z''\to Z'$ be the canonical projections.
Applying \sref{IV.4.5.19.1}, it suffices to show that $q_1^{-1}(\Gamma')=q_2^{-1}(\Gamma')$.
But this is a property which is true if and only if it is true after each base change $\Spec(\boldsymbol{k}(x))\to X$, where $x$ ranges over $X$.
In other words, it suffices to prove the corollary when $X$ is the \emph{spectrum of a field};
but as every $X$-prescheme that is \'{e}tale is automatically \emph{separated} over $X$ \sref{IV.17.6.2},~\emph{c$'$)}, one is reduced to the case considered at the beginning of the proof.
\end{proof}

\begin{remarks}[18.9.6]
\label{IV.18.9.6}
\begin{enumerate}[label=\emph{\roman*)}]
  \item It is possible that, if the residue field of $A$ is of characteristic 0, the functor \sref{IV.18.9.5.1} induces even an \emph{equivalence} of the category of \'{e}tale coverings of $X$ and the category of \'{e}tale coverings of $X'$.
  One can indeed prove this when $A$ is \emph{excellent}, using the resolution of singularities of Hironaka (M.~Artin).
  It is plausible that an analogous statement is still true without restriction on the characteristic of the residue field, provided one restricts to the ``principal Galois'' coverings whose group has order prime to the residual characteristic (cf.~[41]).

  \oldpage[IV-4]{152}
  \item One does not know whether the formal fibres of $A$ are geometrically connected (or even geometrically irreducible) when $A$ is any Henselian local ring.
  It would suffice that, for every Noetherian Henselian local \emph{integral} ring $A$, $\Spec(\hat{A})$ be irreducible.
  One does not know if this is always the case.

  \item One does not know whether, when $A$ is an excellent local ring, its strict Henselization ${}^{hs}\!A$ is excellent.
  To elucidate this question, one can reduce to the case where $A=k[[T_1,\ldots,T_n]]$, $k$ being a field of characteristic $p>0$.
  Taking account of \sref{IV.18.8.17}, the question is whether the formal fibres of ${}^{hs}\!A$ are geometrically regular.
  The answer is affirmative when $n=1$, but is not known for $n=2$.
  One can show that the answer is affirmative whenever, for every scheme $Y_1$ finite over $Y=\Spec(\hat{A})$, the singularities of $Y_1$ can be resolved \sref{IV.7.9.1}.
\end{enumerate}
\end{remarks}

The connectedness assertion made in \sref{IV.18.9.1} generalises in the following way:

\begin{theorem}[18.9.7]
\label{IV.18.9.7}
Let $A$, $B$ be two local Noetherian rings, $\rho:A\to B$ a local homomorphism, $f:\Spec(B)=X\to\Spec(A)=Y$ the corresponding morphism.
Suppose that the following conditions are satisfied:
\begin{enumerate}[label=\emph{\roman*)}]
  \item $f$ is flat and all the fibres $X_y=f^{-1}(y)$ \textup{(}$y\in Y$\textup{)} are geometrically reduced \sref{IV.4.6.2}.
  \item \textup{Either} the ring $A$ is geometrically unibranch \textup{\sref[0]{0.23.2.1}, or else} the ring $A$ is unibranch \sref[0]{0.23.2.1} and $\boldsymbol{k}(B)$ is a primary extension \sref{IV.4.3.1} of $\boldsymbol{k}(A)$.
\end{enumerate}
Then, if $\eta$ is the generic point of $Y$, $X_\eta$ is connected, and for every closed subset $T$ of $X$ such that $f(T)$ is rare in $Y$, $X\setmin T$ is connected.
\end{theorem}

\begin{proof}
The two conclusions stated are in fact equivalent;
indeed, the hypothesis that $f(T)$ is rare implies that $\overline{f(T)}$ does not contain $\eta$.
On the other hand, since $\{\eta\}$ is proconstructible \sref{IV.1.9.6} and $f$ is flat, $X_\eta$ is dense in $X$ \sref{IV.2.3.10}; consequently, if $X_\eta$ is connected, then so is $X\setmin T$, which contains it.
Conversely, as $U$ ranges over the affine open neighbourhoods of $\eta$ in $Y$, $X_\eta$ is the projective limit of the subschemes induced on the open subsets $f^{-1}(U)$ of $X$ \sref{IV.8.1.2},~$a$).
If one supposes that $X\setmin T$ is connected for every closed subset $T$ of $X$ such that $f(T)$ is rare in $Y$, the $f^{-1}(U)$ are connected, and therefore so is $X_\eta$ \sref{IV.8.4.1}.

Let us first show that one can reduce to the case where $A$ is integral and $Y=\Spec(A)$ is geometrically unibranch \sref{IV.6.15.1} (which implies that $A$ is geometrically unibranch, but is not equivalent to this condition \sref{IV.6.15.2}).
It is clear that one can first suppose $Y$ reduced, by considering the morphism $X\times_{Y}Y_{\mathrm{red}}\to Y_{\mathrm{red}}$ obtained from $f$ by base change, which is flat and has the same fibres as $f$.
One can therefore suppose $A$ integral and unibranch.
If $K$ is the field of fractions of $A$, there then exists a \emph{finite} $A$-subalgebra $A''$ of $K$ such that, if $A'$ is the integral closure of $A$, the morphism $\Spec(A')\to\Spec(A'')$ is \emph{radicial} \sref[0]{0.23.2.5}; it follows \sref{IV.6.15.5} that $\Spec(A'')$ is \emph{geometrically unibranch}.
Since $A$ is unibranch, $A'$ is a local ring, and hence so is $A''$.
Let us show that the ring $B''=B\otimes_{A}A''$ is also \emph{local}.
Indeed, $B''$ is a finite $B$-algebra, hence a Noetherian semilocal ring.
If $k''$ is the residue field of $A''$, the maximal ideals of $B''$ are the points of $\Spec(B'')$ lying over the closed point of $\Spec(B)$, hence the points of $\Spec(\boldsymbol{k}(B)\otimes_{\kres(A)}k'')$ \sref[I]{I.3.4.9},
\oldpage[IV-4]{153}
since the closed point of $\Spec(A'')$ is the only point lying over the closed point of $\Spec(A)$.
Now, when $A$ is geometrically unibranch, $k''$ is a radicial extension of $\boldsymbol{k}(A)$, so $\Spec(\boldsymbol{k}(B)\otimes_{\kres(A)}k'')$ is radicial over $\Spec(\boldsymbol{k}(B))$ and consequently consists of a single point.
When $A$ is unibranch and $\boldsymbol{k}(B)$ is a primary extension of $\boldsymbol{k}(A)$, $\Spec(\boldsymbol{k}(B)\otimes_{\kres(A)}k'')$ is irreducible \sref{IV.4.3.2}; since it is a finite discrete space \sref[I]{I.6.4.4}, it again consists of a single point.
Thus $B''$ is \emph{local} in both cases considered.
If $Y''=\Spec(A'')$, $X''=\Spec(B'')=X\times_{Y}Y''$, the morphism $f''=f_{(Y'')}:X''\to Y''$ is flat and has geometrically reduced fibres.
Moreover, if $\eta''$ is the generic point of $Y''$, then $\boldsymbol{k}(\eta'')=\boldsymbol{k}(\eta)=K$ by definition, so the fibres $X_\eta$ and $X''_{\eta''}$ are isomorphic.
It therefore suffices to prove that $X''_{\eta''}$ is connected.
\end{proof}

\begin{remark}[18.9.7.1]
\label{IV.18.9.7.1}
When the ring $A$ is \emph{Japanese}, one can, in the case preceding \sref{IV.18.9.7.2}, take $A''=A'$ by definition \sref[0]{0.23.1.1};
one sees therefore that in this case one would even be able to prove the theorem when $A$ is \emph{integrally closed}.
\end{remark}

\begin{env}[18.9.7.2]
\label{IV.18.9.7.2}
Suppose therefore now that $Y$ is integral and geometrically unibranch.
Note that if $T$ is a closed subset of $X$ such that $f(T)$ is rare in $Y$, $T$ contains no maximal point of $X$ since $f$ is flat \sref{IV.2.3.4}, and is therefore rare in $X$.
Since $X$ is a local scheme, hence connected, to prove that $X\setmin T$ is connected it suffices, by virtue of \sref{IV.15.5.6.1}, to show that for every $x\in X$ such that $f(x)\neq\eta$, $\Spec(\sh{O}_{X,x})\setmin\{x\}$ is \emph{connected}.
Set $y=f(x)$, $A_1=\sh{O}_{Y,y}$, $B_1=\sh{O}_{X,x}$, $Y_1=\Spec(A_1)$, $X_1=\Spec(B_1)$, and let $f_1:X_1\to Y_1$ be the morphism corresponding to the local homomorphism $A_1\to B_1$ induced by $\rho$.
It is clear that $f_1$ is flat, and it follows from \sref{IV.4.6.1} that its fibres are geometrically reduced.
Moreover, $A_1$ is integral and geometrically unibranch \sref{IV.6.15.1}, and $\dim(A_1)\geq 1$.
One is thus reduced to proving the following lemma.
\end{env}

\begin{lemma}[18.9.7.3]
\label{IV.18.9.7.3}
Let $A$, $B$ be two local Noetherian rings, $Y=\Spec(A)$, $X=\Spec(B)$, $\rho:A\to B$ a local homomorphism, $f:X\to Y$ the corresponding morphism.
Suppose that the following conditions are satisfied:
\begin{enumerate}[label=\emph{\roman*)}]
  \item $f$ is flat and all the fibres $X_y$ \textup{(}$y\in Y$\textup{)} are geometrically reduced.
  \item $A$ is integral, geometrically unibranch and $\dim(A)\geq 1$.
\end{enumerate}
Then, if $b$ is the closed point of $X$, $X\setmin\{b\}$ is connected.
\end{lemma}

\begin{proof}
Note first that by virtue of the theorem of Hartshorne \sref{IV.5.10.7}, $X\setmin\{b\}$ is connected if $\operatorname{prof}(B)\geq 2$.
Now one has \sref{IV.6.3.1}
\begin{equation}
  \operatorname{prof}(B)=\operatorname{prof}(A)+\operatorname{prof}(B\otimes_{A}k)
  \label{IV.18.9.7.4}
\tag{18.9.7.4}
\end{equation}
where $k$ denotes the residue field of $A$.
On the other hand, since $A$ is integral and $\dim(A)\geq 1$, one has $\operatorname{prof}(A)\geq 1$, hence $\operatorname{prof}(B)\geq 2$ \emph{unless} $\operatorname{prof}(B\otimes_{A}k)=0$;
moreover, $B\otimes_{A}k$ is reduced by~(i), hence the relation $\operatorname{prof}(B\otimes_{A}k)=0$ means that $B\otimes_{A}k$ is a \emph{field} \sref[0]{0.16.4.7}.
We shall henceforth suppose that this last condition is satisfied.
We shall begin by treating a case where the proof is very simple.

A) \emph{Case where $A$ is integrally closed.}
Then, if $\dim(A)\geq 2$, one has $\operatorname{prof}(A)\geq 2$ \sref[0]{0.16.5.1}, hence $\operatorname{prof}(B)\geq 2$.
One therefore only has to consider the case where $\dim(A)=1$ and where $B\otimes_{A}k$ is a field;
$A$ is then a discrete valuation ring, hence regular, and since $B\otimes_{A}k$
\oldpage[IV-4]{154}
is a field and $f$ is flat, $B$ is regular \sref[0]{0.17.3.3}.
Moreover \sref{IV.6.1.1.1}, one has $\dim(B)=\dim(A)+\dim(B\otimes_{A}k)=\dim(A)=1$, so $B$ is also a discrete valuation ring.
Thus $X\setmin\{b\}$ consists of a single point, which proves the lemma \sref{IV.18.9.7.3} in this case.

One will observe that this proves the statement of the theorem \sref{IV.18.9.7} when one supposes in addition that, in that statement, the ring $A$ is \emph{Japanese}, for by virtue of the remark \sref{IV.18.9.7.1}, one is reduced to the case where $A$ is \emph{integrally closed}, and then, in the preceding reduction \sref{IV.18.9.7.3}, the ring $A_1=\sh{O}_{Y,y}$ is also integrally closed and one can apply the result just proved.

B) \emph{Case where} $\dim(A)\geq 2$.
The result of \sref{IV.18.9.7.3} will follow from the two following lemmas.
\end{proof}

\begin{lemma}[18.9.7.5]
\label{IV.18.9.7.5}
Let $A$ be a semi-local Noetherian ring, reduced, $A'$ its integral closure in its total ring of fractions, $X=\Spec(A)$, $X'=\Spec(A')$;
let $S$ be the set of closed points of $X$ and $U=X\setmin S$.
Suppose that for every point $x'\in X'$ above a point of $S$, $\dim(\sh{O}_{X',x'})\geq 2$;
then $A^{(1)}=\Gamma(U,\sh{O}_{X})$ is integral over $A$ and is an $A$-algebra isomorphic to a subalgebra of $A'$.
\end{lemma}

\begin{proof}
Recall that if the $\mathfrak{p}_i$ ($1\leq i\leq m$) are the minimal prime ideals of $A$, then $A'$ is the direct product of the integral closures $A'_i$ of the integral rings $A_i=A/\mathfrak{p}_i$, so that $X'$ is the sum of the schemes $X'_i=\Spec(A'_i)$ ($1\leq i\leq m$).
If $X_0$ is the sum of the schemes $X_i=\Spec(A_i)$ ($1\leq i\leq m$) and $U_0$ the inverse image of $U$ in $X_0$, the canonical homomorphism $\Gamma(U,\sh{O}_{X})\to\Gamma(U_0,\sh{O}_{X_0})$ is injective since $X$ is reduced and the morphism $X_0\to X$ (and hence also $U_0\to U$) is surjective.
As $\Gamma(U_0,\sh{O}_{X_0})$ is the direct product of the $\Gamma(U_i,\sh{O}_{X_i})$ (where $U_i$ is the inverse image of $U$ in $X_i$), it suffices to prove that for each $i$, $\Gamma(U_i,\sh{O}_{X_i})$ is isomorphic to a sub-$A_i$-algebra of $A'_i$.
In other words, one is reduced to the case where $A$ is a Noetherian integral semilocal ring.
Since $X$ is reduced and the morphism $X'\to X$ is surjective, the homomorphism $\Gamma(U,\sh{O}_{X})\to\Gamma(U',\sh{O}_{X'})$ (where $U'$ is the inverse image of $U$ in $X'$) is injective, and it suffices to verify that the canonical homomorphism $A'=\Gamma(X',\sh{O}_{X'})\to\Gamma(U',\sh{O}_{X'})$ is \emph{bijective}.
Now \sref[I]{I.8.2.1.1}, $\Gamma(U',\sh{O}_{X'})$ is the intersection of the local rings $A'_{\mathfrak{p}'}$, where $\mathfrak{p}'$ ranges over $U'$, and by virtue of the hypothesis, among these local rings are included all those \emph{of height} 1.
But the reasoning of \sref[0]{0.23.2.7} applies also to a Noetherian integral \emph{semilocal} ring, so $A'$ is a semilocal \emph{Krull ring}, and is therefore the intersection of the local rings $A'_{\mathfrak{p}'}$, where $\mathfrak{p}'$ ranges over the set of prime ideals of height 1 of $A'$ (Bourbaki, \emph{Alg. comm.}, chap.~VII, \S~1, no~6, th.~4);
\emph{a fortiori}, $A'$ is the intersection of the $A'_{\mathfrak{p}'}$ for $\mathfrak{p}'\in U'$, which completes the proof of the lemma.
\end{proof}

\begin{remark}[18.9.7.6]
\label{IV.18.9.7.6}
We know \sref[0]{0.23.2.5} that there exists a \emph{finite} sub-$A$-algebra $A''$ of $K$, the field of fractions of $A$, such that the morphism $\Spec(A')\to\Spec(A'')$ is \emph{radicial};
as this morphism is also integral and dominant, hence surjective \sref[II]{II.6.1.10}, it is a \emph{homeomorphism} \sref{IV.2.4.5}.
By virtue of the geometric interpretation of $\dim(\sh{O}_{X',x'})$ \sref{IV.5.1.2}, one has therefore, for every point $x''$ of $X''=\Spec(A'')$ above a point $x\in X$, $\dim(\sh{O}_{X'',x''})=\dim(\sh{O}_{X',x'})\geq 2$, where $x'$ is the unique point of $X'$ above $x$.
\oldpage[IV-4]{155}
This being so, it follows first from \sref[0]{0.16.1.5} and \sref[0]{0.16.1.4.1} that the relation $\dim(\sh{O}_{X',x'})\geq 2$ implies $\dim(\sh{O}_{X,x})\geq 2$.
Conversely, suppose that $\dim(\sh{O}_{X,x})\geq 2$ and that $A$ is \emph{universally catenary} \sref{IV.5.6.2}; then so is $\sh{O}_{X,x}$ \sref{IV.5.6.3}, hence $\dim(\sh{O}_{X'',x''})=\dim(\sh{O}_{X,x})\geq 2$ by virtue of \sref{IV.5.6.10}, and consequently $\dim(\sh{O}_{X',x'})\geq 2$.
The same conclusion holds if $\sh{O}_{X,x}$ is \emph{geometrically unibranch}, for the fibre of the morphism $X'\to X$ at the point $x$ is then reduced to a single point, and $\sh{O}_{X',x'}$ is an \emph{integral} $\sh{O}_{X,x}$-algebra;
it suffices then to apply \sref[0]{0.16.1.5}.
\end{remark}

\begin{lemma}[18.9.7.7]
\label{IV.18.9.7.7}
Let $A$ be a local Noetherian ring, of dimension $\geq 2$, integral and geometrically unibranch, $Y=\Spec(A)$, $y$ the closed point of $Y$.
Let $X$ be a locally Noetherian and connected prescheme and $f:X\to Y$ a flat morphism.
Then, for every closed subset $T\subset f^{-1}(y)$, $X\setmin T$ is connected.
\end{lemma}

\begin{proof}
The same reasoning as at the beginning of \sref{IV.18.9.7.2} reduces to proving that, for every point $x\in f^{-1}(y)$, $\Spec(\sh{O}_{X,x})\setmin\{x\}$ is connected.
One can therefore restrict to the case where $X=\Spec(B)$, $B$ being a local Noetherian ring and the homomorphism $A\to B$ being local;
let $U=Y\setmin\{y\}$, and note that $V=f^{-1}(U)$ is dense in $X$ \sref{IV.2.3.10};
it suffices therefore to prove that $V$ is \emph{connected}.
Keeping the notations of \sref{IV.18.9.7.5}, and setting $A^{(1)}=\Gamma(U,\sh{O}_{Y})$, $Y_1=\Spec(A^{(1)})$, $X_1=X\times_{Y}Y_1=\Spec(B\otimes_{A}A^{(1)})$, we have the commutative diagram
\[
  \xymatrix{
    X \ar[d]_{f} & X_1 \ar[l]_{g_1} \ar[d]^{f_1} \\
    Y & Y_1 \ar[l]^{g}
  }
\]
By virtue of \sref{IV.2.3.1} and of the definition of $A^{(1)}$, one has a canonical isomorphism $\Gamma(V,\sh{O}_{X})\xrightarrow{\sim}B\otimes_{A}A^{(1)}$, and it suffices therefore to prove that this last ring is \emph{local}, since a local ring cannot be a product of two nonzero rings \sref[III]{III.7.8.6.1}.
But we have seen \sref{IV.18.9.7.5} that $A^{(1)}$ is a sub-$A$-algebra of the integral closure $A'$ of $A$, since $\dim(A)\geq 2$.
The hypothesis that $A$ is geometrically unibranch implies that the morphism $\Spec(A')\to\Spec(A)$ is radicial at the point $y$ \sref{IV.6.15.3};
\emph{a fortiori} $g:Y_1\to Y$ is radicial at the point $y$, and consequently $g_1$ is radicial at the closed point $x$ of $X$;
furthermore $g_1$ is an integral morphism, hence the closed points of $X_1$ are above $x$, hence there is only one closed point of $X_1$, which completes the proof of \sref{IV.18.9.7.7}.

It is clear that \sref{IV.18.9.7.7} proves \sref{IV.18.9.7.3} when $\dim(A)\geq 2$ (even without supposing the fibres geometrically reduced).

C) \emph{Case where} $\dim(A)=1$.
One has seen at the beginning of the proof of \sref{IV.18.9.7.3} that one can suppose $\dim(B\otimes_{A}k)=0$, and by flatness \sref{IV.6.1.1.1}, one therefore has
\[
  \dim(B)=\dim(A)+\dim(B\otimes_{A}k)=1.
\]
Since $A$ is integral, $Y$ consists of two points, the generic point $\eta$ and the closed point $a$.
Furthermore, since $f$ is flat, every irreducible component of $X$ dominates $Y$ \sref{IV.2.3.4}, hence $X\setmin\{b\}$, which is the set of maximal points $\xi_i$ of $X$, is equal to
\oldpage[IV-4]{156}
the underlying set of $f^{-1}(\eta)$, and the fibre $f^{-1}(\eta)$ is the sum of the $\Spec(\boldsymbol{k}(\xi_i))$.
The hypothesis on the $X_y$ and the fact that these fibres are of dimension 0, show that they are \emph{geometrically regular} \sref{IV.6.7.6}, and \emph{a fortiori} geometrically normal, in other words $f$ is a \emph{normal morphism}.
But since $A$ is integral and geometrically unibranch, it follows then from \sref{IV.6.15.10} that $B$ is also integral and geometrically unibranch, hence $f^{-1}(\eta)=X\setmin\{b\}$ is reduced to a single point, and \emph{a fortiori} connected;
this concludes the proof of \sref{IV.18.9.7.3} and of \sref{IV.18.9.7}.
\end{proof}

\begin{remark}[18.9.7.8]
\label{IV.18.9.7.8}
In case C) of the proof of \sref{IV.18.9.7.3}, one can avoid making appeal to the delicate result \sref{IV.6.15.10} by reasoning as follows:
since $A$ is integral and of dimension 1, it follows from the theorem of Krull-Akizuki (Bourbaki, \emph{Alg. comm.}, chap.~VII, \S~2, no~5, prop.~5) that its integral closure $A'$ is a \emph{Noetherian} ring.
The same reasoning as at the beginning of the proof of \sref{IV.18.9.7} shows that $B'=B\otimes_{A}A'$ is a local ring;
furthermore, by flatness, $B'$ is contained in the total ring of fractions $R$ of the reduced $B$ \sref{IV.3.3.5}.
Now, the theorem of Krull-Akizuki (Bourbaki, \emph{loc. cit.}) also applies to a local Noetherian \emph{reduced} ring of dimension 1, and shows that for every such ring $B$, everything contained between $B$ and its total ring of fractions is \emph{Noetherian}.
The ring $B'$ being a local Noetherian ring and the morphism $\Spec(B')\to\Spec(A')$ being normal, $B'$ is an integral ring and integrally closed \sref{IV.6.5.4}, and one concludes as at the beginning of the proof of \sref{IV.18.9.7}.
\end{remark}

\begin{corollary}[18.9.8]
\label{IV.18.9.8}
With the notations of \sref{IV.18.9.7}, suppose that condition~\textup{(i)} of \sref{IV.18.9.7} and \emph{one} of the following two conditions hold:
\begin{enumerate}
  \item[\rm\textup{(ii$'$)}] $A$ is a strictly local ring \sref{IV.18.8.2}.
  \item[\rm\textup{(ii$''$)}] $A$ is a Henselian local ring and $\boldsymbol{k}(B)$ is a primary extension of $\boldsymbol{k}(A)$.
\end{enumerate}
Then \emph{all} the fibres $X_y$ \textup{($y\in Y$)} are geometrically connected.
\end{corollary}

\begin{proof}
Let $\mathfrak{p}$ be the prime ideal of $A$ corresponding to the point $y$, and set $A'=A/\mathfrak{p}$ and $B'=B\otimes_{A}A'=B/\mathfrak{p}B$, which are evidently local Noetherian rings;
it is clear that the morphism $f':\Spec(B')\to\Spec(A')$ satisfies condition~(i) of \sref{IV.18.9.7};
and as $y$ is the generic point of $\Spec(A')$ and $\boldsymbol{k}(A')=\boldsymbol{k}(A)$, $\boldsymbol{k}(B')=\boldsymbol{k}(B)$, one sees that it suffices to verify that, in the case~(ii$'$) (resp.~(ii$''$)), $A'$ is geometrically unibranch (resp. unibranch).
But $A'$ is integral, and in the case~(ii$'$) (resp.~(ii$''$)) it is strictly local (resp. Henselian) by virtue of \sref{IV.18.5.10}.
One therefore concludes that $A'$ is geometrically unibranch (resp. unibranch) by virtue of \sref{IV.18.8.15} (resp. \sref{IV.18.6.12}).
\end{proof}

\begin{corollary}[18.9.9]
\label{IV.18.9.9}
Let $A$ be a local Noetherian ring, Henselian, universally Japanese \textup{(i.e.~\sref{IV.7.6.4} and \sref{IV.7.7.2}, whose \emph{formal fibres} \sref{IV.7.3.13} are geometrically reduced)};
then all the formal fibres of $A$ are geometrically connected.
\end{corollary}

\begin{proof}
It suffices to apply \sref{IV.18.9.8} in the case where $B=\hat{A}$, since $B$ is a flat $A$-module and $\boldsymbol{k}(B)=\boldsymbol{k}(A)$.
\end{proof}

\begin{remark}[18.9.10]
\label{IV.18.9.10}
The proofs of \sref{IV.18.9.4} and \sref{IV.18.9.5} use only the fact that the formal fibres of $A$ are geometrically connected.
The conclusions
\oldpage[IV-4]{157}
of these two propositions are therefore still valid when one supposes the local ring $A$ Noetherian, Henselian and universally Japanese.
\end{remark}

\begin{corollary}[18.9.11]
\label{IV.18.9.11}
Let $X$, $Y$ be two locally Noetherian preschemes, $Y$ being supposed geometrically unibranch and $X$ connected.
Let $f:X\to Y$ be a reduced morphism \textup{(i.e. flat and with geometrically reduced fibres; \sref{IV.6.8.1})}.
Then, for every closed subset $T$ of $X$ such that $f(T)$ is rare in $Y$, $X\setmin T$ is connected.
\end{corollary}

\begin{proof}
As $f$ is flat and $f(T)$, being rare, cannot contain a maximal point of $Y$, $T$ cannot contain a maximal point of $X$ \sref{IV.2.3.4}, and hence is rare in $X$.
Utilising \sref{IV.15.5.6.1}, it suffices to show that, for every $x\in T$, $\Spec(\sh{O}_{X,x})\setmin\{x\}$ is connected.
Set $y=f(x)$, which is not a maximal point of $Y$.
Since the ring $\sh{O}_{Y,y}$ is geometrically unibranch, one can apply \sref{IV.18.9.7} to the morphism
\[
  f_1:\Spec(\sh{O}_{X,x})\to\Spec(\sh{O}_{Y,y}),
\]
which satisfies conditions~(i) and~(ii) of this theorem, and to the closed subset $\{x\}$ of $\Spec(\sh{O}_{X,x})$, since $f_1(\{x\})=\{y\}$ is rare in $\Spec(\sh{O}_{Y,y})$.
\end{proof}
\subsection{\'Etale preschemes over a geometrically unibranch or normal prescheme}
\label{subsection:IV.18.10}

\begin{theorem}[18.10.1]
\label{IV.18.10.1}
Let $f:X\to Y$ be a morphism locally of finite presentation, $x$ a point of $X$, $y=f(x)$.
Suppose that $Y$ is integral and geometrically unibranch at the point $y$.
Then, for $f$ to be \'etale at the point $x$, it is necessary and sufficient that the following two conditions be satisfied:
\begin{enumerate}[label=\emph{\roman*)}]
  \item $f$ is unramified at the point $x$;
  \item the homomorphism $\sh{O}_{Y,y}\to\sh{O}_{X,x}$ is injective.
\end{enumerate}

Moreover, if this is the case, $X$ is integral and geometrically unibranch at the point $x$.

Finally, if $\sh{O}_{Y,y}$ is Noetherian, one can, in the preceding criterion, replace the condition \emph{(ii)} by the condition

\textup{(ii \emph{bis})} $\dim\sh{O}_{Y,y}\leq\dim\sh{O}_{X,x}$.
\end{theorem}

\begin{proof}
The fact that, if $f$ is \'etale, $X$ is integral and geometrically unibranch at the point $x$ is a particular case of \sref{IV.17.5.7}.
It is clear on the other hand that the conditions (i) and (ii) are necessary for $f$ to be \'etale at the point $x$, since $\sh{O}_{X,x}$ is then an $\sh{O}_{Y,y}$-module faithfully flat \sref{IV.17.6.1}.
Let us prove that these conditions are sufficient.

Set $A=\sh{O}_{Y,y}$, $A'={}^{hs}A$ \sref{IV.18.8.7}, $Y'=\operatorname{Spec}(A')$, $X'=X\times_{Y}Y'$, $f'=f_{(Y')}:X'\to Y'$.
Let $y'$ be the closed point of $Y'$, $x'$ a point of $X'$ above $x$ and $y'$; since the morphism $g:Y'\to Y$ is \emph{flat} \sref{IV.18.8.8}, it amounts to showing that $f'$ is \'etale at the point $x'$ \sref{IV.17.7.1},~(ii).
The hypothesis (ii) implies that the morphism $\operatorname{Spec}(\sh{O}_{X,x})\to\operatorname{Spec}(\sh{O}_{Y,y})$ is dominant \sref[I]{I.1.2.7}; as $\sh{O}_{Y,y}$ is integral, $\operatorname{Spec}(\sh{O}_{Y,y})$ has a single generic point $y_1$, and there is a generization $x_1$ of $x$ in $X$ above $y_1$.
The projection morphism $X'\to X$ being flat, there exists a generization $x_1'$ of $x'$ above $x_1$ \sref{IV.2.3.4}; set $y_1'=f'(x_1')$. We then have $g(y_1')=y_1$.
Now, the hypothesis on $y$ implies that $A'$ is \emph{integral} \sref{IV.18.8.15}, hence $Y'$ has a single generic point $\eta$, necessarily above $y_1$ since $g$
\oldpage[IV-4]{158}
is flat \sref{IV.2.3.4}; on the other hand, $\eta$ is also the generic point of $g^{-1}(y_1)$ \sref[0\textsubscript{\mdseries I}]{0.2.1.8}, and one knows on the other hand that the fibres of $g$ are \emph{discrete} spaces \sref{IV.18.8.12}, hence $g^{-1}(y_1)=\{\eta\}$, and consequently $y_1'=\eta$.
The morphism $\operatorname{Spec}(\sh{O}_{X',x'})\to\operatorname{Spec}(\sh{O}_{Y',y'})$, restriction of $f'$, is therefore dominant, and since $\sh{O}_{Y',y'}=A'$ is integral, the corresponding homomorphism $\sh{O}_{Y',y'}\to\sh{O}_{X',x'}$ is injective \sref[I]{I.1.2.7}.
On the other hand, $f'$ is unramified at the point $x'$; hence \sref{IV.18.8.3},~$d$), there exists a neighbourhood $U$ of $x'$ in $X'$ such that $f'|U$ is a \emph{closed immersion} of $U$ into $Y'$.
\emph{A fortiori} the homomorphism $\sh{O}_{Y',y'}\to\sh{O}_{X',x'}$ is surjective, and since it is injective, it is therefore \emph{bijective}.
But then $\sh{O}_{X',x'}$ is a \emph{flat} $\sh{O}_{Y',y'}$-module, and $f'$ is therefore \'etale at the point $x'$ \sref{IV.17.6.1}.

Finally, when $\sh{O}_{Y,y}$ is Noetherian, we know that if $f$ is \'etale at the point $x$, we have $\dim\sh{O}_{Y,y}=\dim\sh{O}_{X,x}$ \sref{IV.17.6.4}.
Conversely, suppose the conditions (i) and (ii~bis) satisfied, and let us show that they imply (ii).
Let $\mathfrak{J}$ be the kernel of the canonical homomorphism $\sh{O}_{Y,y}\to\sh{O}_{X,x}$; restricting $Y$ as needed to a neighbourhood of $y$, we may suppose that $\mathfrak{J}=\sh{J}_y$, where $\sh{J}$ is an ideal of $\sh{O}_Y$.
If $Y_1$ is the closed subscheme of $Y$ defined by $\sh{J}$, one can, restricting further $X$ and $Y$ to neighbourhoods of $x$ and $y$ respectively, suppose that $f$ factorises as $X\xrightarrow{f_1}Y_1\to Y$ \sref[I]{I.6.5.1}.
It is clear that $f_1$ is still unramified at the point $x$, hence quasi-finite at this point \sref{IV.17.4.1}; the proof of \sref{IV.5.4.1},~(i) then shows that $\dim(\sh{O}_{X,x})\leq\dim(\sh{O}_{Y,y}/\mathfrak{J})$.
But if one had $\mathfrak{J}\neq 0$, one would conclude that $\dim(\sh{O}_{Y,y}/\mathfrak{J})<\dim(\sh{O}_{Y,y})$ since $Y$ is integral \sref[0]{0.16.1.2.2}; one would then have $\dim\sh{O}_{X,x}<\dim\sh{O}_{Y,y}$, contradicting the hypothesis, which proves (ii).
\end{proof}

\begin{remarks}[18.10.2]
\label{IV.18.10.2}
\medskip\noindent
\begin{enumerate}[label=(\roman*)]
  \item When $A=\sh{O}_{Y,y}$ is Noetherian, and one knows that its completion $\hat{A}$ is integral (for example if $A$ is regular), one can, in the preceding proof, replace $A'$ by $\hat{A}$.
  \item When $Y$ is locally Noetherian, one can give a faster proof of \sref{IV.18.10.1}, without using strict Henselization, but invoking the rather delicate results of \textsection\textsection~14 and 15.
As $f$ is quasi-finite at the point $x$ by hypothesis \sref{IV.17.4.1}, one can, replacing $X$ by an open neighbourhood of $x$, suppose that $f^{-1}(y)=\{x\}$.
The hypothesis (ii) implies, as we have seen, the existence of an irreducible component $Z$ of $X$ containing $x$ and dominating the unique irreducible component $Y_0$ of $Y$ containing $y$.
The morphism $f|Z$ being quasi-finite at the point $x$, hence equidimensional at this point \sref{IV.13.2.2}, it follows from the hypothesis on $Y$ and from the Chevalley criterion \sref{IV.14.4.4} that $g=f|Z$ is universally open at the point $x$.
Furthermore, since $f$ (and hence also $g$) is unramified at the point $x$, $g^{-1}(y)$ is geometrically reduced over $\boldsymbol{k}(y)$ \sref{IV.17.4.1}; as $\sh{O}_{Y,y}$ is integral, one concludes from \sref{IV.15.2.3} that $g$ is \emph{flat} at the point $x$, hence $f$ is \emph{\'etale} at this point \sref{IV.18.4.9}.
  \item With the same notations and hypotheses on $Y$ as in \sref{IV.18.10.1}, suppose only that $f$ is \emph{locally of finite type} (and not necessarily locally of finite presentation).
Then, for $\sh{O}_{X,x}$ to be an $\sh{O}_{Y,y}$-algebra \emph{essentially \'etale} \sref{IV.18.6.1}, it is necessary and sufficient that $f$ satisfy the two following conditions:
\oldpage[IV-4]{159}
\begin{enumerate}[label=\arabic*\textsuperscript{o}]
  \item \itshape $f$ is formally unramified at the point $x$;
  \item the homomorphism $\sh{O}_{Y,y}\to\sh{O}_{X,x}$ is injective.
\end{enumerate}

There is nothing more than proving the sufficiency of these conditions.
Taking into account \sref{IV.18.4.12}, it suffices to show that $f$ is \emph{flat} at the point $x$.
Now, by going back to the proof of \sref{IV.18.10.1} and the notations used in this proof, it suffices to show that $f'$ is flat at the point $x'$ \sref{IV.2.5.1}; by \sref{IV.17.1.3} furthermore $f'$ is formally unramified at the point $x'$.
One can therefore restrict to carrying out the proof when $Y=\operatorname{Spec}(A)$, where $A$ is a \emph{strictly local} ring and $y=f(x)$ the closed point of $Y$.
Using then \sref{IV.18.8.3} we see that $\sh{O}_{Y,y}\to\sh{O}_{X,x}$ is surjective, hence \emph{bijective} by virtue of the hypothesis, and this suffices to show the flatness of $f$ at the point $x$.

Suppose furthermore that $Y$ is \emph{locally integral} at the point $y$ \sref[I]{I.2.1.8}.
Then the conditions $1^{\circ}$ and $2^{\circ}$ above (together with the fact that $Y$ is geometrically unibranch at the point $y$ and $f$ locally of finite type) already imply that $f$ is \emph{\'etale at the point $x$}.
Indeed, this follows from what precedes and from \sref{IV.18.4.13}.
\end{enumerate}
\end{remarks}

\begin{corollary}[18.10.3]
\label{IV.18.10.3}
Let $Y$ be an integral and geometrically unibranch prescheme, $\eta$ its generic point, $X$ a connected prescheme, and $f:X\to Y$ a morphism locally of finite type such that $f^{-1}(\eta)$ is non-empty.
Then, if $f$ is formally unramified, $f$ is \'etale, and $X$ is integral and geometrically unibranch.
\end{corollary}

\begin{proof}
It follows from the hypothesis and from \sref{IV.18.4.13} that $f$ is \'etale at all the points where it is \emph{flat}, and in particular at the points of $f^{-1}(\eta)$, since $\sh{O}_{Y,\eta}=\boldsymbol{k}(\eta)$ is a field.
The open set $U$ of points of $X$ where $f$ is \'etale is therefore \emph{non-empty}.
Let us show on the other hand that for all $x\in\overline{U}$, the homomorphism $\sh{O}_{Y,f(x)}\to\sh{O}_{X,x}$ is \emph{injective}.
Indeed, let $s$ be a section of $\sh{O}_Y$ over an open neighbourhood of $f(x)$, which one may always suppose to be $Y$ itself (the question being local on $X$ and $Y$), and suppose that its image $t\in\Gamma(Y,f_*(\sh{O}_X))=\Gamma(X,\sh{O}_X)$ has zero germ at the point $x$, hence vanishes in a neighbourhood $V$ of $x$; then the restriction of $t$ to the non-empty open set $U\cap V$ is zero. But the restriction $f|(U\cap V)$ is \'etale, hence $W=f(U\cap V)$ is open in $Y$ and consequently contains the generic point $\eta$. Since the homomorphism $\sh{O}_{Y,\eta}\to\sh{O}_{X,z}$ is injective for every point $z\in f^{-1}(\eta)$, the hypothesis that $t|(U\cap V)=0$ implies that $s_\eta=0$; and as $Y$ is integral, this implies $s=0$, proving our assertion.
Applying now \sref{IV.18.10.2},~(iii), we conclude that $f$ is \emph{flat} at the point $x$, in other words $x\in U$, hence $U$ is both open and closed in $X$; as it is non-empty and $X$ is connected, we have $U=X$.

Note now that since $f$ is \'etale, hence locally quasi-finite and $Y$ is irreducible, the maximal points of $X$ belong to $f^{-1}(\eta)$ \sref{IV.2.3.4} and for all $x\in X$, there is a neighbourhood of $x$ that contains only a finite number of these maximal points; in other words, the set of irreducible components of $X$ is locally finite.
But by \sref{IV.18.10.1}, $X$ is integral and geometrically unibranch at every point, hence every point of $X$ belongs to only one irreducible component; and since the irreducible components form a locally finite family, they are open (and closed) in $X$. Since $X$ is connected, it is integral.
\end{proof}

\oldpage[IV-4]{160}
\begin{remark}[18.10.3.1]
\label{IV.18.10.3.1}
If $f:X\to Y$ is locally of finite presentation, \emph{dominant and quasi-compact}, the fibre $f^{-1}(\eta)$ is non-empty by virtue of \sref{IV.1.1.5}.
On the other hand, if one does not suppose $f$ quasi-compact, \emph{$f$ can be unramified and dominant without being \'etale} (still assuming that the hypotheses of \sref{IV.18.10.3} are satisfied for $X$ and $Y$).
This is what the following example shows: one takes for $Y$ the affine plane $\operatorname{Spec}(\mathbf{C}[T,U])$; one considers on the other hand a family $(X_j)_{j\in\mathbf{Z}}$ of preschemes isomorphic to the affine line $\operatorname{Spec}(\mathbf{C}[T])$, and one forms the prescheme obtained by ``gluing together'' $X_{2j}$ and $X_{2j+1}$ at the point $\setmin 1$ and $X_{2j+1}$ and $X_{2j+2}$ at the point $+1$.
Let us define on the other hand $f:X\to Y$ as being equal on $X_{2j}$ to the closed immersion whose image is the line of equation $y=2j$, transforming the point $\setmin 1$ to $(2j\setmin 1,2j)$ and the point $1$ to $(2j+1,2j)$; on $X_{2j+1}$, $f_j$ is the closed immersion whose image is the line of equation $x=2j+1$, transforming the point $\setmin 1$ to $(2j+1,2j)$ and the point $1$ to $(2j+1,2j+2)$.
It is clear that $f$ is a local immersion (hence unramified) and dominant, but the fibre over the generic point of $Y$ is empty and $f$ is not \'etale.
\end{remark}

\begin{corollary}[18.10.4]
\label{IV.18.10.4}
Let $g:Y\to S$ and $h:X\to S$ be two morphisms locally of finite presentation, $f:X\to Y$ an $S$-morphism, $x$ a point of $X$, $y=f(x)$, and $s=g(y)=h(x)$.
Suppose that $h$ is \textup{flat} at the point $x$ and $g^{-1}(s)$ \textup{normal} at the point $y$.
Then, for $f$ to be \'etale at the point $x$, it is necessary and sufficient that $f$ be unramified at the point $x$ and that $\dim_xh^{-1}(s)=\dim_yg^{-1}(s)$.
\end{corollary}

\begin{proof}
The conditions are necessary because if $f$ is \'etale at the point $x$, the same is true of $f_s:h^{-1}(s)\to g^{-1}(s)$, and the conclusion then follows from the fact that $f_s$ is flat and quasi-finite at the point $x$ \sref{IV.6.1.2}.
Conversely, if the conditions of the statement are satisfied, in order to see that $f$ is \'etale at the point $x$, it suffices, by virtue of \sref{IV.17.8.3}, to see that $f_s$ is so.
One can therefore restrict to the case where $S$ is the spectrum of a field.
It then follows from the hypothesis $\dim_xX=\dim_yY$ and from the fact that $f$ is unramified at the point $x$ that one has, by \sref{IV.5.2.3} and \sref{IV.17.4.1}, $\dim\sh{O}_{X,x}=\dim\sh{O}_{Y,y}$. Since $f$ is quasi-finite at the point $x$ and $\sh{O}_{Y,y}$ is integral, we deduce from \sref[0\textsubscript{\mdseries I}]{0.7.4.4} and \sref[0]{0.16.3.10} that the homomorphism $\sh{O}_{Y,y}\to\sh{O}_{X,x}$ is injective; it then suffices to apply \sref{IV.18.10.1}.
\end{proof}

\begin{corollary}[18.10.5]
\label{IV.18.10.5}
With the notations of \sref{IV.18.10.4}, suppose that $h$ is flat and that $g^{-1}(s)$ is normal for all $s\in S$ \textup{(which will be the case for example if $g$ is \emph{smooth})}.
For $f$ to be an open immersion, it is necessary and sufficient that $f$ be a monomorphism and that, for all $x\in X$, $\dim_xh^{-1}(s)=\dim_yg^{-1}(s)$, where $y=f(x)$ and $s=g(y)=h(x)$.
\end{corollary}

\begin{proof}
It remains only to prove the sufficiency of the conditions. It follows from \sref{IV.17.1.3} that a monomorphism locally of finite presentation is unramified; the hypotheses therefore imply that $f$ is \'etale by \sref{IV.18.10.4}, and one concludes by \sref{IV.17.9.1}.
\end{proof}

\begin{lemma}[18.10.6]
\label{IV.18.10.6}
Let $f:X\to Y$ be a flat and locally quasi-finite morphism \textup{($\mathbf{Err}_{\mathrm{III}}$, 20)}.
\begin{enumerate}[label=(\roman*)]
  \item[\textup{(i)}] The set $\mathrm{Max}(X)$ of maximal points of $X$ is in canonical bijection with the set $\coprod_{y\in\mathrm{Max}(Y)}f^{-1}(y)$.
  \item[\textup{(ii)}] If $Y$ is irreducible with generic point $\eta$, the set of irreducible components of $X$ is in canonical bijection with $f^{-1}(\eta)$, and is in particular finite if and only if $f^{-1}(\eta)$ is finite.
\oldpage[IV-4]{161}
  \item[\textup{(iii)}] If the set of irreducible components of $Y$ is locally finite, so is the set of irreducible components of $X$, and in particular $X$ is locally connected.
\end{enumerate}
\end{lemma}

\begin{proof}
The assertion (i) follows immediately from \sref{IV.2.3.4} and \sref[0\textsubscript{\mdseries I}]{0.2.1.6} and from the fact that the fibres $f^{-1}(y)$ are discrete sets.
It is clear that (ii) is a particular case of (i).
Finally, to prove (iii), one can, by virtue of (i), restrict to the case where $f$ is quasi-finite and where $Y$ is irreducible \sref[0\textsubscript{\mdseries I}]{0.2.1.6} and the conclusion follows from (ii) and from the fact that the fibres of $f$ are finite sets \sref[II]{II.6.2.2}.
\end{proof}

\begin{proposition}[18.10.7]
\label{IV.18.10.7}
Let $Y$ be a geometrically unibranch and irreducible prescheme \textup{(resp.\ }integral\textup{)}, $f:X\to Y$ an \'etale morphism.
Then $X$ is isomorphic to a disjoint union of irreducible \textup{(resp.\ }integral\textup{)} preschemes $X_\lambda$, where $\lambda$ ranges over the fibre $f^{-1}(\eta)$ of the generic point $\eta$ of $Y$.
The $X_\lambda$ are geometrically unibranch; if $Y$ is normal, the $X_\lambda$ are normal.
\end{proposition}

\begin{proof}
Indeed, $X$ is geometrically unibranch \sref{IV.17.5.7}, hence every point of $X$ belongs to only one irreducible component. Moreover, by \sref{IV.18.10.6} and \sref{IV.17.6.1}, the family of irreducible components of $X$ is locally finite; hence \sref[0\textsubscript{\mdseries I}]{0.2.1.6} these are the connected components $X_\lambda$ of $X$ ($\lambda\in f^{-1}(\eta)$), and $X$ is their disjoint union. By \sref{IV.17.5.7}, if $Y$ is integral (resp.\ normal), the $X_\lambda$ are integral (resp.\ normal).
\end{proof}

\begin{proposition}[18.10.8]
\label{IV.18.10.8}
Let $Y$ be a normal and integral prescheme, of generic point $\eta$, $K=\boldsymbol{k}(\eta)$ its field of rational functions, $f:X\to Y$ an \'etale morphism.
Suppose that $f^{-1}(\eta)$ is finite and non-empty, so that the $K$-scheme $f^{-1}(\eta)$ is the spectrum of a finite separable $K$-algebra $L$, which is a direct product of finitely many fields $K_i$ \textup{($1\leq i\leq n$)}, finite separable extensions of $K$ \sref{IV.17.6.2}.
Let $Y'$ be the integral closure of $Y$ in $L$, isomorphic to the sum of the integral closures $Y_i$ of $Y$ in $K_i$ \sref[II]{II.6.3.6}.
Then the morphism $f$ factorises uniquely as $X\xrightarrow{f'}Y'\xrightarrow{g}Y$, such that the restriction $f_\eta':f^{-1}(\eta)\to g^{-1}(\eta)$ is the canonical isomorphism.
Moreover $f'$ is a local isomorphism; for $f'$ to be an open immersion, it is necessary and sufficient that $f$ be separated.
\end{proposition}

\begin{proof}
By virtue of \sref{IV.18.10.7}, one can restrict to the case where $X$ is integral and normal, with $L$ as its field of rational functions and $f$ dominant.
The existence and uniqueness of the factorisation considered for $f$ follow from \sref[II]{II.6.3.9}.
As $f'$ is birational and locally quasi-finite, the final assertions follow from \sref{IV.8.12.10}.
\end{proof}

\begin{corollary}[18.10.9]
\label{IV.18.10.9}
Under the hypotheses of \sref{IV.18.10.8}, for $f$ to be a finite morphism \textup{(in other words, for $X$ to be an \emph{\'etale covering} of $Y$)}, it is necessary and sufficient that $X$ be isomorphic to the integral closure $Y'$ of $Y$ in $L$.
\end{corollary}

\begin{proof}
If $f$ is finite, the canonical injection $j:X\to Y'$ (which is an open immersion) is also a finite morphism \sref[II]{II.6.1.5},~(v), hence a closed morphism \sref[II]{II.6.1.10}, and as $j(X)$ is dense in $Y'$ by \sref{IV.18.10.8}, we have $j(X)=Y'$.
Conversely, as $Y$ is normal and $L$ is a direct product of finite separable extensions of $K$, $Y'$ is finite over $Y$ (Bourbaki, \emph{Alg.\ comm.}, chap.~V, \textsection~1, n\textsuperscript{o}~6, cor.~1 of prop.~18), whence the corollary.
\end{proof}

\begin{env}[18.10.10]
\label{IV.18.10.10}
Let $Y$ be a normal and integral prescheme, $K=R(Y)$ its field of rational functions.
We say that a finite-dimensional $K$-algebra $L$ is \emph{unramified over} $Y$ if: $1^{\circ}$ $L$ is a separable $K$-algebra, hence a direct product of finite separable extensions $K_i$ of $K$ ($1\leq i\leq n$); $2^{\circ}$ the integral closure $Y'$ of $Y$ in $L$ (the disjoint union of the integral closures of $Y$ in the $K_i$) is \emph{unramified over} $Y$ (which, by \sref{IV.18.10.3},
\oldpage[IV-4]{162}
is equivalent to saying that $Y'$ is \emph{\'etale over} $Y$).
When $Y=\operatorname{Spec}(A)$, where $A$ is an integral, integrally closed ring and $K$ its field of fractions, one also says (by a slight abuse of language, and when no confusion with the terminology of \sref{IV.17.3.2},~(ii) is to be feared) that $L$ is \emph{unramified over} $A$ instead of ``unramified over $Y$''.
\end{env}

\begin{remark}[18.10.11]
\label{IV.18.10.11}
Instead of saying that $L$ is unramified over $Y$, some authors also say that $L$ is ``unramified over $K$''; we shall not follow this potentially confusing usage.
\end{remark}

One can then express \sref{IV.18.10.9} in the following form:

\begin{corollary}[18.10.12]
\label{IV.18.10.12}
Let $Y$ be a normal and integral prescheme, $K=R(Y)$ its field of rational functions.
Then the functor $X\leadsto R(X)$ is an equivalence between the category of \'etale coverings of $Y$ and the category of finite \'etale $K$-algebras \textup{unramified} over $Y$.
A quasi-inverse functor sends each finite \'etale $K$-algebra $L$ unramified over $Y$ to the integral closure $Y'$ of $Y$ in $L$.
\end{corollary}

\begin{proposition}[18.10.13]
\label{IV.18.10.13}
Let $Y$ be a normal and integral prescheme, $K=R(Y)$ its field of rational functions.
\begin{enumerate}[label=\textup{(\roman*)}]
  \item The field $K$ is a $K$-algebra unramified over $Y$.
  \item Let $L$ be a finite extension of $K$ unramified over $Y$, and let $Z$ be the integral closure of $Y$ in $L$.
Then, if $M$ is an $L$-algebra unramified over $Z$, $M$ is also a $K$-algebra unramified over $Y$ \textup{(``transitivity'' of unramifiedness)}.
  \item Let $Y'$ be a normal and integral prescheme, $K'$ its field of rational functions, $g:Y'\to Y$ a dominant morphism.
If $L$ is a $K$-algebra unramified over $Y$, then $L\otimes_{K}K'$ is a $K'$-algebra unramified over $Y'$ \textup{(``base-change'' property)}.
Furthermore if $Y=\operatorname{Spec}(A)$ and $Y'=\operatorname{Spec}(A')$ are affine, and if $C$ is the integral closure of $A$ in $L$, then $C'=C\otimes_{A}A'$ is the integral closure of $A'$ in $L\otimes_{K}K'$.
\end{enumerate}
\end{proposition}

\begin{proof}
The assertion (i) is immediate, $Y$ being its own normalisation in $K$.
To prove (ii), let $Z'$ be the integral closure of $Z$ in $M$, which is by hypothesis an \'etale covering of $Z$.
As $Z$ is \'etale and separated over $Y$, so is $Z'$ \sref{IV.17.3.3}, and it is clearly the integral closure of $Y$ in $M$.
As $M$ is a finite separable $K$-algebra, it follows from \sref{IV.18.10.10} that $M$ is unramified over $Y$.
Finally, to prove (iii), observe that if $Z$ is the integral closure of $Y$ in $L$, then $Z\times_{Y}Y'$ is \'etale and separated over $Y'$ \sref{IV.17.3.3}, finite over $Y'$ since $Z$ is finite over $Y$, and has $L\otimes_{K}K'$ as its ring of rational functions \sref[I]{I.3.4.9}. Since $Y'$ is normal, so is $Z\times_{Y}Y'$ \sref{IV.17.5.7}; it is therefore the integral closure of $Y'$ in $L\otimes_{K}K'$, which completes the proof.
\end{proof}

\begin{corollary}[18.10.14]
\label{IV.18.10.14}
Let $Y$ and $Y'$ be two normal and integral preschemes, $K$ and $K'$ their respective fields of rational functions, and $g:Y'\to Y$ a dominant morphism, so that $K'$ is an extension of $K$.
Let $L$ be a finite extension of $K$ unramified over $Y$, and let $L_1$ be a composite extension \textup{(Bourbaki, \emph{Alg.}, chap.~VIII, \textsection~8, def.~1) of} $L$ and $K'$.
Then $L_1$ is an extension of $K'$ unramified over $Y'$.
Furthermore, if $Y=\operatorname{Spec}(A)$ and $Y'=\operatorname{Spec}(A')$ are affine, and if $C$ is the integral closure of $A$ in $L$, then the subring $C_1=A[C,A']$ of $L_1$ is the integral closure of $A'$ in $L_1$.
\end{corollary}

\begin{proof}
Indeed, $L\otimes_{K}K'$ is a direct product of finite separable extensions of $K'$, and $L_1$
\oldpage[IV-4]{163}
is one of these extensions; hence the integral closure of $Y'$ in $L\otimes_{K}K'$ is a disjoint union of preschemes, one of which is the integral closure of $Y'$ in $L_1$.
The corollary therefore follows immediately from \sref{IV.18.10.13}.

When, for example, $A$ is the ring of integers of an algebraic number field $K$, and $K'$ and $L$ are algebraic extensions of $K$, there are classical examples showing that the relation $C_1=A[C,A']$ fails when $L$ is not unramified over $A$.
\end{proof}

\begin{env}[18.10.15]
\label{IV.18.10.15}
Suppose that $Y=\operatorname{Spec}(A)$ is affine, $A$ being integral and integrally closed; let $K$ be its field of fractions and $L$ a finite separable extension of $K$, and suppose that the integral closure $C$ of $A$ in $L$ is a \emph{projective} $A$-module of finite type (which will be the case for example if $A$ is a Dedekind ring (Bourbaki, \emph{Alg.\ comm.}, chap.~VII, \textsection~4, n\textsuperscript{o}~10, prop.~22)).
To say that $L$ is unramified over $A$ means that $C$ is \emph{\'etale} \sref{IV.18.10.10}; by \sref{IV.18.2.7},~(ii), this is equivalent to saying that the \emph{discriminant} $d_{C/A}$ of $\operatorname{Spec}(C)$ over $\operatorname{Spec}(A)$ is \emph{invertible} in $A$.
In the particular case where $C$ is a \emph{free} $A$-module and $(x_i)_{1\leq i\leq n}$ a basis of $C$ over $A$, this means that $\det(\operatorname{Tr}_{C/A}(x_ix_j))$ is invertible in $A$.
\end{env}

\begin{theorem}[18.10.16]
\label{IV.18.10.16}
Let $Y$ be an integral prescheme, $\eta$ its generic point, and $f:X\to Y$ a separated morphism locally of finite type.
Suppose that every irreducible component of $X$ dominates $Y$ and that the generic fibre $X_\eta=f^{-1}(\eta)$ is a finite prescheme over $K=\boldsymbol{k}(\eta)$; then $(X_\eta)_{\mathrm{red}}=\operatorname{Spec}(L)$, where $L=\prod_iL_i$ is a product of finite extensions $L_i$ of $K$.
Set
\[
  n=[L:K]=\sum_i[L_i:K],\quad n_s=\sum_i[L_i:K]_s\quad\text{\textup{(sum of separable degrees of the $L_i$)}}.
\]
Let $y$ be a \textup{geometrically unibranch} point of $Y$, and denote by $n(y)$ the sum of the separable degrees over $\boldsymbol{k}(y)$ of the residue fields of the \textup{isolated} points of $f^{-1}(y)$.
Then:
\begin{enumerate}[label=\textup{(\roman*)}]
  \item One has $n(y)\leq n_s\leq n$.
  \item Suppose $X$ is reduced and the point $y$ is normal in $Y$.
For $n(y)=n$, it is necessary and sufficient that there exists a neighbourhood $U$ of $y$ in $Y$ such that the restriction $f^{-1}(U)\to U$ of $f$ is an \'etale and finite morphism.
\end{enumerate}
\end{theorem}

\begin{proof}
A) \emph{Reduction to the case where $f$ is quasi-finite.} ---
One knows \sref{IV.13.1.4} that the set $Z$ of isolated points $x\in X$ of $f^{-1}(f(x))$ is open in $X$ and by hypothesis $Z$ contains $f^{-1}(\eta)$.
Since $f^{-1}(\eta)$ is finite, $Z$ is an increasing filtered union of quasi-compact open sets $V_\lambda$ containing $f^{-1}(\eta)$ (dense in $X$ by hypothesis).
If $n_\lambda(y)$ is the sum of the separable degrees over $\boldsymbol{k}(y)$ of the residue fields of the points of $V_\lambda\cap f^{-1}(y)$, we have $n(y)=\sup n_\lambda(y)$; to prove (i), it will therefore suffice to show that $n_\lambda(y)\leq n_s$.
On the other hand, if $n(y)=n$, there is an index $\lambda$ such that $n_\lambda(y)=n$. Suppose that we have shown that there then exists a neighbourhood $U$ of $y$ in $Y$ such that the restriction $V_\lambda\cap f^{-1}(U)\to U$ of $f$ is an \'etale and finite morphism; since $f$ is separated, the canonical injection $V_\lambda\cap f^{-1}(U)\to f^{-1}(U)$ is then also a finite morphism \sref[II]{II.6.1.5},~(v), hence $V_\lambda\cap f^{-1}(U)$ is closed in $f^{-1}(U)$; but being everywhere dense in $f^{-1}(U)$, it is necessarily equal to it.

B) \emph{Reduction to the case where $X$ is reduced, $f$ is finite, and $Y=\operatorname{Spec}(\sh{O}_{Y,y})$.} ---
We note that the hypothesis on $f$ implies that it is also satisfied for $f_{\mathrm{red}}$, and that the conclusion of (i) concerns only $f_{\mathrm{red}}$.
One can therefore replace $f$ everywhere by $f_{\mathrm{red}}$ and thus suppose $X$ reduced.

\oldpage[IV-4]{164}
On the other hand, by virtue of \sref{IV.18.12.13} (${}^1$), one can (after having replaced $Y$ by an open affine neighbourhood $V$ of $y$ and $f$ by its restriction $f^{-1}(V)\to V$) factorise $f$ as $X\xrightarrow{j}X'\xrightarrow{g}Y$, where $g$ is \emph{finite} and $j$ an open immersion.
Since $X$ is reduced, one can replace $X'$ by the reduced closed sub-prescheme with underlying space $\overline{j(X)}$ \sref[I]{I.5.2.2}, hence suppose $X'$ reduced and such that every irreducible component of $X'$ dominates $Y$.
Moreover, $j(X)$ is an open dense set in $X'$ and has only a finite number of irreducible components by virtue of the hypothesis; hence \sref[0\textsubscript{\mdseries I}]{0.1.2.7} the irreducible components of $X'$ are the closures of those of $j(X)$, and by construction the fibre $g^{-1}(\eta)$ identifies with $f^{-1}(\eta)$.
If one supposes the theorem proved when $f$ is \emph{finite}, one can apply it to $g$, and as $f^{-1}(y)$ identifies with a sub-prescheme of $g^{-1}(y)$, the relation (i) for $g$ implies the same relation for $f$.
Suppose, on the other hand, that $n(y)=n$; then, by what precedes, one can apply the conclusion of (ii) to the morphism $g$, and after replacing $Y$ if necessary by a neighbourhood of $y$, one may suppose that $g$ is \'etale, and consequently so is $f$; moreover, one then has $f^{-1}(y)=g^{-1}(y)$.
As the morphism $g$ is closed and $j(X)$ is open in $X'$, there exists a neighbourhood $U$ of $y$ in $Y$ such that $j(X)\cap g^{-1}(U)=j(f^{-1}(U))$, which proves that (ii) is also true for $f$.

We are thus reduced to the case where $f$ is moreover a \emph{finite} morphism.
It is immediate that, to prove (i), one can restrict to the case where $Y=\operatorname{Spec}(\sh{O}_{Y,y})$.
Let us show that the same is true for proving (ii).
Indeed, suppose it has been proved that the condition $n(y)=n$ implies that $X\times_{Y}\operatorname{Spec}(\sh{O}_{Y,y})$ is \'etale and finite over $\operatorname{Spec}(\sh{O}_{Y,y})$. Since this implies that $f$ is flat and formally unramified at the points of $f^{-1}(y)$, and since $Y$ is \emph{integral}, it follows from \sref{IV.18.4.13} that $f$ is \emph{\'etale at the points of} $f^{-1}(y)$, hence also on a neighbourhood $V$ of $f^{-1}(y)$ in $X$; but since $f$ is finite, hence closed, $V$ contains an open set of the form $f^{-1}(U)$, where $U$ is a neighbourhood of $y$ in $Y$.

C) \emph{Reduction to the case where $Y=\operatorname{Spec}(\sh{O}_{Y,y})$ and $\sh{O}_{Y,y}$ is strictly local.} ---
Set $A=\sh{O}_{Y,y}$, and let $A'={}^{hs}A$ be the strict henselisation of $A$ \sref{IV.18.8.7}, which is an integral, geometrically unibranch local ring \sref{IV.18.8.16}.
Let $Y'=\operatorname{Spec}(A')$, and let $y'$ and $\eta'$ be the closed point and the generic point of $Y'$; $y'$ is above $y$ and since the morphism $g:Y'\to Y$ is flat \sref{IV.18.8.8}, $\eta'$ is above $\eta$ \sref{IV.2.3.4}.
Set $X'=X\times_{Y}Y'$, $f'=f_{(Y')}:X'\to Y'$; $f'$ is finite and since $g$ is flat, every irreducible component of $X'$ dominates $Y'$ \sref{IV.2.3.7}.
If $n'(y')$, $n'$ and $n_s'$ are the numbers defined for $f'$ in the same way as $n(y)$, $n$ and $n_s$ for $f$, we have $n_s'=n_s$ and $n'(y')=n(y)$ \sref[I]{I.6.4.8}, and it therefore amounts to proving the assertion (i) for $f$ or for $f'$.
Moreover, if $X$ is reduced and $y$ is a normal point of $Y$ (hence $A$ is integrally closed), then $A'$ is integrally closed \sref{IV.18.8.12}.
On the other hand, it follows from the definition \sref{IV.18.8.7}, from the remark \sref{IV.8.1.2},~$a$) and from the theorem of the double inductive limit, that $A'$ is the inductive limit of rings $A_\alpha$ such that the morphisms $\operatorname{Spec}(A_\alpha)\to\operatorname{Spec}(A)=Y$ are \'etale; consequently $X'$ is the projective limit of the $X_\alpha=X\otimes_{A}A_\alpha$, which are \'etale over $X$, hence reduced since $X$ is \sref{IV.17.5.7}; by passage to the limit,

\footnotetext{(${}^1$) The reader will verify that \sref{IV.18.10.16} is not used in the proof of \sref{IV.18.12.13}. If $f$ is supposed locally of finite presentation (resp.\ quasi-projective), one can replace \sref{IV.18.12.13} by \sref{IV.8.12.8} (resp.\ \sref{IV.8.12.6}).}

\oldpage[IV-4]{165}
we deduce that $X'$ is reduced \sref{IV.8.7.1}.
This being so, if $n(y)=n$, one has $n_s=n$, and consequently the residue fields at the points of $X_\eta$ are necessarily separable over $\boldsymbol{k}(\eta)$; the residue fields at the points of $X_{\eta'}'$ are then algebraic and separable over $\boldsymbol{k}(\eta')$ \sref{IV.4.6.1}, in other words $n_s'=n'=n$. It follows from what precedes that $n'(y')=n'$.
The morphism $f'$ therefore has the same properties as $f$, and if one proves that $f'$ is \'etale, it will follow from \sref{IV.17.7.1} that $f$ is \'etale.

D) \emph{End of the proof.} ---
Suppose therefore that $A$ is strictly local.
Since $A$ is Henselian and the morphism $f$ is finite, it follows from \sref{IV.18.5.11} that if $x_j$ ($1\leq j\leq m$) are the distinct points of $f^{-1}(y)$, $X$ is the \emph{sum} of open sub-preschemes $\operatorname{Spec}(\sh{O}_{X,x_j})=X_j$, which are the connected components of $X$.
By hypothesis, each of them dominates $Y$, hence meets $f^{-1}(\eta)$, and so the number of points of $f^{-1}(\eta)$ is $\geq m$; \emph{a fortiori} we have $n_s\geq m$.
But since $\boldsymbol{k}(y)$ is separably closed and the $\boldsymbol{k}(x_j)$ are algebraic over $\boldsymbol{k}(y)$, they are necessarily \emph{radicial} extensions of $\boldsymbol{k}(y)$; hence $m=n(y)$, and this proves (i).
Suppose now the hypotheses of (ii) satisfied; the relations $n(y)=n_s=n$ imply first of all, since every irreducible component of $X$ dominates $Y$, that each of the $X_j$ is \emph{irreducible}; moreover, as the $X_j$ are \emph{reduced}, they are \emph{integral}; finally the relations $m=n_s=n$ imply that for the generic point $\xi_j$ of $X_j$, $\boldsymbol{k}(\xi_j)$ is a separable extension of $\boldsymbol{k}(\eta)$ of degree $1$, hence isomorphic to $\boldsymbol{k}(\eta)$; in other words, the restriction $f|X_j:X_j\to Y$ of $f$ is \emph{finite and birational}; moreover, $X_j$ is \emph{integral} and by hypothesis $Y$ is \emph{normal}; hence \sref{IV.8.12.10.1} $f|X_j$ is an \emph{isomorphism}, which proves that $f$ is \emph{\'etale}, and finishes the proof of \sref{IV.18.10.16}.
\end{proof}

\begin{remark}[18.10.17]
\label{IV.18.10.17}
The method of ``\'etale localisation'' used in the proof of \sref{IV.18.10.16} also allows one to improve the results of \sref{IV.15.5.1} by eliminating the Noetherian hypothesis.
One has first of all the following result which generalises \sref{IV.15.5.2}:
\end{remark}

\begin{env}[18.10.17.1]
\label{IV.18.10.17.1}
\emph{Let $f:X\to Y$ be a separated and quasi-finite morphism, $y$ a point of $Y$ such that $f$ is universally open at the points of $f^{-1}(y)$.
Then, for every generisation $y'$ of $y$, we have $n(y)\leq n(y')$.}
\end{env}

\begin{proof}
One can replace $Y$ by the reduced closed sub-prescheme $Z$ with underlying space $\overline{\{y'\}}$ and $X$ by $f^{-1}(Z)$, hence suppose $Y$ integral with generic point $y'$.
Using \sref{IV.18.12.13} as in \sref{IV.18.10.16},~B), one can suppose that $X$ is an open dense subset of a prescheme $X'$ and that $f$ is the restriction to $X$ of a \emph{finite} morphism $g:X'\to Y$.
Let us show that $g^{-1}(y')=f^{-1}(y')$.
Indeed, let $i:g^{-1}(y')\to X'$ be the canonical injection, which is a flat morphism since $y'$ is the generic point of $Y$ \sref[I]{I.3.6.5}; one can write $f^{-1}(y')=i^{-1}(X)$.
On the other hand, the canonical injection $j:X\to X'$ is a morphism of finite type since the composite $g\circ j=f$ is of finite type and $g$ is separated \sref{IV.1.5.4}; hence $X$ is a retrocompact open subset of $X'$, and consequently pro-constructible in $X'$ \sref{IV.1.9.5},~(v).
One concludes therefore from \sref{IV.2.3.10} that $i^{-1}(\overline{X})=\overline{i^{-1}(X)}$, in other words $f^{-1}(y')$ is dense in $g^{-1}(y')$; but as $g^{-1}(y')$ is discrete, we necessarily have $f^{-1}(y')=g^{-1}(y')$.

We are thus reduced to proving that the sum of the numbers $[\boldsymbol{k}(x):\boldsymbol{k}(y)]_s$ where $x$ ranges over the points of $g^{-1}(y)$ where $g$ is universally open, is at most equal to
\oldpage[IV-4]{166}
$n(y')$ (for the morphism $g$).
Using then the fact that the property of being universally open at a point is preserved by base change, one shows, as in \sref{IV.18.10.16},~C), that one can reduce to the case where $Y=\operatorname{Spec}(\sh{O}_{Y,y})$, with $\sh{O}_{Y,y}$ \emph{strictly local}.
Then, with the notation of \sref{IV.18.10.16},~D), $X$ is the sum of the open sub-preschemes $X_j=\operatorname{Spec}(\sh{O}_{X,x_j})$, where $x_j$ ranges over $g^{-1}(y)$; the hypothesis that $g$ is \emph{open} at one of these points $x_j$ implies that the corresponding sub-prescheme $X_j$ dominates $Y$ \sref{IV.1.10.3}; one concludes as in \sref{IV.18.10.16},~D).

We then deduce from \sref{IV.18.10.17.1} that \emph{the conclusions of} \sref{IV.15.5.1} \emph{are still valid when one no longer supposes $Y$ locally Noetherian, but only that $f$ is separated, quasi-finite and of finite presentation}.
\end{proof}

Indeed, to prove assertion (i) of \sref{IV.15.5.1}, one remarks that, since $f$ is of finite presentation, the set $E$ of points $z\in Y$ such that $n(z)\geq n(y)$ is locally constructible \sref{IV.9.7.9}; to show that $y$ is interior to $E$, it suffices, by virtue of \sref{IV.1.10.1}, to see that every generisation $y'$ of $y$ belongs to $E$, which is nothing other than \sref{IV.18.10.17.1}.

To prove assertion (ii) of \sref{IV.15.5.1}, one may, since $f$ is of finite presentation, use \sref{IV.8.10.5},~(xii) applied according to the method of \sref{IV.8.1.2},~\emph{a)}, and thus first suppose that $Y=\operatorname{Spec}(\sh{O}_{Y,y})$. Applying next \sref{IV.2.7.1},~(vii), one sees that after the faithfully flat base change $Y'\to Y$, where $Y'=\operatorname{Spec}({}^{hs}\sh{O}_{Y,y})$, one may suppose that the ring $\sh{O}_{Y,y}$ is \emph{strictly local}.
Applying \sref{IV.18.5.11},~\emph{c)}, we see then that if $x_j$ ($1\leq j\leq n$) are the points of $f^{-1}(y)$, $X$ is the \emph{sum} of $n$ open sub-preschemes $X_j=\operatorname{Spec}(\sh{O}_{X,x_j})$ which are \emph{finite} over $Y$, and of an open prescheme $X''$.
If one proves that $X''=\varnothing$, one will have shown that $f$ is finite, hence proper. Now the fields $\boldsymbol{k}(x_j)$, being algebraic extensions of $\boldsymbol{k}(y)$, are \emph{radicial}, and consequently $n(y)=n$; since $f$ is open at each point $x_j$, $X_j$ dominates $Y$ \sref{IV.1.10.3}, hence the restriction of $f$ to $X_j$, being a \emph{finite} morphism, is \emph{surjective} \sref[II]{II.6.1.10}. The hypothesis that $z\leadsto n(z)$ is constant on $Y$ therefore implies that $X''\cap f^{-1}(z)=\varnothing$ for every $z\in Y$, that is, $X''=\varnothing$.

Finally, the proof of (iii) reproduces without change that given in \sref{IV.15.5.1}.

\begin{env}[18.10.17.2]
\label{IV.18.10.17.2}
By analogous methods using strict henselisation, one can, in most of the results of \textsection\textsection~14 and 15, eliminate the Noetherian hypotheses.
\end{env}

\begin{lemma}[18.10.18]
\label{IV.18.10.18}
Let $X$, $Y$ be two preschemes, $f:X\to Y$ a birational morphism (${}^1$), $x$ a point of $X$.
For $f$ to be an isomorphism local at the point $x$, it is necessary and sufficient that $f$ be \'etale at the point $x$.
For $f$ to be an open immersion, it is necessary and sufficient that $f$ be \'etale and separated.
\end{lemma}

\footnotetext{(${}^1$) We understand by this the notion defined in \sref{IV.6.15.4}, but where we do not suppose $X$ and $Y$ reduced.}

\begin{proof}
The conditions stated being trivially necessary, all that is needed is to see that they are sufficient.
For the first assertion, the question is local on $X$ and $Y$; hence one may suppose that $f$ is \'etale and that $X$ and $Y$ are affine, so that $f$ is \emph{separated}, and one is reduced to proving the second assertion.
By virtue of \sref{IV.17.9.1}, it is a question of proving that $f$ is \emph{radicial}.
Let $y\in Y$, and let us prove that $f^{-1}(y)$ is radicial on $\boldsymbol{k}(y)$.
The base change $\operatorname{Spec}(\sh{O}_{Y,y})\to Y$ not changing the fact that $f$ is \'etale, separated, and birational, one can restrict to the case where $Y=\operatorname{Spec}(A)$, with $A=\sh{O}_{Y,y}$.
\oldpage[IV-4]{167}
Let $A'={}^{hs}A$, which is a local ring and a faithfully flat $A$-module \sref{IV.18.8.8}, the homomorphism $A\to A'$ being local.
Set $Y'=\operatorname{Spec}(A')$, $X'=X_{(Y')}$, $f'=f_{(Y')}:X'\to Y'$, so that $f'$ is \'etale and separated; moreover, as the morphism $Y'\to Y$ is flat, the reasoning of \sref{IV.6.15.4.1} shows that $f'$ is also birational.
If $y'$ is the closed point of $Y'$, it suffices, by \sref{IV.2.6.1},~(v), to prove that $f'^{-1}(y')$ is radicial over $\boldsymbol{k}(y')$; since $f'$ is \'etale, it is enough to prove that $f'^{-1}(y')$ contains \emph{at most one point} \sref{IV.17.6.1},~$c'$).

Suppose therefore $Y=\operatorname{Spec}(A)$, $A$ strictly local, $f$ separated, \'etale and birational, and let us show that $f^{-1}(y)$ can contain only a single point.
Indeed, if $f^{-1}(y)$ contained two distinct points $x_1$, $x_2$, there would exist two $Y$-sections $u_1$, $u_2$ of $f$ such that $u_1(y)=x_1$ and $u_2(y)=x_2$ \sref{IV.18.8.1}; but since $f$ is \'etale and separated, $u_1$ and $u_2$ would be isomorphisms of $Y$ onto two connected components (open) of $X$ \sref{IV.17.9.4} and there would then be two maximal points of $X$ at least above every maximal point of $Y$, which contradicts the hypothesis that $f$ is \emph{birational}.
\end{proof}

\begin{proposition}[18.10.19]
\label{IV.18.10.19}
Let $Y$ be a reduced prescheme whose set of irreducible components is locally finite, $f:X\to Y$ a separated morphism, locally of finite type and formally unramified, $g$ a $Y$-rational section of $f$ \textup{(i.e.}\ a $Y$-rational map from $Y$ to $X$\textup{)}, $U$ the domain of definition of $g$ \sref[I]{I.7.2.1}, and let $Z$ be the closure of $g(U)$ in $X$.
Then, for every point $y\in Y\setmin U$ such that $Y$ is geometrically unibranch at the point $y$, one has $Z_y=Z\cap f^{-1}(y)=\varnothing$.
In particular, if $Y$ is geometrically unibranch, $g(U)$ is a subset of $X$ that is both open and closed.
If, for every \emph{irreducible} closed subset $T$ of $X$, $f(T)$ is closed in $Y$, $g$ is defined at every geometrically unibranch point of $Y$.
\end{proposition}

\begin{proof}
Since $Y$ is reduced and $f$ separated, it follows from \sref[I]{I.7.2.2} that there exists a $U$-section $u$ of $f^{-1}(U)$ belonging to the class $g$; moreover, since $f$ is formally unramified, \sref{IV.17.4.1.2} shows that $u$ is an isomorphism of $U$ onto the prescheme induced on the open set $u(U)$ of $X$.
If one denotes also by $Z$ the reduced sub-prescheme of $X$ having $Z$ for underlying space, $g(U)=u(U)$, being reduced, is induced also by $Z$, so the restriction $f'=f|Z$ of $f$ is a \emph{birational} morphism of $Z$ to $Y$, moreover formally unramified, since $f$ is \sref{IV.17.1.3}.
The question being local on $Y$, and since $Y$ is reduced and geometrically unibranch (hence integral) at the point $y$, one may, after replacing $Y$ by an open neighbourhood of $y$, restrict to the case where $Y$ is \emph{integral}; then $u(U)$ is irreducible, hence so is $Z$, and consequently $Z$ is integral as well.
Since $f'$ is dominant, \sref[I]{I.8.2.7} shows that the homomorphism $\sh{O}_{Y,y}\to\sh{O}_{Z,x}$ is injective for every $x\in Z_y$, and one concludes from \sref{IV.18.10.2},~(iii) that $f'$ is an \emph{\'etale} morphism. Being separated and birational, it is a local isomorphism at $x$ by \sref{IV.18.10.18}; but this would imply that the $U$-section $u$ extends to an open set strictly larger than $U$, contrary to the definition of $U$.
We therefore necessarily have $Z_y=\varnothing$, which establishes the first assertion; the second is an obvious consequence.
Moreover, under the hypotheses of the last assertion of the statement, as $Z$ is closed and irreducible, $f(Z)$ is closed in $Y$, hence equal to $Y$, i.e.\ $Z_y\neq\varnothing$ for all $y\in Y$, which finishes the proof.
\end{proof}

\begin{remark}[18.10.20]
\label{IV.18.10.20}
Let $Y$ be a locally Noetherian prescheme.
One says that a morphism $f:X\to Y$ is \emph{essentially proper} if it is \emph{locally of finite type} and if, for every $Y$-scheme of the form $Y'=\operatorname{Spec}(A)$, where $A$ is a discrete valuation ring, the canonical map \sref[II]{II.7.3.2.2} relative to the $Y'$-prescheme $X\times_{Y}Y'$ is bijective.
\oldpage[IV-4]{168}
The reasoning of \sref[II]{II.7.3.8} proves that this is the same as saying that $f$ (supposed locally of finite type) is \emph{separated} and that for every base change $Y''\to Y$, where $Y''$ is \emph{locally Noetherian}, the image by $f_{(Y'')}: X_{(Y'')}\to Y''$ of every closed \emph{irreducible} subset of $X''$ is \emph{closed}.
To say that $f$ is \emph{proper} (for $Y$ locally Noetherian) means therefore that $f$ is \emph{essentially proper and of finite type} \sref[II]{II.7.3.8}; but one encounters important examples of essentially proper morphisms that are not of finite type (for example certain ``Picard preschemes'' or certain ``preschemes of N\'eron-Severi'' (chap.~VI)).
It follows immediately from \sref{IV.18.10.19} that if $Y$ is \emph{locally Noetherian, reduced, and geometrically unibranch}, and if $f:X\to Y$ is \emph{essentially proper and unramified}, then \emph{every $Y$-rational section of $f$ is defined everywhere}.
\end{remark}

The following proposition generalises the smoothness criterion \sref{IV.17.15.5} for a morphism $X\to Y$ when $Y$ is not the spectrum of a field:

\begin{proposition}[18.10.21]
\label{IV.18.10.21}
Let $f:X\to Y$ be a morphism locally of finite presentation, $x$ a point of $X$, $y=f(x)$, and suppose that the ring $\sh{O}_{Y,y}$ is \emph{integral} and \emph{geometrically unibranch}.
Let $r$ be the minimum number of generators of the $\sh{O}_{X,x}$-module $(\varOmega^1_{X/Y})_x$ \textup{(equal to the \emph{rank} of the $\boldsymbol{k}(x)$-vector space $(\varOmega^1_{X/Y})_x\otimes_{\sh{O}_{X,x}}\boldsymbol{k}(x)$)}.
Then, for $f$ to be smooth at the point $x$, it is necessary and sufficient that, if $y'$ is the generic point of the unique irreducible component of $Y$ containing $y$, there exists a generisation $x'$ of $x$ such that $f(x')=y'$ and $\dim_{x'}f^{-1}(y')\geq r$.
\end{proposition}

\begin{proof}
If $f$ is smooth at the point $x$, it is also so in an open neighbourhood $U$ of $x$, hence in every generisation $x'$ of $x$, and $\dim_{x'}(f^{-1}(f(x')))=r$ in each of these points, by virtue of \sref{IV.17.10.2}.
As $f$ is moreover flat in $U$, there exists a generisation of $x$ above every generisation of $y$ \sref{IV.2.3.4}, which proves that the condition is necessary.
To show that it is sufficient, let us first note that by \sref{IV.17.5.1}, it suffices to prove that the morphism deduced from $f$ by the base change $\operatorname{Spec}(\sh{O}_{Y,y})\to Y$ is smooth at the point $x$; one can therefore suppose that $Y=\operatorname{Spec}(\sh{O}_{Y,y})$.

The question is local on $X$; hence, by \sref{IV.16.4.22} and \sref[0\textsubscript{\mdseries I}]{0.5.2.2}, one may suppose that $X$ is affine and that there exist $r$ sections $s_i$ ($1\leq i\leq r$) of $\sh{O}_X$ over $X$ such that the sections $d_{X/Y}(s_i)$ generate $\varOmega^1_{X/Y}$.
Let
\[
  g:X\to Z=Y[T_1,\ldots,T_r]=\mathbf{V}_{Y}^r
\]
be the $Y$-morphism corresponding to the homomorphism $\sh{O}_{Y}^r\to f_*(\sh{O}_X)$ of $\sh{O}_Y$-Modules defined by the $s_i$ (considered as sections of $f_*(\sh{O}_X)$ over $Y$) \sref[II]{II.1.2.7}.
We have the exact sequence \sref{IV.16.4.19}
\[
  g^*(\varOmega^1_{Z/Y})\to\varOmega^1_{X/Y}\to\varOmega^1_{X/Z}\to 0.
\]
As the $dT_i$ ($1\leq i\leq r$) generate $\varOmega^1_{Z/Y}$, it follows from the definition of $g$ that the homomorphism $g^*(\varOmega^1_{Z/Y})\to\varOmega^1_{X/Y}$ is surjective, hence $\varOmega^1_{X/Z}=0$, and $g$ is consequently \emph{unramified} \sref{IV.17.4.1}.
We shall now see that $g$ is \emph{\'etale} at the point $x$, and taking into account \sref{IV.17.3.8}, this will prove that $f$ is smooth at the point $x$.
As $\sh{O}_{Y,y}$ is assumed to be integral and geometrically
\oldpage[IV-4]{169}
unibranch, so is $\sh{O}_{Z,z}$ by \sref{IV.17.5.7}, where $z=g(x)$, since the structural morphism $Z\to Y$ is smooth \sref{IV.17.3.8}.
By virtue of \sref{IV.18.10.1}, it suffices to show that the homomorphism $\sh{O}_{Z,z}\to\sh{O}_{X,x}$ is injective; it amounts to the same thing to see that the morphism $\operatorname{Spec}(\sh{O}_{X,x})\to\operatorname{Spec}(\sh{O}_{Z,z})$ corresponding is \emph{dominant}, since $\sh{O}_{Z,z}$ is integral \sref[I]{I.1.2.7}.
Let $y'$ be the generic point of $Y=\operatorname{Spec}(\sh{O}_{Y,y})$, and let $z'$ be the generic point of the integral scheme $Z=\operatorname{Spec}(\sh{O}_{Y,y}[T_1,\ldots,T_r])$; if $h:Z\to Y$ is the structural morphism, then $h(z')=y'$.
It suffices to prove that there exists in $f^{-1}(y')$ a generisation $x'$ of $x$ such that the image of $x'$ by the morphism $g_{y'}:f^{-1}(y')\to h^{-1}(y')$ equals $z'$.
Moreover, $z'$ is also the generic point of $h^{-1}(y')=\operatorname{Spec}(\boldsymbol{k}(y')[T_1,\ldots,T_r])$; since the morphism $g_{y'}$ is unramified, hence quasi-finite \sref{IV.17.4.2}, so is its restriction to every irreducible component of $f^{-1}(y')$; since $f^{-1}(y')$ and $h^{-1}(y')$ are locally Noetherian and $\dim(h^{-1}(y'))=r$, it follows from \sref{IV.5.4.1} that the restriction of $g_{y'}$ to every irreducible component \emph{of dimension} $\geq r$ of $f^{-1}(y')$ is a dominant morphism.
Now, there exists by hypothesis such a component containing a generisation of $x$; \emph{a fortiori} its generic point is a generisation of $x$, whence the conclusion.
\end{proof}

\subsection{Application to Noetherian local algebras complete over a field}
\label{subsection:IV.18.11}

The following lemma generalizes \sref[0]{0.21.9.1} and \sref[0]{0.21.9.2}:

\begin{lemma}[18.11.1]
\label{IV.18.11.1}
Let $k$ be a field, let $A$ be a complete Noetherian local ring that is a $k$-algebra, and let $K$ be the residue field of $A$.
\begin{enumerate}[label=\textup{(\roman*)}]
  \item If the $K$-vector space $\Omega^1_{K/k}$ is of finite rank, the $A$-module $\hat{\Omega}^1_{A/k}$ is of finite type.
  \item If, moreover, $A$ is formally smooth as a $k$-algebra (for the adic topology), $\hat{\Omega}^1_{A/k}$ is a free $A$-module of rank
\[
  \dim(A)+\operatorname{rg}_{K}(\Omega^1_{K/k})\setmin\operatorname{rg}_{K}(\Upsilon_{K/k}).
\]
\end{enumerate}
Furthermore, for every subfield $k_0$ of $k$ such that $\Omega^1_{k/k_0}$ is a finite-dimensional $k$-vector space, $\hat{\Omega}^1_{A/k_0}$ is a free $A$-module of rank
\[
  \dim(A)+\operatorname{rg}_{K}(\Omega^1_{K/k})\setmin\operatorname{rg}_{K}(\Upsilon_{K/k})+\operatorname{rg}_{k}(\Omega^1_{k/k_0}).
\]
\end{lemma}

\begin{proof}
Assertion (i) is recalled here, having been proved in \sref[0]{0.20.7.15}.
To prove (ii), note that this assertion was in fact established by the proof of \sref[0]{0.21.9.2}; the statement of \sref[0]{0.21.9.2} uses the finiteness of the transcendence degree of $K$ over $k$ only through Cartier's equality \sref[0]{0.21.7.1}.
\end{proof}

\begin{lemma}[18.11.2]
\label{IV.18.11.2}
Let $k$ be a field of characteristic $p>0$, and let $C$ be an integral complete Noetherian local $k$-algebra whose residue field is a finite extension of $k$; let $n$ be its dimension and $L$ its field of fractions.
Then $\Omega^1_{L/k}$ and $\Upsilon_{L/k}$ are $L$-vector spaces of finite rank, and one has
\begin{equation}
\label{IV.18.11.2.1}
\tag{18.11.2.1}
  \operatorname{rg}_{L}(\Omega^1_{L/k})\setmin\operatorname{rg}_{L}(\Upsilon_{L/k})\geq n.
\end{equation}
\oldpage[IV-4]{170}
Suppose furthermore that $[k:k^p]<+\infty$.
Then one has the equality
\begin{equation}
\label{IV.18.11.2.2}
\tag{18.11.2.2}
  \operatorname{rg}_{L}(\Omega^1_{L/k})\setmin\operatorname{rg}_{L}(\Upsilon_{L/k})=n.
\end{equation}
Furthermore, $\Omega^1_{C/k}$ is then isomorphic to $\hat{\Omega}^1_{C/k}$, and hence, by \sref{IV.18.11.1}, is a $C$-module of finite type.
\end{lemma}

\begin{proof}
Indeed, by \sref[0]{0.19.8.9}, there exists a $k$-subalgebra $C_0$ of $C$ isomorphic to the formal power series algebra $k[[T_1,\ldots,T_n]]$ and such that $C$ is finite over $C_0$.
Consequently, $L$ is a \emph{finite} extension of the field of fractions $L_0=k((T_1,\ldots,T_n))$ of $C_0$.
Applying \sref[0]{0.20.6.17.1} to the prime subfield of $k$ and to the three fields $k\subset L_0\subset L$ (and taking \sref[0]{0.20.6.21},~(i), into account) gives the exact sequence of $L$-vector spaces
\[
  0\to\Upsilon_{L_0/k}\otimes_{L_0}L\to\Upsilon_{L/k}\to\Upsilon_{L/L_0}\to\Omega^1_{L_0/k}\otimes_{L_0}L\to\Omega^1_{L/k}\to\Omega^1_{L/L_0}\to 0.
\]
Since $L$ is a finite extension of $L_0$, the $L$-vector spaces $\Omega^1_{L/L_0}$ and $\Upsilon_{L/L_0}$ have finite and equal ranks, by Cartier's equality \sref[0]{0.21.7.1}. Since $L_0$ is separable over $k$ \sref[0]{0.21.9.6.4}, one has $\Upsilon_{L_0/k}=0$ \sref[0]{0.20.6.3}. The preceding exact sequence therefore already shows that $\Upsilon_{L/k}$ has finite rank and that, in every case \sref[0\textsubscript{\mdseries III}]{0.11.10.2},
\[
  \operatorname{rg}_{L}(\Omega^1_{L/k})\setmin\operatorname{rg}_{L}(\Upsilon_{L/k})=\operatorname{rg}_{L_0}(\Omega^1_{L_0/k})\setmin\operatorname{rg}_{L_0}(\Upsilon_{L_0/k}).
\]
It therefore suffices to prove \sref{IV.18.11.2.1} and \sref{IV.18.11.2.2} after replacing $C$ and $L$ by $C_0$ and $L_0$. Since $L_0$ is separable over $k$ \sref[0]{0.21.9.6.4}, one has $\Upsilon_{L_0/k}=0$ \sref[0]{0.20.6.3}; on the other hand,
\[
  \Omega^1_{L_0/k}=\Omega^1_{C_0/k}\otimes_{C_0}L_0
  \qquad\text{\sref[0]{0.20.5.9}.}
\]
We know from \sref[0]{0.21.9.4} that, if $C_1=k[[T_1^p,\ldots,T_n^p]]$, then $\hat{\Omega}^1_{C_0/k}$ identifies with $\Omega^1_{C_0/C_1}$. Moreover,
$\Omega^1_{C_0/k}=\Omega^1_{C_0/k[C_0^p]}$ \sref[0]{0.21.1.5}, and, since $C_0^p=k^p[[T_1^p,\ldots,T_n^p]]$, one has $k[C_0^p]\subset C_1$, with \emph{equality when} $[k:k^p]<+\infty$. The exact sequence
\[
  \Omega^1_{C_0/k[C_0^p]}\longrightarrow\Omega^1_{C_0/C_1}\longrightarrow 0
  \qquad\text{\sref[0]{0.20.5.7}}
\]
therefore yields an exact sequence $\Omega^1_{L_0/k}\to\hat{\Omega}^1_{C_0/k}\otimes_{C_0}L_0\to0$. Since $\hat{\Omega}^1_{C_0/k}$ is a free $C_0$-module of rank $n$ \sref[0]{0.21.9.3}, it follows that $\operatorname{rg}_{L_0}(\Omega^1_{L_0/k})\geq n$ in every case, with equality when $[k:k^p]<+\infty$. This proves \sref{IV.18.11.2.1} and \sref{IV.18.11.2.2}.

Finally, when $[k:k^p]<+\infty$, to see that $\Omega^1_{C/k}\to\hat{\Omega}^1_{C/k}$ is an isomorphism, it suffices to prove that $\Omega^1_{C/k}$ is a finite-type $C$-module, since $C$ is a complete Noetherian local ring \sref[0\textsubscript{\mdseries I}]{0.7.3.3}. Since $\Omega^1_{C/k}\cong\Omega^1_{C/k[C^p]}$ and $k[C_0^p]\subset k[C^p]$, this reduces to proving that $C$ is finite over $k[C_0^p]$ \sref[0]{0.20.4.7}; this follows because $C$ is finite over $C_0$, while $C_0$ is finite over $k[C_0^p]$ under the hypothesis $[k:k^p]<+\infty$.
\end{proof}

\begin{lemma}[18.11.3]
\label{IV.18.11.3}
Let $k$ be a field, let $p$ be its characteristic exponent, and suppose that $[k:k^p]<+\infty$. Let $A$ be a \textup{complete} Noetherian local $k$-algebra whose residue field is a \textup{finite} extension of $k$. Let $\mathfrak{p}$ be a prime ideal of $A$ such that $A_{\mathfrak{p}}$ is \textup{geometrically regular} over $k$ \sref{IV.6.7.6}, so that there is a unique minimal prime ideal $\mathfrak{q}$ of $A$ contained in $\mathfrak{p}$. Then $(\hat{\Omega}^1_{A/k})_{\mathfrak{p}}$ is a free $A_{\mathfrak{p}}$-module of rank $\dim(A/\mathfrak{q})$.
\end{lemma}

\begin{proof}
Since $A$ is Noetherian and $\operatorname{Spec}(A)$ is regular (and hence, a fortiori, integral) at $\mathfrak{p}$, the point $\mathfrak{p}$ lies on only one irreducible component of $\operatorname{Spec}(A)$. Thus it contains
\oldpage[IV-4]{171}
only one minimal ideal $\mathfrak{q}$ of $A$, and moreover $\mathfrak{q}_{\mathfrak{p}}=0$. Put $B=A/\mathfrak{q}$. The sequence \sref[0]{0.20.7.20}
\[
  \mathfrak{q}/\mathfrak{q}^2\longrightarrow
  \hat{\Omega}^1_{A/k}\widehat{\otimes}_{A}B
  \longrightarrow\hat{\Omega}^1_{B/k}\longrightarrow0
\]
is \emph{exact}. Indeed, $\hat{\Omega}^1_{A/k}$ is a finite-type $A$-module by \sref{IV.18.11.1}, and hence
\[
  \hat{\Omega}^1_{A/k}\widehat{\otimes}_{A}B
  =\hat{\Omega}^1_{A/k}\otimes_{A}B
  =\hat{\Omega}^1_{A/k}/\mathfrak{q}.\hat{\Omega}^1_{A/k}.
\]
As this $B$-module is of finite type, every one of its submodules is closed \sref[0\textsubscript{\mdseries I}]{0.7.3.5}, taking \sref[0]{0.20.4.5} into account. The image of $\mathfrak{q}/\mathfrak{q}^2$ is dense in the kernel of the displayed homomorphism \sref[0]{0.20.7.20}, and is therefore equal to that kernel. Localizing at $\mathfrak{p}$ also shows that the canonical homomorphism $(\hat{\Omega}^1_{A/k})_{\mathfrak{p}}\to(\hat{\Omega}^1_{B/k})_{\mathfrak{p}/\mathfrak{q}}$ is \emph{bijective}. We may consequently reduce to the case $\mathfrak{q}=0$, that is, to the case where $A$ is \emph{integral}. We distinguish two cases.

I) $p>1$. Apply \sref{IV.18.11.2} to the quotient algebra $C=A/\mathfrak{p}$, whose field of fractions $K$ is the residue field of $A_{\mathfrak{p}}$. Thus
\begin{equation}
\label{IV.18.11.3.1}
\tag{18.11.3.1}
  \operatorname{rg}_{K}(\Omega^1_{K/k})\setmin
  \operatorname{rg}_{K}(\Upsilon_{K/k})=\dim(A/\mathfrak{p}).
\end{equation}
Now $A_{\mathfrak{p}}$ is \emph{formally smooth} as a $k$-algebra for its $\mathfrak{p}A_{\mathfrak{p}}$-adic topology \sref[0]{0.19.6.6}; hence $\Omega^1_{A_{\mathfrak{p}}/k}$ is formally projective \sref[0]{0.20.4.9} for that topology \sref[0]{0.20.4.5}. Moreover,
$\Omega^1_{A_{\mathfrak{p}}/k}=(\Omega^1_{A/k})_{\mathfrak{p}}$ \sref[0]{0.20.5.9} is of finite type by \sref{IV.18.11.2} applied to $A$. For every integer $j$, the module
\[
  \Omega^1_{A_{\mathfrak{p}}/k}/\mathfrak{p}^{j+1}\Omega^1_{A_{\mathfrak{p}}/k}
\]
is therefore projective over $A_{\mathfrak{p}}/\mathfrak{p}^{j+1}A_{\mathfrak{p}}$, of rank
$m=\operatorname{rg}_{K}(\Omega^1_{A_{\mathfrak{p}}/k}\otimes_{A_{\mathfrak{p}}}K)$ \sref[0]{0.19.2.4}. It follows from \sref[0\textsubscript{\mdseries III}]{0.10.2.1} and \sref[0\textsubscript{\mdseries III}]{0.10.1.3} that $\Omega^1_{A_{\mathfrak{p}}/k}$ is a \emph{free $A_{\mathfrak{p}}$-module of rank $m$}. Let $A'=(A_{\mathfrak{p}})^{{\widehat{\,}}}$ be the \emph{completion} of $A_{\mathfrak{p}}$ for its $\mathfrak{p}A_{\mathfrak{p}}$-adic topology. The algebra $A'$ is again formally smooth over $k$ for its adic topology \sref[0]{0.19.3.6}, and \sref[0]{0.20.7.14} and \sref[0]{0.20.4.5} give $\hat{\Omega}^1_{A'/k}=\hat{\Omega}^1_{A_{\mathfrak{p}}/k}$. Thus $\hat{\Omega}^1_{A'/k}$ is a \emph{free} $A'$-module \emph{of rank $m$}. By \sref{IV.18.11.1} and $\dim(A')=\dim(A_{\mathfrak{p}})$ \sref[0]{0.16.2.4},
\begin{equation}
\label{IV.18.11.3.2}
\tag{18.11.3.2}
  m=\dim(A_{\mathfrak{p}})+
  \bigl(\operatorname{rg}_{K}(\Omega^1_{K/k})\setmin
  \operatorname{rg}_{K}(\Upsilon_{K/k})\bigr).
\end{equation}
By \sref{IV.18.11.3.1}, this is $m=\dim(A/\mathfrak{p})+\dim(A_{\mathfrak{p}})$. Since the complete Noetherian local ring $A$ is a quotient of a regular ring \sref[0]{0.19.8.8}, one has
$\dim(A)=\dim(A/\mathfrak{p})+\dim(A_{\mathfrak{p}})$ by \sref[0]{0.16.5.12}, which proves the assertion in this case.

II) $p=1$. As above, $\dim(A)=\dim(A/\mathfrak{p})+\dim(A_{\mathfrak{p}})$. Put $n=\dim(A)$, $r=\dim(A/\mathfrak{p})$, and $s=\dim(A_{\mathfrak{p}})$. We shall prove that $(\hat{\Omega}^1_{A/k})_{\mathfrak{p}}$ is a free $A_{\mathfrak{p}}$-module of rank $n$. We first prove the following lemma.

\begin{lemma}[18.11.3.3]
\label{IV.18.11.3.3}
Let $A$ be a Noetherian local ring and $\mathfrak{p}$ a prime ideal such that $A_{\mathfrak{p}}$ is Cohen--Macaulay. For every system of parameters $(t_i)_{1\leq i\leq s}$ of $A_{\mathfrak{p}}$, there exists a sequence $(x_i)_{1\leq i\leq s}$ of elements of $\mathfrak{p}$ such that the $x_i$ form part of a system of parameters of $A$
\oldpage[IV-4]{172}
and, for every $i$, $x_i/1$ is congruent to $t_i$ modulo $t_1A_{\mathfrak{p}}+\ldots+t_{i-1}A_{\mathfrak{p}}$. In particular, if $A_{\mathfrak{p}}$ is regular and the $t_i$ form a \textup{regular} system of parameters of $A_{\mathfrak{p}}$, the same is true of the $x_i/1$.
\end{lemma}

\begin{proof}[Proof of \sref{IV.18.11.3.3}]
This lemma is itself a consequence of the following one.

\begin{lemma}[18.11.3.4]
\label{IV.18.11.3.4}
Let $A$ be a Noetherian local ring, let $\mathfrak{p}$ be a prime ideal of $A$, and let $x\in A$. Then there exists $x'\in A$ such that $x'/1=x/1$ in $A_{\mathfrak{p}}$ and $x'$ belongs to none of the minimal prime ideals of $A$ not contained in $\mathfrak{p}$.
\end{lemma}
We first show how \sref{IV.18.11.3.4} implies \sref{IV.18.11.3.3}. Multiplying the $t_i$, if necessary, by units of $A_{\mathfrak{p}}$, we may assume that $t_i=x_i/1$ with $x_i\in\mathfrak{p}$ for $1\leq i\leq s$. We argue by induction on $s$. Since $t_1$ forms part of a system of parameters of $A_{\mathfrak{p}}$, it belongs to no minimal prime ideal of $A_{\mathfrak{p}}$ \sref[0]{0.16.5.5}; hence no $x_1'\in A$ with $x_1'/1=t_1$ can belong to a minimal prime ideal of $A$ contained in $\mathfrak{p}$. By \sref{IV.18.11.3.4}, we may moreover choose $x_1'$ such that $x_1'/1=t_1$ (so $x_1'\in\mathfrak{p}$) and $x_1'$ belongs to \emph{no} minimal prime ideal of $A$. Thus $x_1'$ forms part of a system of parameters of $A$ \sref[0]{0.16.3.4} and \sref{IV.16.3.7}. We then apply induction to $A'=A/x_1'A$ and $\mathfrak{p}'=\mathfrak{p}/x_1'A$. Since $A'_{\mathfrak{p}'}=A_{\mathfrak{p}}/t_1A_{\mathfrak{p}}$, this ring is Cohen--Macaulay \sref[0]{0.16.5.5}, and the images $t_i'$ ($2\leq i\leq s$) form a system of parameters of it. The induction hypothesis applied to $A'$ and the $t_i'$ completes the construction. The last assertion follows because, in $A_{\mathfrak{p}}$, a system of parameters that generates the maximal ideal is a regular system of parameters \sref[0]{0.17.1.1}.

It remains to prove \sref{IV.18.11.3.4}. Let $(\mathfrak{p}'_k)_{1\leq k\leq r}$ be the minimal prime ideals of $A$ not contained in $\mathfrak{p}$ and containing $x$, and let $(\mathfrak{p}_j)_{1\leq j\leq n}$ be the other minimal prime ideals of $A$; we may assume $r\geq1$. No $\mathfrak{p}'_k$ contains $\bigcap_{1\leq j\leq n}\mathfrak{p}_j$, so there exists
$y\in\bigcap_{1\leq j\leq n}\mathfrak{p}_j$ belonging to none of the $\mathfrak{p}'_k$ (Bourbaki, \emph{Alg.\ comm.}, chap.~II, \textsection~1, no.~1, prop.~2). Moreover, $y/1$ belongs to every minimal prime ideal of $A_{\mathfrak{p}}$, and is therefore nilpotent. If $(y/1)^h=0$, the element $x'=x+y^h$ has the required properties: $x'/1=x/1$; since $y^h\notin\mathfrak{p}'_k$ and $x\in\mathfrak{p}'_k$, one has $x'\notin\mathfrak{p}'_k$ for $1\leq k\leq r$; and if $\mathfrak{p}_j$ is a minimal prime ideal not contained in $\mathfrak{p}$ with $x\notin\mathfrak{p}_j$, then $x'\notin\mathfrak{p}_j$ because $y\in\mathfrak{p}_j$.
\end{proof}

We return to the proof of \sref{IV.18.11.3} when $p=1$. By \sref{IV.18.11.3.3}, there exists a regular system of parameters $(t_i)_{1\leq i\leq s}$ of $A_{\mathfrak{p}}$ such that $t_i=x_i/1$, where the $x_i$ ($1\leq i\leq s$) belong to $\mathfrak{p}$ and form part of a system of parameters $(x_j)_{1\leq j\leq n}$ of $A$. Put $A_0=k[[T_1,\ldots,T_n]]$. Since $A$ is a complete $k$-algebra and the $x_j$ belong to its maximal ideal $\mathfrak{m}$, there is a local $k$-homomorphism $u:A_0\to A$ such that $u(T_j)=x_j$ for $1\leq j\leq n$ (Bourbaki, \emph{Alg.\ comm.}, chap.~III, \textsection~4, no.~5, prop.~6). If $\mathfrak{n}$ is the ideal generated by the $x_j$, it is an ideal of definition of $A$; hence \sref[0\textsubscript{\mdseries I}]{0.7.4.4} and \sref{IV.7.4.3} show that $u$ makes $A$ \emph{finite} over $A_0$.

Put $\mathfrak{p}_0=\sum_{j=1}^{s}A_0T_j$, $B_0=(A_0)_{\mathfrak{p}_0}$, and $B=A\otimes_{A_0}B_0$, so that $u_{\mathfrak{p}_0}:B_0\to B$ makes
\oldpage[IV-4]{173}
$B$ finite over $B_0$. Moreover, if $\mathfrak{p}'$ is the ideal of $B$ generated by $\mathfrak{p}$, then $B_{\mathfrak{p}'}=A_{\mathfrak{p}}$; since $\mathfrak{p}'$ contains $\mathfrak{p}_0B$ by construction, it lies over the maximal ideal $\mathfrak{p}_0B_0$ of $B_0$. The morphism $\operatorname{Spec}(B)\to\operatorname{Spec}(B_0)$ is \emph{unramified} at $\mathfrak{p}'$: indeed, $\boldsymbol{k}(\mathfrak{p}')$ is a finite extension of the characteristic-zero field $\boldsymbol{k}(\mathfrak{p}_0)$, hence is \emph{necessarily separable}, and $B_{\mathfrak{p}'}/\mathfrak{p}_0B_{\mathfrak{p}'}=\boldsymbol{k}(\mathfrak{p}')$ by the choice of the $x_j$ for $1\leq j\leq s$ \sref{IV.17.4.1}.

We now record the following lemma.

\begin{lemma}[18.11.3.5]
\label{IV.18.11.3.5}
Let $k$ be a field, and let $R$ and $S$ be complete Noetherian local $k$-algebras such that, if $K$ is the residue field of $S$, $\Omega^1_{K/k}$ has finite rank over $K$. Let $u:R\to S$ be a $k$-homomorphism making $S$ finite over $R$. Then $\hat{\Omega}^1_{S/R}=\Omega^1_{S/R}$, and the sequence
\begin{equation}
\label{IV.18.11.3.6}
\tag{18.11.3.6}
  \hat{\Omega}^1_{R/k}\widehat{\otimes}_{R}S\xrightarrow{v}\hat{\Omega}^1_{S/k}\xrightarrow{w}\Omega^1_{S/R}\to 0
\end{equation}
(cf.\ \sref[0]{0.20.7.17.3}) is exact.
\end{lemma}

\begin{proof}
The first assertion follows because $\Omega^1_{S/R}$ is a finite-type $S$-module \sref[0]{0.20.4.7} and by \sref[0\textsubscript{\mdseries I}]{0.7.3.6}. Moreover, \sref[0]{0.20.7.17.3} says that the image of $v$ is dense in $\ker(w)$ and that $w$ is surjective. By \sref{IV.18.11.1}, however, $\hat{\Omega}^1_{S/k}$ is a finite-type $S$-module, and therefore every one of its submodules is closed \sref[0\textsubscript{\mdseries I}]{0.7.3.5}. The lemma follows.
\end{proof}

Apply the lemma with $R=A_0$ and $S=A$. The module $\hat{\Omega}^1_{A_0/k}$ is free of rank $n$ over $A_0$ \sref[0]{0.21.9.3}; hence
$\hat{\Omega}^1_{A_0/k}\widehat{\otimes}_{A_0}A=\hat{\Omega}^1_{A_0/k}\otimes_{A_0}A$, and this is a free $A$-module of rank $n$. Since $\operatorname{Spec}(B)$ is unramified over $\operatorname{Spec}(B_0)$ at $\mathfrak{p}'$, one has
$(\Omega^1_{A/A_0})_{\mathfrak{p}}=\Omega^1_{B_{\mathfrak{p}'}/B_0}=0$ \sref{IV.16.4.15} and \sref{IV.17.4.1}. Localizing \sref{IV.18.11.3.6}, applied to $A_0$ and $A$, at $\mathfrak{p}$ therefore gives a \emph{surjective} homomorphism
\[
  v_{\mathfrak{p}}:(\hat{\Omega}^1_{A_0/k}\otimes_{A_0}A)_{\mathfrak{p}}
  \longrightarrow(\hat{\Omega}^1_{A/k})_{\mathfrak{p}}.
\]
Thus $(\hat{\Omega}^1_{A/k})_{\mathfrak{p}}$ has a system of $n$ generators. On the other hand, $\hat{\Omega}^1_{A/k}$ has \emph{rank $n$} by \sref[0]{0.21.9.5}, which applies to the complete integral ring $A$ by Cohen's theorem \sref[0]{0.19.8.8},~(ii), and because the field of fractions of $A$ has characteristic zero. Hence $(\hat{\Omega}^1_{A/k})_{\mathfrak{p}}$ also has rank $n$; its quotient by its torsion submodule, being generated by $n$ elements, is necessarily \emph{free}. It follows at once that the $n$ generators obtained above form a \emph{free} system, which proves the assertion.
\end{proof}

\begin{lemma}[18.11.4]
\label{IV.18.11.4}
Let $k$ be a field, $A$ a $k$-algebra, local Noetherian and complete, whose residue field is a finite extension of $k$.
Let $k'$ be an arbitrary extension of $k$, and set $A'=A\widehat{\otimes}_kk'$.
Then:
\begin{enumerate}[label=\textup{(\roman*)}]
  \item $A'$ is a semi-local Noetherian ring, complete, composed as a direct product of complete local rings $A_i'$ \textup{($1\leq i\leq r$)} which are $A$-modules faithfully flat, and whose residue fields are finite extensions of $k'$; if $\mathfrak{m}$ is the maximal ideal of $A$, $\mathfrak{m}A'$ is an ideal of definition of $A'$; and $\dim(A_i')=\dim(A)$ for all $i$.
  \item For every $i$, $\hat{\Omega}^1_{A_i'/k'}$ is canonically isomorphic to $\hat{\Omega}^1_{A/k}\otimes_{A}A_i'$.
\end{enumerate}
\end{lemma}

\begin{proof}
\oldpage[IV-4]{174}
The first assertions follow immediately from \sref{IV.7.5.5} and \sref[0\textsubscript{\mdseries I}]{0.6.6.2}; the fact that $\mathfrak{m}A'$ is an ideal of definition of $A'$ also follows from \sref{IV.7.5.5}, since $K$ is the residue field of $A$, $K\otimes_k k'$ is a finite $k'$-algebra.
Finally, $A_i'/\mathfrak{m}A_i'$ is one of the Artinian local rings occurring as direct factors of $K\otimes_k k'$ \sref{IV.7.5.5}, and is therefore zero-dimensional. Since $A_i'$ is flat over $A$, the equality $\dim(A_i')=\dim(A)$ follows from \sref{IV.6.1.2}.

(ii) By the hypothesis on $K$, $\hat{\Omega}^1_{A/k}$ is an $A$-module of finite type \sref{IV.18.11.1}; hence $\hat{\Omega}^1_{A/k}\otimes_{A}A'$ is complete and identifies with the completed tensor product $\hat{\Omega}^1_{A/k}\widehat{\otimes}_{A}A'$. By associativity of completed tensor products \sref[0\textsubscript{\mdseries I}]{0.7.7.4},
\[
  \hat{\Omega}^1_{A/k}\otimes_{A}A'
  =\hat{\Omega}^1_{A/k}\otimes_{A}(A\widehat{\otimes}_kk')
  \simeq\hat{\Omega}^1_{A/k}\widehat{\otimes}_kk'.
\]
Now $\Omega^1_{A/k}\otimes_kk'$ identifies with $\Omega^1_{A''/k'}$, where $A''=A\otimes_kk'$ \sref[0]{0.20.5.5}; and since $A'$ is by definition the separated completion of $A''$, the separated completion of $\Omega^1_{A''/k'}$ identifies by construction with that of $\Omega^1_{A'/k'}$ \sref[0]{0.20.7.4}. In other words, there is a canonical isomorphism
\[
  \hat{\Omega}^1_{A'/k'}\xrightarrow{\sim}
  \hat{\Omega}^1_{A/k}\otimes_{A}A'.
\]
The conclusion of (ii) now follows because $\Omega^1_{A'/k'}$ is the direct sum of the $\Omega^1_{A_i'/k'}$ \sref[0]{0.20.4.13}. Indeed, if $\mathfrak r$ is the radical of $A'$, its $\mathfrak r$-adic topology on $\Omega^1_{A'/k'}$ is the product of the $\mathfrak m_i'$-adic topologies on the $\Omega^1_{A_i'/k'}$, where $\mathfrak m_i'$ is the maximal ideal of $A_i'$. Thus the separated completion $\hat{\Omega}^1_{A'/k'}$ is the product of the separated completions $\hat{\Omega}^1_{A_i'/k'}$, and it remains only to apply \sref[0]{0.20.4.5}.
\end{proof}

\begin{proposition}[18.11.5]
\label{IV.18.11.5}
Let $k$ be a field, $A$ a $k$-algebra, local Noetherian and complete, whose residue field is a finite extension of $k$, $\mathfrak{p}$ a prime ideal of $A$ distinct from the maximal ideal $\mathfrak{m}$, such that there exists a minimal prime ideal $\mathfrak{q}\subset\mathfrak{p}$ for which $\dim(A/\mathfrak{q})=\dim(A)$ \textup{(which will in particular hold when $A$ is \emph{equidimensional})}, $m$ an integer $\geq 0$.
The following conditions are equivalent:
\begin{enumerate}[label=\emph{\alph*)}]
  \item The $A_\mathfrak{p}$-module $(\hat{\Omega}^1_{A/k})_\mathfrak{p}$ admits a system of $m$ generators.
  \item There exists a local $k$-homomorphism $u:B\to A$, where $B=k[[T_1,\ldots,T_m]]$, making $A$ a finite $B$-algebra, and such that the corresponding morphism $\operatorname{Spec}(A)\to\operatorname{Spec}(B)$ is unramified at $\mathfrak{p}$.
\end{enumerate}
\end{proposition}

\begin{proof}
To prove that $b)$ implies $a)$, note that by Lemma \sref{IV.18.11.3.5} one has the exact sequence
\[
  \hat{\Omega}^1_{B/k}\otimes_BA\to\hat{\Omega}^1_{A/k}\to\Omega^1_{A/B}\to 0
\]
since $A$ is a $B$-algebra finite; localising at $\mathfrak{p}$ and noting that by hypothesis $(\Omega^1_{A/B})_\mathfrak{p}=0$ \sref{IV.17.4.1}, one obtains a surjective homomorphism $(\hat{\Omega}^1_{B/k}\otimes_BA)_\mathfrak{p}\to(\hat{\Omega}^1_{A/k})_\mathfrak{p}$, and the conclusion follows because $\hat{\Omega}^1_{B/k}$ is a free $B$-module of rank $m$ \sref[0]{0.21.9.3}.

To prove that $a)$ implies $b)$, let us first prove the following lemmas.
\end{proof}

\begin{lemma}[18.11.5.1]
\label{IV.18.11.5.1}
Let $k$ be a field, $A$ a $k$-algebra, local Noetherian and complete, whose residue field is a finite extension of $k$, $\mathfrak{p}$ a prime ideal of $A$, $\mathfrak{q}$ a minimal prime ideal of $A$ contained in $\mathfrak{p}$.
Then the minimum number of generators of the $A_\mathfrak{p}$-module $(\hat{\Omega}^1_{A/k})_\mathfrak{p}$ is at least equal to $\dim(A/\mathfrak{q})$.
\end{lemma}

\oldpage[IV-4]{175}

\begin{proof}
Set $n=\dim(A/\mathfrak{q})$, and let $m$ be the minimum number of generators of the $A_\mathfrak{p}$-module $(\hat{\Omega}^1_{A/k})_\mathfrak{p}$, which is equal to $\operatorname{rg}_{\kres(\mathfrak{p})}((\hat{\Omega}^1_{A/k})_\mathfrak{p}\otimes_{A_\mathfrak{p}}\boldsymbol{k}(\mathfrak{p}))$ (Bourbaki, \emph{Alg.\ comm.}, chap.~II, \textsection~3, n\textsuperscript{o}~2, cor.~2 of prop.~4).
Note that $(\hat{\Omega}^1_{A/k})_\mathfrak{q}=(\hat{\Omega}^1_{A/k})_\mathfrak{p}\otimes_{A_\mathfrak{p}}A_\mathfrak{q}$, hence the minimum number of generators of the $A_\mathfrak{q}$-module $(\hat{\Omega}^1_{A/k})_\mathfrak{q}$ is \emph{at most equal} to $m$.
It suffices therefore to consider the case where $\mathfrak{p}=\mathfrak{q}$ is minimal.
In second place, let us show one can restrict to the case where $k$ is \emph{algebraically closed}.
Indeed, let $k'$ be an algebraic closure of $k$, and set $A'=A\widehat{\otimes}_kk'$, which is composed as a direct product of $k'$-local algebras $A_i'$ ($1\leq i\leq r$), whose residue fields are \emph{isomorphic to $k'$}, and which are $A$-modules faithfully flat \sref{IV.18.11.4}.
Applying \sref{IV.2.3.4} and \sref{IV.6.1.1}, one sees that in a given $A_i'$, there exists a minimal prime ideal $\mathfrak{q}_i'$ above $\mathfrak{q}$ such that $n'=\dim(A_i'/\mathfrak{q}_i')\geq\dim(A/\mathfrak{q})=n$.
On the other hand by \sref{IV.18.11.4}, the minimum $m'$ of generators of the $(A_i')_{\mathfrak{q}_i'}$-module $(\hat{\Omega}^1_{A_i'/k'})_{\mathfrak{q}_i'}$ is \emph{at most equal to $m$}; whence our assertion.

Suppose therefore that $k$ is algebraically closed, and let us show that we may further reduce to the case $\mathfrak q=0$. Indeed, $\Omega^1_{(A/\mathfrak q)/A}=0$ \sref[0]{0.20.4.12}, so the exact sequence \sref{IV.18.11.3.6} gives a surjection
\[
  (\hat{\Omega}^1_{A/k})\otimes_{A}(A/\mathfrak q)
  \longrightarrow \hat{\Omega}^1_{(A/\mathfrak q)/k}.
\]
It follows at once that the minimum number of generators of the $K$-vector space $(\hat{\Omega}^1_{(A/\mathfrak q)/k})\otimes_{A/\mathfrak q}K$, where $K$ is the fraction field of $A/\mathfrak q$, is at most $m$.

We may thus assume that $A$ is integral. Its fraction field $K$ is separable over the algebraically closed field $k$, and hence geometrically regular over $k$ \sref{IV.6.7.6}. We may therefore apply \sref{IV.18.11.3} with $\mathfrak p=0$; since $k=k^p$, this gives $m=n$.
\end{proof}

\begin{lemma}[18.11.5.2]
\label{IV.18.11.5.2}
Let $k$ be a field, $A$ a $k$-algebra, local Noetherian and complete, whose residue field is a finite extension of $k$, $\mathfrak{m}$ its maximal ideal, $\mathfrak{p}\subset\mathfrak{m}$ a prime ideal of $A$.
Let $m$ be the minimum number of generators of the $A_\mathfrak{p}$-module $(\hat{\Omega}^1_{A/k})_\mathfrak{p}$, and let $(x_i)_{1\leq i\leq m}$ be a family of elements of $\mathfrak{m}$ not belonging to $\mathfrak{p}$.
Then there exist $m$ invertible elements $u_i$ \textup{(}$1\leq i\leq m$\textup{)} of $A$ such that if $y_i=u_ix_i$, the canonical images of the $dy_i$ in $(\hat{\Omega}^1_{A/k})_\mathfrak{p}$ generate this $A_\mathfrak{p}$-module.
\end{lemma}

\begin{proof}
For all $x\in A$, denote by $\delta(x)$ the canonical image of $dx$ $(=d_{A/k}x)$ in $(\hat{\Omega}^1_{A/k})_\mathfrak{p}\otimes_{A_\mathfrak{p}}\boldsymbol{k}(\mathfrak{p})=\hat{\Omega}^1_{A/k}\otimes_{A}\boldsymbol{k}(\mathfrak{p})$; as this is a $\boldsymbol{k}(\mathfrak{p})$-vector space of dimension $m$ by hypothesis, it suffices (by virtue of the Nakayama lemma) to prove that one can determine the $u_i$ such that the $\delta(u_ix_i)$ form a free system.
We reason by induction and suppose for an integer $r<m$, we have determined the $u_i$ ($1\leq i\leq r$) such that the $\delta(u_ix_i)$ for $1\leq i\leq r$ are linearly independent over $\boldsymbol{k}(\mathfrak{p})$; if $\delta(x_{r+1})$ is not a linear combination of the $\delta(u_ix_i)$ for $1\leq i\leq r$, one takes $u_{r+1}=1$ to continue the induction.
In the contrary case, for $u\in A$ the element $\delta(ux_{r+1})$ is the canonical image of $(du)x_{r+1}+u(dx_{r+1})$. The second term has image in the span of the $\delta(u_ix_i)$ for $1\leq i\leq r$, while the image of $x_{r+1}$ in $\boldsymbol{k}(\mathfrak p)$ is nonzero. It is therefore enough to find a unit $u\in A$ such that $\delta(u)$ does not belong to that span.
\oldpage[IV-4]{176}
Now \sref[0]{0.20.7.15} and \sref[0\textsubscript{\mdseries I}]{0.7.2.9} show that the $\delta(x)$, for $x\in A$, generate the $\boldsymbol{k}(\mathfrak p)$-vector space $(\hat{\Omega}^1_{A/k})_{\mathfrak p}\otimes_{A_{\mathfrak p}}\boldsymbol{k}(\mathfrak p)$. Since $r<m$, there is therefore a $z\in A$ for which $\delta(z)$ is not in the span of the $\delta(u_ix_i)$. If $z\notin\mathfrak m$, take $u=z$; otherwise $1+z$ is invertible and $\delta(1+z)=\delta(z)$ because $d(1+z)=dz$, so take $u=1+z$. This completes the proof of \sref{IV.18.11.5.2}.
\end{proof}

\begin{proof}[Proof of {\sref{IV.18.11.5}}, continued]
We now return to the proof of the implication \emph{a)} $\Rightarrow$ \emph{b)} in \sref{IV.18.11.5}.
Let us first note that since $\mathfrak{p}\neq\mathfrak{m}$, there exists a system of parameters $(x_i)_{1\leq i\leq n}$ of $A$ (with $n=\dim(A)$) such that $x_i\notin\mathfrak{p}$ for all $i$, without which $A/\mathfrak{p}$ would be of finite length, and $\mathfrak{p}$ would be maximal, contrary to the hypothesis.
But if for example $x_1\notin\mathfrak{p}$, it suffices, for each index $i$ such that $x_i\in\mathfrak{p}$, to replace $x_i$ by $x_1+x_i$, to have a system of parameters none of whose elements belong to $\mathfrak{p}$.
The hypothesis \emph{a)} and the relation $\dim(A/\mathfrak{q})=n$ imply, by virtue of \sref{IV.18.11.5.1}, that $m\geq n$; one can therefore consider a family $(x_i)_{1\leq i\leq m}$ of elements $x_i\notin\mathfrak{p}$, whose $n$ first form a system of parameters of $A$.
Multiplying furthermore the $x_i$ by invertible elements $u_i$ of $A$, one can, by \sref{IV.18.11.5.2}, suppose that the images of $dx_i$ ($1\leq i\leq m$) in $(\hat{\Omega}^1_{A/k})_\mathfrak{p}$ generate this $A_\mathfrak{p}$-module, and the multiplication by the $u_i$ has not altered the fact that the $x_i$ for $1\leq i\leq n$ form a system of parameters.
We shall consider the local $k$-homomorphism $u:B\to A$ such that $u(T_i)=x_i$ for $1\leq i\leq m$ (Bourbaki, \emph{Alg.\ comm.}, chap.~III, \textsection~4, n\textsuperscript{o}~5, prop.~6); as the $x_i$ generate an ideal of definition of $A$, it follows from \sref[0\textsubscript{\mdseries I}]{0.7.4.4} and \sref{IV.7.4.3} that $u$ makes $A$ a $B$-algebra \emph{finite}.
We therefore have the exact sequence \sref{IV.18.11.3.5}
\[
  (\hat{\Omega}^1_{B/k})\otimes_BA\xrightarrow{v}\hat{\Omega}^1_{A/k}\xrightarrow{w}\Omega^1_{A/B}\longrightarrow 0
\]

But the $d_Ax_i$ are the canonical images by $v$ of the elements $d_BT_i\otimes 1$ \sref[0]{0.20.5.2.6}.
Localising the preceding exact sequence at $\mathfrak p$, we see that
\[
  v_\mathfrak p:(\hat{\Omega}^1_{B/k})\otimes_BA_\mathfrak p
  \longrightarrow(\hat{\Omega}^1_{A/k})_\mathfrak p
\]
is \emph{surjective}. Consequently
\[
  0=(\Omega^1_{A/B})_\mathfrak p
   =\Omega^1_{A_\mathfrak p/B_\mathfrak r},
\]
where $\mathfrak r=u^{-1}(\mathfrak p)$ is the inverse image of $\mathfrak p$ in $B$; cf.\ \sref{IV.16.4.15}. By \sref{IV.17.4.1}, the morphism $\operatorname{Spec}(A)\to\operatorname{Spec}(B)$ is therefore unramified at $\mathfrak p$, which proves $b)$.
\end{proof}

\begin{remark}[18.11.6]
\label{IV.18.11.6}
In the statement of \sref{IV.18.11.5}, the hypothesis $\mathfrak{p}\neq\mathfrak{m}$ cannot be omitted.
Indeed, for every local $k$-homomorphism $u:B\to A$, where $B=k[[T_1,\ldots,T_r]]$ ($r$ any integer) making $A$ a $B$-algebra finite, $\mathfrak{m}$ is the only point of $\operatorname{Spec}(A)$ above the maximal ideal $\mathfrak{n}$ of $B$; the morphism $\operatorname{Spec}(A)\to\operatorname{Spec}(B)$ cannot be unramified at $\mathfrak{m}$ unless $A/\mathfrak{n}A$ is a field, an extension separable of $k$ \sref{IV.17.4.1}, which implies that the residue field $K$ of $A$ is a separable extension of $k$.
If this condition is not satisfied, the conclusion of \sref{IV.18.11.5} can never be satisfied for $\mathfrak{p}=\mathfrak{m}$, whatever the integer $m$.
\end{remark}

\begin{corollary}[18.11.7]
\label{IV.18.11.7}
Let $k$ be a field, $A$ a $k$-algebra, local Noetherian and complete, equidimensional, whose residue field is a finite extension of $k$, $\mathfrak{p}$ a prime ideal of $A$.
Then the minimum number of generators of the $A_\mathfrak{p}$-module $(\hat{\Omega}^1_{A/k})_\mathfrak{p}$ is at least equal to $\dim(A)$.
\end{corollary}

\begin{proof}
This is a special case of \sref{IV.18.11.5.1}.
\end{proof}

\oldpage[IV-4]{177}

\begin{corollary}[18.11.8]
\label{IV.18.11.8}
Let $k$ be a field, $A$ a $k$-algebra, local Noetherian and complete, \textup{integral}, whose residue field is a finite extension of $k$.
If $K$ is the field of fractions of $A$, one has $\operatorname{rg}_{K}((\hat{\Omega}^1_{A/k})\otimes_{A}K)\geq\dim(A)$.
\end{corollary}

\begin{proof}
It suffices to set $\mathfrak{p}=0$ in \sref{IV.18.11.7}.
\end{proof}

\begin{corollary}[18.11.9]
\label{IV.18.11.9}
Let $k$ be a field, $A$ a $k$-algebra, local Noetherian and complete, of dimension $n$, whose residue field is a finite extension of $k$.
Let $\mathfrak{p}$ be a prime ideal of $A$ distinct from the maximal ideal, and containing a minimal prime ideal $\mathfrak{q}$ such that $\dim(A/\mathfrak{q})=n$.
Suppose that the $A_\mathfrak{p}$-module $(\hat{\Omega}^1_{A/k})_\mathfrak{p}$ admits $n$ generators.
Then:
\begin{enumerate}[label=\textup{(\roman*)}]
  \item The $A_\mathfrak{p}$-module $(\hat{\Omega}^1_{A/k})_\mathfrak{p}$ is free of rank $n$.
  \item There exists a local $k$-homomorphism $u:B\to A$, where $B=k[[T_1,\ldots,T_n]]$, making $A$ a $B$-algebra finite, and such that the corresponding morphism $\operatorname{Spec}(A)\to\operatorname{Spec}(B)$ is \'etale at the point $\mathfrak{p}$.
  \item The $k$-algebra $A_\mathfrak{p}$ is geometrically regular over $k$.
\end{enumerate}
\end{corollary}

\begin{proof}
We first prove (ii). By \sref{IV.18.11.5}, there is a local $k$-homomorphism
\[
  u:B=k[[T_1,\ldots,T_n]]\longrightarrow A
\]
making $A$ a \emph{finite} $B$-algebra and such that $\operatorname{Spec}(A)\to\operatorname{Spec}(B)$ is unramified at $\mathfrak p$; put $\mathfrak r=u^{-1}(\mathfrak p)$. The hypothesis on $\mathfrak q\subset\mathfrak p$, together with the fact that $A$ is a quotient of a regular ring \sref[0]{0.19.8.8}, gives
\[
  \dim(A_\mathfrak p)=n\setmin\dim(A/\mathfrak p)
\]
by \sref[0]{0.16.5.12}; likewise $\dim(B_\mathfrak r)=n\setmin\dim(B/\mathfrak r)$. Since the morphism is unramified at $\mathfrak p$, its fibre over $\mathfrak r$ has dimension zero, so \sref[0]{0.16.3.9} gives $\dim(A/\mathfrak p)\leq\dim(B/\mathfrak r)$ and hence $\dim(B_\mathfrak r)\leq\dim(A_\mathfrak p)$. It follows from \sref{IV.18.10.1} that the morphism is \'etale at $\mathfrak p$.

Assertion (iii) follows because $A_\mathfrak p$ is formally smooth over $B_\mathfrak r$, while $B_\mathfrak r$ is formally smooth over $k$, for their preadic topologies \sref[0]{0.19.3.4} and \sref[0]{0.19.3.5}. Thus $A_\mathfrak p$ is formally smooth over $k$ for its preadic topology, and consequently geometrically regular over $k$ \sref[0]{0.22.5.8}.

Finally, we prove (i). Let $k'$ be an algebraically closed extension of $k$, and consider the semilocal $k'$-algebra $A'=A\widehat{\otimes}_kk'$. Regarding $A$ as a finite $B$-module via $u$, there is a canonical isomorphism
\[
  A'\simeq A\otimes_{B}(B\widehat{\otimes}_kk')
\]
by \sref{IV.7.5.7.1}; that result does not require the residue field of $B$ to be finite over $k$. Put $B'=B\widehat{\otimes}_kk'$, which identifies canonically with $k'[[T_1,\ldots,T_n]]$.

Since $\operatorname{Spec}(A)\to\operatorname{Spec}(B)$ is finite and \'etale at $\mathfrak p$, the morphism $\operatorname{Spec}(A')\to\operatorname{Spec}(B')$ is finite and \'etale at every point $\mathfrak p'$ over $\mathfrak p$ \sref{IV.17.3.3}. Moreover, $A'$ is the direct product of local rings $A_i'$ of dimension $n$ \sref{IV.18.11.4}, and $\mathfrak p'$ identifies with a prime ideal $\mathfrak p_i'$ of one of them. The preceding argument shows that $(A_i')_{\mathfrak p_i'}$ is geometrically regular over $k'$. Since $k'$ is perfect, \sref{IV.18.11.3} shows that $(\hat{\Omega}^1_{A_i'/k'})_{\mathfrak p_i'}$ is a free $(A_i')_{\mathfrak p_i'}$-module of finite type.

But \sref{IV.18.11.4} gives
\[
  (\hat{\Omega}^1_{A_i'/k'})_{\mathfrak p_i'}
  \simeq (\hat{\Omega}^1_{A/k})_\mathfrak p
  \otimes_{A_\mathfrak p}(A_i')_{\mathfrak p_i'},
\]
and $(A_i')_{\mathfrak p_i'}$ is faithfully flat over $A_\mathfrak p$. Therefore $(\hat{\Omega}^1_{A/k})_\mathfrak p$ is a \emph{free} $A_\mathfrak p$-module of finite type \sref{IV.2.5.2};

\oldpage[IV-4]{178}

and \sref{IV.18.11.5.1}, together with $\dim(A/\mathfrak q)=n$, shows that it is \emph{of rank $n$}.
\end{proof}

\begin{theorem}[18.11.10]
\label{IV.18.11.10}
Let $k$ be a field of characteristic exponent $p$, $A$ a $k$-algebra, local Noetherian and complete, whose residue field is a finite extension of $k$, $\mathfrak{p}$ a prime ideal of $A$ distinct from the maximal ideal.
The following conditions are equivalent:
\begin{enumerate}[label=\emph{\alph*)}]
  \item For every extension $k'$ of $k$, and every prime ideal $\mathfrak{p}'$ of $A'=A\widehat{\otimes}_kk'$ above $\mathfrak{p}$, $A'_{\mathfrak{p}'}$ is a regular ring.
\item[\textup{a$'$)}] There exists a \textup{perfect} extension $k'$ of $k$ and a prime ideal $\mathfrak{p}'$ of $A'=A\widehat{\otimes}_kk'$ above $\mathfrak{p}$ such that $A'_{\mathfrak{p}'}$ is regular.
  \item Let $n$ be the largest of the dimensions of the irreducible components of $\operatorname{Spec}(A)$ to which $\mathfrak{p}$ belongs; then $(\hat{\Omega}^1_{A/k})_\mathfrak{p}$ is an $A_\mathfrak{p}$-module free of rank $n$.
\end{enumerate}
Suppose furthermore that $n=\dim(A)$ \textup{(which will be the case if $A$ is \emph{equidimensional})}; then the preceding conditions are also equivalent to:
\begin{enumerate}[label=\emph{\alph*)}]
  \setcounter{enumi}{2}
  \item There exists a local $k$-homomorphism $u:B\to A$, where $B=k[[T_1,\ldots,T_n]]$, making $A$ a $B$-algebra finite, and such that the corresponding morphism $\operatorname{Spec}(A)\to\operatorname{Spec}(B)$ is \'etale at the point $\mathfrak{p}$.
\end{enumerate}
Each of the conditions \textup{a)}, \textup{a$'$)}, and \textup{b)} implies the following:
\begin{enumerate}[label=\emph{\alph*)}]
  \setcounter{enumi}{3}
  \item The ring $A_\mathfrak{p}$ is geometrically regular over $k$.
\end{enumerate}
Furthermore, if $[k:k^p]<+\infty$, condition \textup{d)} is equivalent to \textup{a)}, \textup{a$'$)}, and \textup{b)}.
\end{theorem}

\begin{proof}
That $d)$ implies $b)$ when $[k:k^p]<+\infty$ is precisely \sref{IV.18.11.3}. It is clear that $a)$ implies $a')$. Let us show that $c)$ implies $a)$ when $\dim(A)=n$. With the notation of $a)$, one then has
\[
  A'=A\otimes_{B}B',
  \qquad
  B'=B\widehat{\otimes}_kk'=k'[[T_1,\ldots,T_n]]
\]
by \sref{IV.7.5.7.1}. The morphism $\operatorname{Spec}(A')\to\operatorname{Spec}(B')$ is finite and \'etale at $\mathfrak p'$, and the argument proving \sref{IV.18.11.9},~(iii), shows that $A'_{\mathfrak p'}$ is regular.

That $b)$ implies $c)$ when $\dim(A)=n$ follows from \sref{IV.18.11.9}. Let us prove that $a')$ implies $b)$ under the same hypothesis. The ring $A'$ is flat over $A$ \sref{IV.18.11.4}, and $\dim(A')=\dim(A)=n$. By \sref{IV.2.3.4} and \sref{IV.6.1.1}, there is a minimal prime $\mathfrak q'$ of $A'$, contained in $\mathfrak p'$ and lying over a minimal prime $\mathfrak q$ of $A$, such that
\[
  \dim(A'/\mathfrak q')=\dim(A/\mathfrak q)=n.
\]
Moreover, \sref{IV.18.11.4} gives
\[
  (\hat{\Omega}^1_{A'/k'})_{\mathfrak p'}
  \simeq(\hat{\Omega}^1_{A/k})_\mathfrak p
  \otimes_{A_\mathfrak p}A'_{\mathfrak p'}.
\]
Since $k'$ is perfect, the regularity of $A'_{\mathfrak p'}$ implies that it is \emph{geometrically regular} over $k'$ \sref{IV.6.7.7}. As $k'^p=k'$, we may apply the implication $d)\Rightarrow b)$ to $A'$ and $\mathfrak p'$. Thus $(\hat{\Omega}^1_{A'/k'})_{\mathfrak p'}$ is a free $A'_{\mathfrak p'}$-module of rank $n$; faithful flatness \sref{IV.18.11.4} and \sref{IV.2.5.2} then shows that $(\hat{\Omega}^1_{A/k})_\mathfrak p$ is a free $A_\mathfrak p$-module of rank $n$.

This proves the equivalence of $a)$, $a')$, $b)$, and $c)$ when $\dim(A)=n$.
It remains to prove that $a)$, $a')$, and $b)$ are still equivalent in the general case.
Let $\mathfrak J$ be the kernel of the canonical homomorphism $A\to A_\mathfrak p$, and put $A_1=A/\mathfrak J$. Necessarily $\mathfrak J\subset\mathfrak p$; and if $\mathfrak p_1=\mathfrak p/\mathfrak J$, the canonical homomorphism $A_\mathfrak p\to(A_1)_{\mathfrak p_1}$ is \emph{bijective}. It follows from \sref[I]{I.6.5.4} that the canonical injection $\operatorname{Spec}(A_1)\to\operatorname{Spec}(A)$ is a \emph{local isomorphism} at $\mathfrak p_1$, and $\mathfrak J_\mathfrak p=0$. As in \sref{IV.18.11.3}, there is an exact sequence
\[
  \mathfrak J/\mathfrak J^2
  \longrightarrow (\hat{\Omega}^1_{A/k})\otimes_{A}A_1
  \longrightarrow \hat{\Omega}^1_{A_1/k}\longrightarrow0
\]\oldpage[IV-4]{179}
After localising at $\mathfrak p$, this yields an isomorphism
\[
  (\hat{\Omega}^1_{A/k})_\mathfrak p
  \xrightarrow{\sim}
  (\hat{\Omega}^1_{A_1/k})_{\mathfrak p_1}.
\]
Thus condition $b)$ for $A$ and $\mathfrak p$ is \emph{equivalent} to condition $b)$ for $A_1$ and $\mathfrak p_1$.

On the other hand, with the notation of $a)$, there is a canonical isomorphism
\[
  A_1'=A_1\widehat{\otimes}_kk'
  \simeq A'/\mathfrak JA'
\]
by \sref{IV.7.5.7.1}, where the hypothesis on the residue field of $B$ is superfluous. If $\mathfrak p'$ is a prime of $A'$ above $\mathfrak p$, every element of $\mathfrak JA'$ annihilates an element of $A'\setmin\mathfrak p'$, and hence $\mathfrak JA'\subset\mathfrak p'$. Putting $\mathfrak p_1'=\mathfrak p'/\mathfrak JA'$, the prime $\mathfrak p_1'$ lies above $\mathfrak p_1$ and $(A_1')_{\mathfrak p_1'}$ identifies canonically with $A'_{\mathfrak p'}$. This proves that condition $a)$ (respectively $a')$) for $A$ and $\mathfrak p$ is \emph{equivalent} to the same condition for $A_1$ and $\mathfrak p_1$.

All minimal prime ideals of $A$ contained in $\mathfrak p$ contain $\mathfrak J$, since $\operatorname{Spec}(A_1)\to\operatorname{Spec}(A)$ is a local isomorphism at $\mathfrak p_1$. Moreover, the primes in $\operatorname{Ass}_{A}(A/\mathfrak J)$ are precisely those in $\operatorname{Ass}(A)$ which are contained in $\mathfrak p$ (Bourbaki, \emph{Alg.\ comm.}, chap.~IV, \textsection~1, n\textsuperscript{o}~2, prop.~6). Hence all minimal primes of $A_1$ are contained in $\mathfrak p_1$, and consequently $\dim(A_1)=n$. It remains only to apply to $A_1$ and $\mathfrak p_1$ what was proved above.
\end{proof}

\begin{remarks}[18.11.11]
\label{IV.18.11.11}
\begin{enumerate}[label=(\roman*)]
  \item The equivalence of conditions $d)$ and $b)$ in \sref{IV.18.11.10} is no longer valid when one no longer supposes that $[k:k^p]<+\infty$.
Indeed, in the example of \sref[0]{0.22.7.7},~(ii), the ring $B=A/\mathfrak{q}$ is integral and of dimension $1$; on the other hand, the sequence
\[
  (\mathfrak{q}/\mathfrak{q}^2)\otimes_{A}L\xrightarrow{j}\hat{\Omega}^1_{A/k}\otimes_{A}L\to\hat{\Omega}^1_{B/k}\otimes_{B}L\to 0
\]
is exact (same demonstration as in \sref{IV.18.11.3}), and we know that \emph{$j$ is not injective}, hence $j=0$ since $(\mathfrak{q}/\mathfrak{q}^2)\otimes_AL$ is of rank $1$.
We conclude that $\operatorname{rg}_{L}(\hat{\Omega}^1_{B/k}\otimes_{B}L)=2$; as $B_\mathfrak{p}$ is a $k$-algebra geometrically regular, one sees that the condition $d)$ does not imply here the condition $b)$.
  \item The notation being that of \sref{IV.18.11.10}, suppose that $n=\dim(A)$ and let $(x_i)_{1\leq i\leq n}$ be a system of parameters of $A$ not belonging to $\mathfrak{p}$; one deduces a local $k$-homomorphism $u:B\to A$ such that $u(T_i)=x_i$ for $1\leq i\leq n$, making $A$ a finite $B$-algebra.
For the corresponding morphism $\operatorname{Spec}(A)\to\operatorname{Spec}(B)$ to be \emph{\'etale} at the point $\mathfrak{p}$, it is necessary and sufficient that the images of $d_Ax_i$ in $(\hat{\Omega}^1_{A/k})_\mathfrak{p}$ form a \emph{system of generators} of this $A_\mathfrak{p}$-module.
Indeed, one has seen in the proof of \sref{IV.18.11.5} that if the morphism $\operatorname{Spec}(A)\to\operatorname{Spec}(B)$ is not ramified at $\mathfrak{p}$, the canonical homomorphism $(\hat{\Omega}^1_{B/k}\otimes_BA)_\mathfrak{p}\to(\hat{\Omega}^1_{A/k})_\mathfrak{p}$ is surjective and the images of the elements $d_BT_i\otimes 1$, which generate $(\hat{\Omega}^1_{B/k})\otimes_BA$, are the $d_Ax_i$; whence the \emph{necessity} of the condition.
Conversely, the same reasoning shows that if this condition is satisfied, the morphism $\operatorname{Spec}(A)\to\operatorname{Spec}(B)$ is not ramified at $\mathfrak{p}$, and the reasoning of \sref{IV.18.11.9} proves in fact that this morphism is \emph{\'etale} at $\mathfrak{p}$.
\end{enumerate}
\end{remarks}

\begin{corollary}[18.11.12]
\label{IV.18.11.12}
Let $k$ be a field of characteristic exponent $p$ such that $[k:k^p]<+\infty$, $A$ a $k$-algebra, local Noetherian and complete and integral, which is not a field, and whose residue field is a finite extension of $k$.
Then there exists a finite radicial extension $k'$ of $k$ such that, setting $A'=A\otimes_kk'$, the $k'$-algebra $A'_{\mathrm{red}}$ and the prime ideal $0$ of this algebra satisfy the equivalent conditions \textup{a)}, \textup{a$'$)}, \textup{b)}, \textup{c)}, and \textup{d)} of \sref{IV.18.11.10}.
In particular, if $n=\dim(A)=\dim(A')$,
\oldpage[IV-4]{180}
there exists a local $k'$-homomorphism $B'=k'[[T_1,\ldots,T_n]]\to A'_\mathrm{red}$ making $A'_\mathrm{red}$ a finite $B'$-algebra, and such that the field of fractions $K'$ of $A'_\mathrm{red}$ is a finite \textup{separable} extension of the field of fractions $k'((T_1,\ldots,T_n))$ of $B'$.
\end{corollary}

\begin{proof}
Since $\operatorname{Spec}(A')\to\operatorname{Spec}(A)$ is radicial and finite, $A'$ is a complete local ring and its nilradical is the unique prime above the zero ideal of $A$. By flatness, $A'$ identifies with a subring of $K\otimes_kk'$, which is also the total ring of fractions of $A'$; hence
\[
  K'=(K\otimes_kk')_{\mathrm{red}}.
\]
In view of the equivalence \emph{d)}$\Leftrightarrow$\emph{c)} in \sref{IV.18.11.10}, it is therefore enough to find a finite radicial extension $k'$ of $k$ such that $(K\otimes_kk')_{\mathrm{red}}$ is \emph{separable} over $k'$ \sref{IV.6.7.6}.

By \sref[0]{0.19.8.9}, under the present hypotheses there is a sub-$k$-algebra $C=k[[T_1,\ldots,T_m]]$ of $A$ such that $A$ is finite over $C$. Thus $K$ is a \emph{finite} extension of the fraction field $K_1=k((T_1,\ldots,T_m))$ of $C$, and $K_1$ is \emph{separable} over $k$ \sref[0]{0.21.9.6.4}. We may write
\[
  K\otimes_kk^{p^{-\infty}}
  =K\otimes_{K_1}(K_1\otimes_kk^{p^{-\infty}}),
\]
and $K_1\otimes_kk^{p^{-\infty}}$ is a \emph{field}, radicial over $K_1$. Consequently $K\otimes_kk^{p^{-\infty}}$ is finite over a field and is therefore \emph{Artinian}. The conclusion follows from the more general lemma below.
\end{proof}

\begin{lemma}[18.11.12.1]
\label{IV.18.11.12.1}
Let $k$ be a field of characteristic $p>0$, $K$ an extension of $k$.
The following conditions are equivalent:
\begin{enumerate}[label=\emph{\alph*)}]
  \item The ring $K\otimes_kk^{p^{-\infty}}$ is Artinian.
  \item There exists an extension $k'$ of $k$ of finite type such that $(K\otimes_kk')_{\mathrm{red}}$ is a $k'$-algebra separable.
  \item There exists a finite radicial extension $k'$ of $k$ such that $(K\otimes_kk')_\mathrm{red}$ is a separable extension of $k'$.
\end{enumerate}
\end{lemma}

\begin{proof}
Let us first show that $c)$ implies $a)$.
The ring $A=K\otimes_kk'$ is then local Artinian. If $\mathfrak N$ is its nilradical, its residue field $L=A/\mathfrak N$ is separable over $k'$, and hence $L\otimes_{k'}k^{p^{-\infty}}$ is a field. Since this field is
\[
  (A\otimes_{k'}k^{p^{-\infty}})/(\mathfrak N\otimes_{k'}k^{p^{-\infty}}),
\]
the ideal $\mathfrak N\otimes_{k'}k^{p^{-\infty}}$ is the nilradical of $K\otimes_kk^{p^{-\infty}}$. The ideal $\mathfrak N$ is finitely generated and nilpotent, so this nilradical is finitely generated; it follows that $K\otimes_kk^{p^{-\infty}}$ is Artinian.

Conversely, let us prove that $a)$ implies $c)$.
Let $\mathfrak R$ be the nilradical of the local Artinian ring $B=K\otimes_kk^{p^{-\infty}}$. It is generated by finitely many elements
\[
  y_i=\sum_j x_{ij}\otimes\xi_{ij},
  \qquad x_{ij}\in K,\quad \xi_{ij}\in k^{p^{-\infty}}.
\]
Let $k'$ be the finite radicial extension of $k$ generated by the $\xi_{ij}$, and let $\mathfrak R_0$ be the ideal of $B_0=K\otimes_kk'$ generated by the $y_i$. The $y_i$ are nilpotent in $B_0$; moreover
\[
  \mathfrak R_0\otimes_{k'}k^{p^{-\infty}}=\mathfrak R,
  \qquad
  \mathfrak R\cap B_0=\mathfrak R_0.
\]
Thus $\mathfrak R_0$ is exactly the nilradical of $B_0$. Since
\[
  B/\mathfrak R=(B_0/\mathfrak R_0)\otimes_{k'}k^{p^{-\infty}}
\]
is reduced, \sref{IV.4.6.1} shows that $B_0/\mathfrak R_0=(K\otimes_kk')_{\mathrm{red}}$ is separable over $k'$.

It is clear that $c)$ implies $b)$.
Conversely, suppose $b)$ holds. There is a separable extension $k_1$ of $k$ such that $k'$ is finite radicial over $k_1$; put $K_1=K\otimes_kk_1$, which is a field. Applying the equivalence of $a)$ and $c)$ to the extension $K_1/k_1$, we see that $K_1\otimes_{k_1}k_1^{p^{-\infty}}$ is Artinian. But
\oldpage[IV-4]{181}
\[
  K\otimes_kk_1^{p^{-\infty}}
  =(K\otimes_kk^{p^{-\infty}})
   \otimes_{k^{p^{-\infty}}}k_1^{p^{-\infty}},
\]
so $K\otimes_kk^{p^{-\infty}}$ is Artinian as well (Bourbaki, \emph{Alg.\ comm.}, chap.~I, \textsection~3, n\textsuperscript{o}~5, cor.\ of prop.~8). Thus $b)$ implies $a)$, completing the proofs of \sref{IV.18.11.12.1} and \sref{IV.18.11.12}.
\end{proof}


\subsection{Applications of \'etale localisation to quasi-finite morphisms (generalisations of earlier results)}
\label{subsection:IV.18.12}

The results of this section were communicated by P.~Deligne.

\begin{theorem}[18.12.1]
\label{IV.18.12.1}
Let $f:X\to Y$ be a morphism locally of finite type, let $x$ be a point of $X$, and put $y=f(x)$.
Suppose that $x$ is an isolated point of the space $f^{-1}(y)$.
Then there exist an \'etale morphism $Y'\to Y$, a point $x'$ of $X'=X\times_{Y}Y'$ above $x$, and an open neighbourhood $V'$ of $x'$ in $X'$ such that, if $f'=f_{(Y')}:X'\to Y'$, the restriction $f'|V'$ is finite.
If, moreover, $f$ is separated, then $V'$ is both open and closed, and $X'$ is therefore the disjoint union of two subschemes induced on open subsets of $X'$, one of which is finite over $Y'$ and contains $x'$.
\end{theorem}

\begin{proof}
The last assertion follows because, if $f$ is separated, so is $f'$. Thus, if $j':V'\to X'$ denotes the canonical injection, the finiteness of $f'|V'=f'\circ j'$ implies that $j'$ is finite \sref[II]{II.6.1.5}; consequently $V'=j'(V')$ is closed in $X'$ \sref[II]{II.6.1.10}.

The question is local on $X$, so we may suppose that $Y=\operatorname{Spec}(A)$ and $X=\operatorname{Spec}(B)$ are affine, with $B$ an $A$-algebra of finite type. Write $B=C/\mathfrak{J}$, where $C=A[T_1,\ldots,T_r]$ and $\mathfrak{J}$ is an ideal of $C$. Since $f$ is then \emph{separated}, we may also suppose that $f^{-1}(y)=\{x\}$.

Let $(\mathfrak{J}_\lambda)$ be the filtered family of finitely generated ideals of $C$ contained in $\mathfrak{J}$; then $\mathfrak{J}$ is their union. Regard $X$ and the schemes $X_\lambda=\operatorname{Spec}(C/\mathfrak{J}_\lambda)$ as closed subschemes of $Z=\operatorname{Spec}(C)$. On the underlying spaces, $X=\bigcap_\lambda X_\lambda$. If $f_\lambda:X_\lambda\to Y$ is the structural morphism, then
\[
  f^{-1}(y)=\bigcap_\lambda f_\lambda^{-1}(y).
\]
The sets $f_\lambda^{-1}(y)$ are closed in the \emph{Noetherian} space $Z_y=\operatorname{Spec}(\boldsymbol{k}(y)[T_1,\ldots,T_r])$; hence, for some $\lambda$, $f_\lambda^{-1}(y)=f^{-1}(y)=\{x\}$.

We may replace $X$ by such an $X_\lambda$. Indeed, if the assertion is proved for $X_\lambda$, let $x'$ lie above $x$ in $Z'=Z\times_{Y}Y'$ and put $X_\lambda'=p^{-1}(X_\lambda)$, where $p:Z'\to Z$ is the projection (so also $X'=p^{-1}(X)$). If $V_\lambda'$ is an open neighbourhood of $x'$ in $X_\lambda'$ finite over $Y'$, then the restriction to the closed subscheme $X'\cap V_\lambda'$ is finite as well. We have therefore reduced to the case where $f$ is of \emph{finite presentation}.

The set of points of $X$ which are isolated in their fibres is open in $X$ \sref{IV.13.1.4}; after replacing $X$ by an open neighbourhood of $x$, we may therefore suppose that $f$ is \emph{quasi-finite}. Let $A''$ be the Henselization of $\sh{O}_{Y,y}$ and put $Y''=\operatorname{Spec}(A'')$. If $X''=X\times_{Y}Y''$ and $f''=f_{(Y'')}:X''\to Y''$, then $f''$ is separated, quasi-finite, and of finite presentation. Since $A''$ is henselian, every point $x''\in X''$ above $x$ has an open-and-closed neighbourhood $V''$ in $X''$ which is \emph{finite} over $Y''$ \sref{IV.18.5.11},~\emph{c)}. The immersion $V''\to X''$ is open and closed, hence quasi-compact and of finite presentation \sref{IV.1.6.2}; consequently $f''|V''$ is of finite presentation.

By definition, $A''$ is the filtered direct limit of strictly essentially \'etale $\sh{O}_{Y,y}$-algebras $A_\lambda$ \sref{IV.18.6.5}. Each $A_\lambda$ is itself the filtered direct limit of \emph{\'etale} $\sh{O}_{Y,y}$-algebras $B_{\lambda\mu}$ \sref{IV.18.6.1}, \sref{IV.8.1.2},~$a$). Since the $B_{\lambda\mu}$ may be assumed finitely presented over $\sh{O}_{Y,y}$, it follows again from \sref{IV.8.1.2},~\emph{a)} and \sref{IV.17.7.8} that they are filtered direct limits of \'etale $A$-algebras. The double direct-limit theorem therefore shows that $Y''$ is the projective limit of a filtered family $(Y_\alpha')$ of affine schemes \'etale over $Y$.

Let $x_\alpha'$ be the projection of $x''$ in $X_\alpha'=X\times_{Y}Y_\alpha'$. Since $V''$ is quasi-compact, after increasing $\alpha$ we may write
\[
  V''=V_\alpha'\times_{Y_\alpha'}Y'',
\]
where $V_\alpha'$ is an open neighbourhood of $x_\alpha'$ in $X_\alpha'$ \sref{IV.8.2.11}. Finally, since $V''$ is of finite presentation over $Y''$, \sref{IV.8.10.5},~(x) shows that, for a suitable $\alpha$, $V_\alpha'$ is \emph{finite} over $Y_\alpha'$, as required.
\end{proof}

\begin{remark}[18.12.2]
\label{IV.18.12.2}
The preceding proof shows (taking \sref{IV.18.6.2} into account) that $Y'$ may moreover be chosen so that, if $y'=f'(x')$, the homomorphism $\boldsymbol{k}(y)\to\boldsymbol{k}(y')$ is \emph{bijective}.
\end{remark}

\begin{corollary}[18.12.3]
\label{IV.18.12.3}
Let $f:X\to Y$ be a morphism locally of finite type and separated, $y$ a point of $Y$ such that the subspace $f^{-1}(y)$ is finite and discrete.
Then there exist an \'etale morphism $Y'\to Y$, a point $y'\in Y'$ above $y$, such that $\boldsymbol{k}(y')=\boldsymbol{k}(y)$, and a decomposition of $X'$ as the sum of two sub-preschemes $X'_1$, $X'_2$ induced on open subsets of $X'$, such that the restriction of $f'=f_{(Y')}:X'\to Y'$ to $X'_1$ is a finite morphism and that $X'_2\cap f'^{-1}(y')=\varnothing$.
\end{corollary}

\begin{proof}
If $n$ is the number of points of $f^{-1}(y)$, one argues by induction on $n$, the corollary being trivial for $n=0$.
Let $x$ be a point of $f^{-1}(y)$; by virtue of \sref{IV.18.12.1} and \sref{IV.18.12.2} there are an \'etale morphism $Y_1\to Y$ and a point $y_1$ of $Y_1$ above $y$ such that $\boldsymbol{k}(y_1)=\boldsymbol{k}(y)$ and, if one sets $S_1=X\times_Y Y_1$, $f_1=f_{(Y_1)}:S_1\to Y_1$, there exists a point $x_1$ of $f_1^{-1}(y_1)$
\oldpage[IV-4]{182}
such that $S_1$ is the sum of two open subsets $V_1$, $X_1$, $V_1$ being finite over $Y_1$ and a neighbourhood of $x_1$.
By virtue of the relation $\boldsymbol{k}(y_1)=\boldsymbol{k}(y)$, the fibre $f_1^{-1}(y_1)$ in $S_1$ is isomorphic to $f^{-1}(y)$, so $X_1\cap f_1^{-1}(y_1)$ is finite, discrete and has $n\setmin 1$ points.
Since $f_1|X_1$ is locally of finite type and separated, one applies the induction hypothesis to this morphism: there are an \'etale morphism $Y'\to Y_1$, a point $y'\in Y'$ above $y_1$ such that $\boldsymbol{k}(y')=\boldsymbol{k}(y_1)$ and, if $X'=X\times_Y Y'$ and if $p:X'\to S_1$ is the canonical projection, a decomposition of $p^{-1}(X_1)$ as the sum of two sub-preschemes $U'$ and $X'_2$ induced on open subsets of $X'$ such that $X'_2\cap f'^{-1}(y')=\varnothing$ and that $U'$ is finite over $Y'$.
Moreover $V'=p^{-1}(V_1)$ is finite over $Y'$, and $X'$ is the sum of $U'$, $V'$ and $X'_2$; one therefore answers the question by taking for $X'_1$ the sum of $U'$ and $V'$.
\end{proof}

The following corollary improves \sref{IV.8.11.1}.

\begin{corollary}[18.12.4]
\label{IV.18.12.4}
For a morphism $f:X\to Y$ to be finite, it is necessary and sufficient that it be separated, universally closed, locally of finite type, and that $f^{-1}(y)$ be a finite discrete space for every $y\in Y$. In particular, every proper quasi-finite morphism is finite.
\end{corollary}

\begin{proof}
The necessity is clear. Conversely, apply \sref{IV.18.12.3} at an \emph{arbitrary} point $y\in Y$. With the notation of that corollary, $f'$ is closed. Since $X_2'$ is closed in $X'$, the set $f'(X_2')$ is closed in $Y'$ and does not contain $y'$. There is therefore an open neighbourhood $U'$ of $y'$ in $Y'$ such that $U'\to Y$ is of finite type and
\[
  f'^{-1}(U')=X\times_{Y}U'
\]
is finite over $U'$.

Let $U$ be the image of $U'$ in $Y$. It is open by \sref{IV.11.3.1}, since $Y'\to Y$ is \'etale, hence flat and locally of finite presentation. Moreover
\[
  X\times_{Y}U'
  =f^{-1}(U)\times_{U}U'.
\]
The morphism $U'\to U$ is faithfully flat and quasi-compact; hence \sref{IV.2.7.1},~(xv) implies that $f^{-1}(U)\to U$ is finite. Since $y$ was arbitrary, $f$ is finite.
\end{proof}

\begin{remark}[18.12.5]
\label{IV.18.12.5}
Observe that the proof of \sref{IV.18.12.4} does not use universal closedness in its full form, but only the fact that $f_{(Y')}$ is closed for every \emph{\'etale} morphism $Y'\to Y$.
\end{remark}

\begin{corollary}[18.12.6]
\label{IV.18.12.6}
Let $f:X\to Y$ be a morphism of preschemes.
The following conditions are equivalent:
\begin{enumerate}[label=\emph{\alph*)}]
  \item $f$ is a closed immersion.
  \item $f$ is a proper monomorphism.
  \item $f$ is proper and, for all $y\in Y$, the $\boldsymbol{k}(y)$-prescheme $X_y=f^{-1}(y)$ is radicial and geometrically reduced on $\boldsymbol{k}(y)$ (i.e.\ empty or $\boldsymbol{k}(y)$-isomorphic to $\operatorname{Spec}(\boldsymbol{k}(y))$).
\end{enumerate}
\end{corollary}

\begin{proof}
This follows from \sref{IV.18.12.4} in the same way that \sref{IV.8.11.5} follows from \sref{IV.8.11.1}.
\end{proof}

\begin{proposition}[18.12.7]
\label{IV.18.12.7}
Let $f:X\to Y$ be a morphism locally of finite type, and let $y$ be a point of $Y$.
There exists an open neighbourhood $U$ of $y$ such that the restriction $f^{-1}(U)\to U$ is a closed immersion if and only if $X_y$ is a radicial and geometrically reduced $\boldsymbol{k}(y)$-scheme (that is, empty or $\boldsymbol{k}(y)$-isomorphic to $\operatorname{Spec}(\boldsymbol{k}(y))$), and there exists an open neighbourhood $V$ of $y$ such that $f^{-1}(V)\to V$ is universally closed.
\end{proposition}

The following lemma will be used in the proof.

\begin{lemma}[18.12.7.1]
\label{IV.18.12.7.1}
Let $f:X\to Y$ be a closed morphism, and let $y\in Y$ be a point such that every neighbourhood of $X_y$ contains an affine neighbourhood of $X_y$ (a condition which is always satisfied when $X_y$ is empty or consists of a single point).
Then there exists an affine open neighbourhood $U$ of $y$ such that $f^{-1}(U)\to U$ is affine.
\end{lemma}

\begin{proof}
Let $U_0$ be an affine open neighbourhood of $y$ in $Y$, and let $V$ be an affine open neighbourhood of $X_y$ contained in $f^{-1}(U_0)$. Since $f$ is closed, there is an affine open neighbourhood $U\subset U_0$ of $y$ such that $f^{-1}(U)\subset V$. The restriction $g:V\to U_0$ of $f$ is affine, and hence so is its restriction $f^{-1}(U)\to U$ \sref[II]{II.1.2.5}.
\end{proof}

\begin{proof}[Proof of {\hyperref[IV.18.12.7]{Proposition 18.12.7}}]
The conditions are plainly necessary. To prove their sufficiency, we may suppose first that $X_y\ne\varnothing$, and, by the second condition, that $f$ is universally closed. By {\hyperref[IV.18.12.7.1]{Lemma~18.12.7.1}}, after shrinking $Y$ about $y$, we may also suppose that $f$ is \emph{affine}, and hence separated.

We may now suppose that $X$ and $Y$ are affine. Since $f$ is of finite type, it is \emph{proper}. It remains, by \sref{IV.18.12.4} and \sref{IV.18.12.6}, to make $f$ \emph{quasi-finite} over a neighbourhood of $y$. There is an open neighbourhood $V$ of the one-point set $X_y$ such that $f|V$ is quasi-finite \sref{IV.13.1.4}. Since $f$ is closed, there is an open neighbourhood $U$ of $y$ such that $f^{-1}(U)\subset V$. Thus $f^{-1}(U)\to U$ is proper and quasi-finite, hence \emph{finite} by \sref{IV.18.12.4}; its fibres satisfy condition~c) of \sref{IV.18.12.6}, so it is a closed immersion. The case $X_y=\varnothing$ is immediate by the same closedness argument.
\end{proof}

\begin{proposition}[18.12.8]
\label{IV.18.12.8}
Let $f:X\to Y$ be a morphism of preschemes.
For $f$ to be integral, it is necessary and sufficient that $f$ be affine and universally closed.
\end{proposition}

\begin{proof}
The necessity follows from \sref[II]{II.6.1.10}. Conversely, suppose that $f$ is affine and universally closed. We may take $Y=\operatorname{Spec}(A)$ and $X=\operatorname{Spec}(B)$; it is enough to prove that every $b\in B$ is integral over $A$ \sref[II]{II.6.1.1}.

Let $B'=A[b]\subset B$ and set $X'=\operatorname{Spec}(B')$. Then $f$ factors as
\[
  X\xrightarrow{g}X'\xrightarrow{h}Y,
\]
where $h$ is of finite type and $g$ is \emph{dominant} because $B'\to B$ is injective \sref[I]{I.1.2.7}. Since $h$ is separated and $h\circ g$ is universally closed, $g$ is universally closed \sref[II]{II.5.4.3}, \sref[II]{II.5.4.9}; being dominant, it is therefore \emph{surjective}. It follows from the same results that $h$ is universally closed as well.

\oldpage[IV-4]{183}
Thus $h$ is \emph{proper}, since it is of finite type and separated. For every $y\in Y$, the fibre $h^{-1}(y)$ is proper and affine over $k(y)$, hence \emph{finite} \sref[III]{III.4.4.2}. Thus $h$ is \emph{quasi-finite}, and \sref{IV.18.12.4} shows that $h$ is \emph{finite}. Consequently $B'$ is a finite $A$-algebra and $b$ is integral over $A$.
\end{proof}

\begin{remark}[18.12.9]
\label{IV.18.12.9}
It may be that a morphism $f:X\to Y$ is integral whenever it is separated and universally closed and, for every $y\in Y$, the induced morphism $f^{-1}(y)\to\operatorname{Spec}(\boldsymbol{k}(y))$ is integral.
To establish this it would be enough to prove that these conditions imply that $f$ is affine, or equivalently that every fibre $X_y$ is contained in an affine open neighbourhood.
\end{remark}

\begin{corollary}[18.12.10]
\label{IV.18.12.10}
Every injective universally closed morphism $f:X\to Y$ is integral.
\end{corollary}

\begin{proof}
It suffices, by \sref{IV.18.12.8}, to prove that $f$ is affine, which follows from the lemma \sref{IV.18.12.7.1} and the hypothesis.
\end{proof}

\begin{corollary}[18.12.11]
\label{IV.18.12.11}
Let $f:X\to Y$ be a morphism of preschemes \textup{(resp.}\ a locally of finite type morphism\textup{)}.
For $f$ to be a universal homeomorphism \sref{IV.2.4.2}, it is necessary and sufficient that $f$ be integral \textup{(resp.}\ finite\textup{)}, radicial, and surjective.
\end{corollary}

\begin{proof}
One knows that the conditions are sufficient \sref{IV.2.4.5}; they are necessary by virtue of \sref{IV.18.12.10} and \sref{IV.18.12.4}.
\end{proof}

The following proposition improves \sref{IV.8.11.2}.

\begin{proposition}[18.12.12]
\label{IV.18.12.12}
Every morphism $f:X\to Y$ which is quasi-finite and separated is quasi-affine.
\end{proposition}

\begin{proof}
Put $\mathscr{A}=f_*(\sh{O}_X)$. This is a quasi-coherent $\sh{O}_Y$-algebra because $f$ is quasi-compact and separated \sref{IV.1.7.4}. Let $Z=\operatorname{Spec}(\mathscr{A})$ \sref[II]{II.1.3.1}; then $f$ has the canonical factorisation
\[
  X\xrightarrow{g}Z\xrightarrow{h}Y,
\]
where $h$ is affine and $g$ corresponds to the identity automorphism of $\mathscr{A}$ \sref[II]{II.1.2.7}. It is enough to prove that $g$ is an \emph{open immersion}, or equivalently \sref{IV.17.9.1} that $g$ is \emph{\'etale} and \emph{radicial}. It is plainly enough to establish this at every point of $f^{-1}(y)$, for each $y\in Y$.

Fix $y\in Y$ and apply \sref{IV.18.12.3} to $f$, retaining its notation. The morphism $Y'\to Y$ is flat, and $f$ is quasi-compact and separated; hence there is a canonical isomorphism
\[
  f'_*(\sh{O}_{X'})
  =\mathscr{A}\otimes_{\sh{O}_Y}\sh{O}_{Y'}
  =:\mathscr{A}'
\]
by \sref{IV.2.3.1}. Thus $Z'=Z\times_{Y}Y'$ identifies with $\operatorname{Spec}(f'_*(\sh{O}_{X'}))$, and the canonical factorisation
\[
  X'\xrightarrow{g'}Z'\xrightarrow{h'}Y'
\]
is obtained from the preceding one by base change \sref[II]{II.1.2.7}.

The decomposition $X'=X_1'\amalg X_2'$ yields a direct-product decomposition $\mathscr{A}'=\mathscr{A}_1'\times\mathscr{A}_2'$ and hence $Z'=Z_1'\amalg Z_2'$, where $Z_i'=\operatorname{Spec}(\mathscr{A}_i')$ and $g'(X_i')\subset Z_i'$. Since $X_1'$ is finite over $Y'$, it is affine over $Y'$, and therefore $g'|X_1':X_1'\to Z_1'$ is an \emph{isomorphism}. Moreover $X_2'\cap f'^{-1}(y')=\varnothing$. It follows that $g'$ is \emph{\'etale} and \emph{radicial} at every point of $f'^{-1}(y')$.

Since $Y'\to Y$ is flat and locally of finite presentation, \sref{IV.17.7.4} shows that $g$ is \'etale at the projections of the points of $f'^{-1}(y')$, that is, at all points of $f^{-1}(y)$ \sref[I]{I.3.5.2}. On the other hand, the induced map
\[
  g'_{y'}:f'^{-1}(y')\longrightarrow h'^{-1}(y')
\]
is radicial; since $\boldsymbol{k}(y')=\boldsymbol{k}(y)$, it is the scalar extension of $g_y:f^{-1}(y)\to h^{-1}(y)$. Hence $g_y$ is radicial as well. Thus $g$ is \'etale and radicial at every point of $X$, which proves the proposition.
\end{proof}

The following assertion likewise improves \sref{IV.8.12.6}.

\begin{corollary}[18.12.13]
\label{IV.18.12.13}
(``Main Theorem'' of Zariski).
Let $Y$ be a quasi-compact and quasi-separated scheme, and let $f:X\to Y$ be a quasi-finite separated morphism.
Then $f$ admits a factorisation
\[
  X\xrightarrow{g}Z\xrightarrow{u}Y
\]
where $g$ is an open immersion (necessarily quasi-compact) and $u$ is finite.
\end{corollary}

\begin{proof}
Indeed, it follows from \sref{IV.18.12.12} that $f$ is \emph{quasi-affine}, hence quasi-projective.
One can therefore apply \sref{IV.8.12.8}, where the hypothesis on $\sh{O}_Y$-Module ample can in fact be replaced by the sole hypothesis that $Y$ is quasi-compact and quasi-separated: indeed, it follows from \sref{IV.8.12.3} that the existence of the factorisation announced is a \emph{local} property on $Y$, and it suffices to prove it when $Y$ is \emph{affine}.
\end{proof}

\begin{remark}[18.12.14]
\label{IV.18.12.14}
There is also a proof of \sref{IV.18.12.13} analogous to that of \sref{IV.18.12.12}. It does not use \sref{IV.8.12.8}, but it does use \sref{IV.18.12.3}, and hence \sref{IV.18.5.11}, which itself uses the ``Main Theorem'' in its local form \sref{IV.8.12.9}.

Retain the notation of the proof of \sref{IV.18.12.12}, and let $\mathscr{B}$ be the quasi-coherent $\sh{O}_Y$-algebra which is the \emph{integral closure} of $\sh{O}_Y$ in $\mathscr{A}$ \sref[II]{II.6.3.2}. Put $T=\operatorname{Spec}(\mathscr{B})$. By \sref{IV.8.12.3}, it is enough to prove that the $Y$-morphism $g:X\to T$ corresponding to the canonical injection $\mathscr{B}\to\mathscr{A}$ is an immersion; for this it is enough, by \sref{IV.17.9.1}, to show that $g$ is \'etale and radicial.

With the notation of \sref{IV.18.12.12}, we may suppose that $Y=\operatorname{Spec}(C)$ and $Y'=\operatorname{Spec}(C')$ are affine, with $C'$ an \emph{\'etale} $C$-algebra \sref{IV.17.3.2}. Write $\mathscr{A}=\widetilde{A}$, where $A$ is a $C$-algebra, put $A'=A\otimes_CC'$, and write $\mathscr{B}=\widetilde{B}$, where $B$ is the integral closure of $C$ in $A$. Recall that $A'$ is isomorphic to the \emph{product} $A_1'\times A_2'$, with $A_1'$ a \emph{finite} $C'$-algebra. It is enough to prove that
\[
  B'=B\otimes_CC',
\]
\oldpage[IV-4]{184}
which identifies by flatness with a $C'$-subalgebra of $A'$, contains $A_1'$. Indeed, $B'$ would then decompose as $A_1'\times A_2''$, where $A_2''$ is a $C'$-subalgebra of $A_2'$. Thus $T'=T\times_YY'$ would be the disjoint union of $Z_1'=\operatorname{Spec}(A_1')$ and $T_2''=\operatorname{Spec}(A_2'')$, and $g'|X_1'$ would be an \emph{isomorphism from} $X_1'$ onto $Z_1'$. One would then conclude as in \sref{IV.18.12.12}.

To prove that $B'$ contains $A_1'$, it plainly suffices to prove the following proposition, which partially extends \sref{IV.6.14.4} when the Noetherian hypotheses are removed.
\end{remark}

\begin{proposition}[18.12.15]
\label{IV.18.12.15}
Let $C$ be a ring, let $C'$ be an \textup{\'etale} $C$-algebra, let $A$ be a $C$-algebra, and let $B$ be the integral closure of $C$ in $A$. Put $A'=A\otimes_CC'$ and $B'=B\otimes_CC'$, where $B'$ is identified with a subalgebra of $A'$. Then $B'$ is the integral closure of $C'$ in $A'$.
\end{proposition}

\begin{proof}
Consider the filtered increasing family of finite-type $C$-subalgebras of $A$ and argue as in \sref{IV.6.14.4},~II. We may first suppose that $A$ is a \emph{finite-type} $C$-algebra, say $A=E/\mathfrak{J}$, where $E=C[T_1,\ldots,T_r]$. Then $A'=E'/\mathfrak{J}E'$, where $E'=C'[T_1,\ldots,T_r]$.

Let $(\mathfrak{J}_\lambda)$ be the filtered increasing family of finitely generated ideals of $E$ contained in $\mathfrak{J}$, and let $B_\lambda$ be the integral closure of $C$ in $E/\mathfrak{J}_\lambda$. The reasoning of \sref{IV.5.13.4} shows that $B=\varinjlim_\lambda B_\lambda$. Likewise, the integral closure of $C'$ in $E'/\mathfrak{J}E'$ is the direct limit of the integral closures of $C'$ in $E'/\mathfrak{J}_\lambda E'$. If the latter closure equals $B_\lambda\otimes_CC'$ for every $\lambda$, then
\[
  B'=\varinjlim
  (B_\lambda\otimes_CC')
\]
is the integral closure of $C'$ in $A'$. We are therefore reduced to the case where $A$ is \emph{finitely presented} over $C$.

We next reduce to the case where $C$ is \emph{Noetherian}. Write $C$ as the filtered union of its finite-type $\mathbb{Z}$-subalgebras $C_\alpha$. By \sref{IV.17.7.8}, for some $\alpha$ there is an \emph{\'etale} $C_\alpha$-algebra $C_\alpha'$ such that
\[
  C'=C_\alpha'\otimes_{C_\alpha}C.
\]
For $\beta\geq\alpha$, put $C_\beta'=C_\alpha'\otimes_{C_\alpha}C_\beta$; then $C'=\varinjlim_{\beta\geq\alpha}C_\beta'$. Since $A$ is finitely presented over $C$, after increasing $\alpha$ we may write
\[
  A=A_\alpha\otimes_{C_\alpha}C,
\]
with $A_\alpha$ a finite-type $C_\alpha$-algebra. Put $A_\beta=A_\alpha\otimes_{C_\alpha}C_\beta$; then $A=\varinjlim_{\beta\geq\alpha}A_\beta$ \sref{IV.8.9.1}. If $B_\beta$ is the integral closure of $C_\beta$ in $A_\beta$, the reasoning of \sref{IV.5.13.4} gives $B=\varinjlim B_\beta$. Similarly,
\[
  A'=\varinjlim_{\beta\geq\alpha}A_\beta',
  \qquad
  A_\beta'
   =A_\beta\otimes_{C_\beta}C_\beta'
   =A_\alpha\otimes_{C_\alpha}C_\beta'.
\]
The same argument shows that it is enough to prove that $B_\alpha'=B_\alpha\otimes_{C_\alpha}C_\alpha'$ is the integral closure of $C_\alpha'$ in $A_\alpha'$. Thus we may suppose that $C$ is Noetherian.

In that case the proposition is a special case of \sref{IV.6.14.4}. One should note, however, that when $C'$ is \emph{\'etale} over the Noetherian ring $C$, the proof of \sref{IV.6.14.4} no longer requires the difficult theorem \sref{IV.6.14.1}. Indeed, the successive reductions in that proof lead, \emph{without using} \sref{IV.6.14.1}, to the case in which $C$ is \emph{integral} and $A$ is its field of fractions. Then $B$ is normal, and since $\operatorname{Spec}(B')\to\operatorname{Spec}(B)$ is \emph{\'etale}, $B'$ is normal by \sref{IV.17.5.7}. The argument then uses only the elementary lemma \sref{IV.6.14.1.1}, not the difficult part of \sref{IV.6.14.1}.
\end{proof}

\begin{proposition}[18.12.16]
\label{IV.18.12.16}
Let $g:Y\to S$ be a quasi-compact morphism, and let $f:X\to Y$ be separated and quasi-finite.
For every invertible $\sh{O}_Y$-Module $\mathscr{L}$ which is ample relative to $S$ \sref[II]{II.4.6.1}, the invertible $\sh{O}_X$-Module $f^*(\mathscr{L})$ is ample relative to $S$.
\end{proposition}

\begin{proof}
By \sref{IV.18.12.12}, the morphism $f$ is quasi-affine; hence $\sh{O}_X$ is ample for $f$ by \sref[II]{II.5.1.6}. It follows from \sref[II]{II.4.6.13},~(ii) that
\[
  \sh{O}_X\otimes_{\sh{O}_X}f^*(\mathscr{L}^{\otimes n})
  =f^*(\mathscr{L}^{\otimes n})
\]
is ample relative to $S$ for all sufficiently large $n$. Since
\[
  f^*(\mathscr{L}^{\otimes n})=(f^*\mathscr{L})^{\otimes n},
\]
this implies that $f^*(\mathscr{L})$ is ample relative to $S$ \sref[II]{II.4.6.9},~(i).
\end{proof}

\begin{corollary}[18.12.17]
\label{IV.18.12.17}
Let $Z$ be an affine scheme, $h:X\to Z$ a morphism of finite type, $\mathscr{L}$ an invertible $\sh{O}_X$-Module.
With the notation of \textup{\sref[II]{II.4.5.2}}, the conditions of \textup{\sref[II]{II.4.5.2}} \textup{(which are equivalent to the fact that $\mathscr{L}$ is \emph{ample})} are also equivalent to each of the following:
\begin{enumerate}
  \item[\textup{b$''$)}] $h$ is separated, one has $G(\varepsilon)=X$, and the canonical morphism $u:G(\varepsilon)\to\operatorname{Proj}(S)$ has finite and discrete fibres.
  \item[\textup{b$'''$)}] One has $G(\varepsilon)=X$ and the canonical morphism $u$ is radicial.
\end{enumerate}
\end{corollary}

\begin{proof}
If $Z=\operatorname{Spec}(A)$, the ring $S$ is canonically equipped with a structure of graded $A$-algebra, hence one has a structural morphism $g:\operatorname{Proj}(S)\to Z$, which is separated \sref[II]{II.2.4.2}, and one has $h=g\circ u$.

Since every radicial morphism is separated \sref{IV.1.8.7.1}, \emph{b$'''$)} implies that $h$ is separated, hence implies \emph{b$''$)}.
Since condition \emph{b)} of \sref[II]{II.4.5.2} evidently implies \emph{b$'''$)}, it remains to see that \emph{b$''$)} implies that $\mathscr{L}$ is ample.
Since $h=g\circ u$ is of finite type by hypothesis and $g$ is separated, $u$ is also of finite type \sref[I]{I.6.3.4}, so hypothesis \emph{b$''$)} implies that $u$ is quasi-finite and separated \sref[I]{I.5.5.1}.
To prove that $\mathscr{L}$ is ample, we shall apply \sref{IV.18.12.16}.
Set $Y=\operatorname{Proj}(S)$.
Since $X$ is quasi-compact, there exist an integer $n>0$ and a finite number of elements $f_j\in S_n$ such that the inverse images $u^{-1}(D_+(f_j))$ cover $X$; one may therefore regard $u$ as a quasi-finite and separated morphism from $X$ to the open subset $Y'=\bigcup_j D_+(f_j)$ of $Y$.
Consider then the invertible $\sh{O}_{Y'}$-Module $\mathscr{L}'=\sh{O}_Y(n)|Y'$ \sref[II]{II.2.5.8}; it is very ample relative to $Z$ \sref[II]{II.4.4.3}, and $u^*(\mathscr{L}')$ is by definition none other than $\mathscr{L}^{\otimes n}$
\oldpage[IV-4]{185}
\sref[II]{II.3.7.9},~(i).
By virtue of \sref{IV.18.12.16}, $\mathscr{L}^{\otimes n}$ is therefore ample relative to $Z$ (which is equivalent to saying that it is ample \sref[II]{II.4.6.6}), and consequently so is $\mathscr{L}$ \sref[II]{II.4.5.6},~(i).
\end{proof}

We leave it to the reader to state in the same way conditions equivalent to those of \sref[II]{II.4.6.3} for relatively ample $\sh{O}_X$-Modules.
